\documentclass{article}
\usepackage[utf8]{inputenc}
% Margin
\usepackage[a4paper,margin=1in]{geometry}
\usepackage{xcolor}
\usepackage{url}
\usepackage{makecell}
% For internal + external links
\usepackage{hyperref}
% Customise clickable links, e.g. colour
\hypersetup{
    colorlinks=true,
    citecolor=blue,
    filecolor=blue,
    % Black because table of contents should be black instead of blue
    % Will be updated to blue after table of contents
    linkcolor=black,
    urlcolor=blue
}
% Required for inserting images
\usepackage{graphicx}
\graphicspath{{img/}}
\usepackage{float}
\usepackage{placeins}  % Optional but recommended
\usepackage{amsthm}  % Required for self-defined environment (must be placed before amsmath)
\usepackage{amsmath}
% Correct character mapping
\usepackage[T1]{fontenc}
% Required for quote in code
\usepackage{textcomp}
\usepackage{listings}
\lstset{
    % Ensure ' stays straight and does not curl to be LaTeX '
    upquote=true,
    % Help with character spacing
    columns=fullflexible,
    % Force the ' to be a standard quote
    literate={'}{{\textquotesingle}}1 {"}{{\textquotedbl}}1
}
\lstdefinelanguage{ini}{
    morecomment=[l]{;},
    morecomment=[l]{\#},
    morestring=[b]",
    alsoletter={_},
    keywords={
        initial_connection,
        full_direct,
        partial_direct,
        unconnected
    },
    keywordstyle=\color{blue}\bfseries,
    commentstyle=\color{teal},
    stringstyle=\color{yellow},
    backgroundcolor=\color{lightgray},
    frame=single,
}
\lstset{ 
    language=Python,                      % Set language to Python
    basicstyle=\ttfamily\footnotesize,    % Set font to monospaced and small size
    keywordstyle=\color{blue}\bfseries,   % Keywords in blue and bold
    commentstyle=\color{teal},            % Comments in teal
    stringstyle=\color{yellow},           % Strings in yellow
    backgroundcolor=\color{lightgray},    % Light gray background
    numbers=left,                         % Line numbers on the left
    numberstyle=\tiny\color{gray},        % Line number style
    stepnumber=1,                         % Line number step (every line)
    numbersep=5pt,                        % Distance from line numbers to code
    showstringspaces=false,               % Don't show spaces in strings
    breaklines=true,                      % Break lines when necessary
    breakatwhitespace=true,               % Break lines at whitespace
    showtabs=false,                       % Don't show tab characters
    captionpos=b,                         % Position caption at the bottom
    tabsize=4,                            % Set tab size
    escapeinside={\%*}{*)},               % Enable escape sequences for LaTeX inside code
}
\usepackage{newtxtext,newtxmath}  % Times New Roman text + math
% Self-define environment
\newtheorem{theorem}{Theorem}
% No identation for any paragraphs
% \usepackage{indentfirst}
\setlength{\parindent}{0pt}
% Space between paragraphs
\setlength{\parskip}{1em}


% ================================================
\title{COMP30040 Third Year Project Laboratory 2025/26 Full Year}
\author{Hin Yui Jacob Yip}
\date{September 2025}


\begin{document}


\maketitle


% Generate a table of contents
\newpage
\tableofcontents
\newpage


% Change colour of links from black to blue
\hypersetup{
    linkcolor=blue
}


\section{Abstract}
Machine Learning has been a hot topic for discussion nowadays as keywords like "LLM", "GPT" etc. are seen often. Most traditional machine learning models share a common feature that the model runs independently. 
In this project, an advanced technique called \textbf{ensemble learning} is introduced which requires a committee of models. This dissertation addresses the benefits of ensemble learning compared to a single model, both mathematically and empirically through several mathematical derivations and experiments. 
This dissertation also introduces some niche algorithms, which can be combined with ensemble learning (e.g. NEAT). 
Key project achievements include mathematically and empirically proving that an ensemble model performs better than a single MLP, the importance of diversity in ensemble learning, combining NEAT with ensemble learning and an algorithm created by the author that converts a dynamic NEAT neural network to a GPU-runnable format in inference mode. 


\section{Declaration}
No portion of the work referred to in this dissertation has
been submitted in support of an application for another
degree or qualification of this or any other university or
other institute of learning. 


\section{Copyright}

\begin{enumerate}
    \item The author of this thesis (including any appendices and/or schedules to this thesis) owns certain copyright or related rights in it (the ``Copyright'') and s/he has given The University of Manchester certain rights to use such Copyright, including for administrative purposes. 
    
    \item Copies of this thesis, either in full or in extracts and whether in hard or electronic copy, may be made \textbf{only} in accordance with the Copyright, Designs and Patents Act 1988 (as amended) and regulations issued under it or, where appropriate, in accordance with licensing agreements which the University has from time to time. This page must form part of any such copies made. 
    
    \item The ownership of certain Copyright, patents, designs, trade marks and other intellectual property (the ``Intellectual Property'') and any reproductions of copyright works in the thesis, for example graphs and tables (``Reproductions''), which may be described in this thesis, may not be owned by the author and may be owned by third parties. Such Intellectual Property and Reproductions cannot and must not be made available for use without the prior written permission of the owner(s) of the relevant Intellectual Property and/or Reproductions. 
    
    \item Further information on the conditions under which disclosure, publication and commercialisation of this thesis, the Copyright and any Intellectual Property and/or Reproductions described in it may take place is available in the University IP Policy (see \url{http://documents.manchester.ac.uk/DocuInfo.aspx?DocID=24420}), in any relevant Thesis restriction declarations deposited in the University Library, The University Library's regulations (see \url{http://www.library.manchester.ac.uk/about/regulations/}) and in The University's policy on presentation of Theses. 
\end{enumerate}


\section{Acknowledgement}
I would like to express my sincere gratitude to my supervisor, Professor Gavin Brown, for his invaluable guidance and unwavering support throughout this project. I also thank my second marker for their time in reviewing and evaluating this work. 
I would also like to acknowledge the assistance given by Research IT and the use of the Computational Shared Facility at The University of Manchester. 


\section{Introduction}
There are six stages in this project. One experiment had been performed for each stage to study different behaviours of ensemble models. The six stages are designed to be \textbf{sequentially structured}, with each phase depends on work from previous stages. Below contains the summaries of each stage: 

\begin{enumerate}
    \item \textbf{Section~\ref{sec:stage_1}} \\
    This stage studies the behaviour of a MLP with different hyperparameters. It serves as a preparatory step for more critical experiments later. 

    \item \textbf{Section~\ref{sec:stage_2}} \\
    This stage studies the behaviour of a traditional ensemble model (i.e. an arithmetic-mean-aggregator) with different ensemble sizes $m$. Its codebase is dependent on the development in Section~\ref{sec:stage_1}. 

    \item \textbf{Section~\ref{sec:stage_3}} \\
    This stage studies the behaviour of a NCL ensemble model (where a penalty term $p_i$ is added to the loss function of base learners to encourage diversity) with different values of $\lambda$, which controls the weighting of $p_i$. Its codebase is dependent on the development in Section~\ref{sec:stage_2}. 

    \item \textbf{Section~\ref{sec:stage_4}} \\
    This stage studies the behaviour of an ensemble model that combines NEAT and NCL. It has made reference to the niche algorithm proposed in Sheng et al.\ \cite{sheng2017niching}. 

    \item \textbf{Section~\ref{sec:stage_5}} \\
    This stage studies the effect of choosing different aggregators on ensemble performance. It requires all ensemble models developed in Section~\ref{sec:stage_2}, Section~\ref{sec:stage_3} and Section~\ref{sec:stage_4}. 

    \item \textbf{Section~\ref{sec:stage_6}} \\
    This stage studies the impact of two different recursive architectures on the ensemble performance. It requires all ensemble models and aggregators developed in Section~\ref{sec:stage_2}, Section~\ref{sec:stage_3}, Section~\ref{sec:stage_4} and Section~\ref{sec:stage_5}. 
\end{enumerate}


\section{Notation and Terminology}
\begin{itemize}
    \item $\lambda$ — Correlation penalty coefficient used in this work
    \item $\gamma$ — Correlation penalty coefficient used in Brown et al.\ \cite{brown2005managing}, which is equivalent to $\lambda$ in this work
    \item $\lambda^*$ — Optimal correlation penalty coefficient
    \item $\rho$ — Diversity coefficient (see Appendix~\ref{appendix:diversity_coefficient} for the formula of $\rho$)
    \item $\mathbf{x}$ — Input of a regression prediction
    \item $y$ — True value of a given input $\mathbf{x}$
    \item $\hat{y}$ — Prediction of 1 base learner
    \item $\bar{y}$ — Mean prediction of all base learners
    \item $y^*$ - The Bayes Model
    \item $\hat{y^*}$ - Prediction of the Bayes Model
    \item $p_i$ - The correlation penalty term added to MSE of the $i$th base learner in NCL
    \item $m$ - Ensemble size
    \item $\mathcal{S}$ - The set of all possible training datasets
    \item $\sigma$ - Niche radius for NEAT-NCL in Section~\ref{sec:stage_4}
    \item Idempotent - An idempotent mathematical function gives the same result when applied multiple times as it does when applied once. Formally, a function $f$ is idempotent if $f(f(x)) = f(x)$ for all $x$ in its domain. 
    \item MLP — Multi-Layer Perceptron
    \item MSE — Mean Squared Error (see Appendix~\ref{appendix:mse} for the formula of MSE)
    \item NCL - Negative Correlation Learning
    \item NEAT - NeuroEvolution of Augmenting Topologies
    \item RMSE - Root Mean Squared Error (see Appendix~\ref{appendix:rmse} for the formula of RMSE)
\end{itemize}


\section{Background and Related Work}


\subsection{Viewing supervised learning systematically}
A supervised learning system comprises four essential elements: \textbf{data}, \textbf{model}, \textbf{loss} and \textbf{algorithm}. The experimental designs in this project are framed by this conceptual structure. 

\paragraph{Data}
Data consists of examples defining the desired system behavior, used for both training and evaluation. This project utilises the \textit{Boston Housing dataset}. 

\paragraph{Model}
The model is the predictive component. While this project focuses on a single \textit{MLP} and various ensemble methods, supervised learning models are not limited to neural networks and include alternatives such as decision trees and logistic regression. 

\paragraph{Loss}
The loss function mathematically assesses model performance by scoring the accuracy of predictions against the ground truth. This project employs MSE and its variants. 

\paragraph{Algorithm}
The algorithm is the optimisation procedure used to update model parameters (or \textit{weights} in neural networks) based on the loss output. This project utilised gradient descent and genetic algorithm (in Section~\ref{sec:stage_4}). 


\subsection{Bias Variance Decomposition for squared loss}
\begin{theorem}\label{theorem:bias_variance_decomposition}
    Given a learning alrothim that produces a predictor (its prediction is $\hat{y}$) from a random training set $\mathcal{S}$, the expected risk is given as: 
    \begin{align*}
        \underbrace{\mathbb{E}_{\mathcal{S}}\left[\mathbb{E}_{\mathbf{x}y}\left[(y - \hat{y})^2\right]\right]}_{\text{expected risk}} = \underbrace{\mathbb{E}_{\mathbf{x}y}\left[(y - \hat{y^*})^2\right]}_{\text{noise}} + \underbrace{\mathbb{E}_{\mathbf{x}y}\left[\left(\hat{y^*} - \mathbb{E}_{\mathcal{S}}[\hat{y}]\right)^2\right]}_{\text{bias}} + \underbrace{\mathbb{E}_{\mathbf{x}y}\left[\mathbb{E}_{\mathcal{S}}\left[(\hat{y} - \mathbb{E}_{\mathcal{S}}[\hat{y}])^2\right]\right]}_{\text{variance}}
    \end{align*}
\end{theorem}
Refer to Appendix~\ref{appendix:bias_variance_decomposition} for its proof. 


\subsection{What is ensemble learning? }
Ensemble learning involves training a ``committee'' of models (i.e. \textit{base learners}) rather than a single individual model. While common algorithms include \textit{Bagging} and \textit{Boosting}, this project explores \textit{NCL}. The fundamental principle of ensembling is to encourage base learners to capture diverse aspects of the training data. As detailed in Section~\ref{sec:stage_2}, this increased diversity enhances overall system performance compared to a single-model approach. 


\subsection{What is Negative Correlation Learning? }
\label{subsec:background_ncl}
\begin{theorem}\label{theorem:ncl_differentiation_incorrect}
    When $\bar{y}$ is assumed to have a constant value with respect to $\hat{y}_i$, the derivative of $p_i$ with respect to $\hat{y}_i$ is given below: 
    \begin{align*}
        \frac{\partial{p_i}}{\partial{\hat{y}_i}} = -(\hat{y}_i - \bar{y})
    \end{align*}
\end{theorem}
Refer to Appendix~\ref{appendix:ncl_differentiation_incorrect} for its proof. 

\begin{theorem}\label{theorem:ncl_differentiation_correct}
    When $\bar{y}$ is not assumed to have a constant value with respect to $\hat{y}_i$, the derivative of $p_i$ with respect to $\hat{y}_i$ is given below: 
    \begin{align*}
        \frac{\partial{p_i}}{\partial{\hat{y}_i}} = \frac{-2 (m - 1)}{m} (\hat{y}_i - \bar{y})
    \end{align*}
\end{theorem}
, where $m$ is the ensemble size. 
\\
Refer to Appendix~\ref{appendix:ncl_differentiation_correct} for its proof. 

\textbf{Negative Correlation Learning} (\textit{NCL}) is a key ensemble method explored in this project. Unlike traditional approaches such as bagging and boosting, which promote diversity passively through data sub-sampling, NCL enforces diversity actively by adding a penalty term to each base learner's loss function. 
Originally proposed by Liu et al.\ \cite{liu1999negative}, the loss function of the $i$th base learner is $\mathcal{L}_i = \frac{1}{2} (\hat{y}_i - y)^2 - \lambda p_i$, where $p_i$ is the correlation-penalty term that penalises similarity between predictions. The initial formula assumed $p_i = (\hat{y}_i - \bar{y}) \sum_{j \neq i} (\hat{y}_j - \bar{y})$, with $\bar{y}$ treated as a constant with respect to $\hat{y}_i$. Based on this assumption, Liu et al.\ \cite{liu1999negative} proposed Theorem~\ref{theorem:ncl_differentiation_incorrect}. 
This assumption was later shown to be incorrect. Consequently, the penalty term was revised to $p_i = -(\hat{y}_i - \bar{y})^2$ (see Appendix~\ref{appendix:correlation_penalty_function_conversion}) with an updated derivation given in Theorem~\ref{theorem:ncl_differentiation_correct}. 
This revision  affected the upper bound of $\lambda$. As shown by Brown et al.\ \cite{brown2005managing}, $\lambda_{upper\_bound} = \frac{m}{m - 1}$ under an independence assumption between $\bar{y}$ and $\hat{y}_i$. Otherwise, $\lambda_{upper\_bound} = \frac{m^2}{2(m - 1)^2}$. This project adopts the latter. 


\subsection{What is NEAT? }
\label{subsec:background_neat}
\textbf{NeuroEvolution of Augmenting Topologies} (\textit{NEAT}) is a genetic algorithm-based method for evolving neural networks. Unlike standard approaches that utilise a fixed architecture, NEAT evolves both weights and topologies by dynamically adding or removing nodes and connections while employing speciation to protect structural innovations. Consequently, the hierarchy of a NEAT base learner is represented as a directed graph of nodes and edges. 
While traditional NEAT is fundamentally an evolutionary algorithm and does not utilise gradient-based optimisation, this project integrates NEAT with NCL and incorporates gradient descent to refine network parameters. 

\subsubsection{NEAT Notation and Terminology}
The following definitions establish the core NEAT terminology used in this report: 

\paragraph{Genome}
A \textbf{genome} is the encoded representation of a neural network, analogous to its ``DNA''. While various methods exist, such as bit-string representation, this project utilises \textbf{graph encoding} to more naturally reflect the network's topology. Within this framework, a genome consists of a directed graph where nodes represent neurons and edges represent weighted connections. 

\paragraph{Phenome}
A \textbf{phenome} is the functional neural network expressed by its corresponding genome. Utilising a direct encoding scheme, NEAT explicitly represents every structural element such as nodes and edges as a specific gene. Think of the genome as the ``DNA'' while the phenome as the resulting ``organism''. 

\paragraph{Crossover}
It is the process of combining 2 parent genomes to produce an offspring genome. There are 3 types of genes. 

\begin{enumerate}
    \item \textbf{Matching Gene} \\
    A gene that has the same innovation number and exist in both parent genomes. They will be randomly selected to pass down to the offspring genome. The minimum and maximum innovation numbers of all matching genes form the \textbf{shared historical range}. 

    \item \textbf{Disjoint Gene} \\
    A non-matching gene within the shared historical range that appears in 1 parent genome only. Disjoint genes of the less-fit parent genome are ignored during crossover. 

    \item \textbf{Excess Gene} \\
    A non-matching gene beyond the shared historical range that appears in 1 parent genome only. Excess genes of the less-fit parent genome are ignored during crossover. 
\end{enumerate}

Figure~\ref{fig:background_neat_crossover} illustrates a NEAT crossover between two parent genomes. In this example, the phenome processes three input features (nodes~1-3) to produce a single output (node~4), with bold arrows indicating the direction of data flow. 
Each gene is represented by a rectangle detailing its innovation number, connection node IDs and activation status. Blue and green backgrounds identify genes from Parent~1 and Parent~2 respectively. 
During crossover, genes are aligned by their innovation numbers: 

\begin{enumerate}
    \item \textbf{Matching Gene} \\
    Randomly selected from either parent. 

    \item \textbf{Disjoint and Excess Gene} \\
    Inherited solely from the more fit parent (Parent~2) while those from the less fit parent (Parent~1) are skipped. 
\end{enumerate}

As shown, within the shared historical range $[1, 7]$, genes~4, 8 and 9 from Parent~1 are ignored. Consequently, the resulting offspring primarily inherits its genetic structure from the more fit parent. 

% 1 Image
\begin{figure}[H]
    \centering
    \includegraphics[width=0.64\textwidth]{img/background/neat-crossover.png}
    \caption{Example of NEAT crossover between 2 parent genomes}
    \label{fig:background_neat_crossover}
\end{figure}

\paragraph{Speciation}
Speciation partitions the population into distinct groups based on topological similarity, enabling networks to compete primarily within their own niches. This mechanism protects structural innovation by ensuring that nascent, lower-fitness genomes are not prematurely outcompeted by established, high-fitness individuals. Protection is further maintained by setting an appropriate \textit{elitism} threshold within each \textbf{species} (a subgroup of genetically similar genomes). 

\paragraph{Innovation number}
An innovation number is a globally unique identifier assigned to each new node or connection created through mutation. This historical marking tracks the origin of genetic features, allowing the system to identify homologous genes across structurally diverse networks. By enabling this alignment, innovation numbers facilitate efficient, meaningful recombination without the need for complex topological mapping. 

\paragraph{Genetic distance}
Genetic distance (also known as \textit{compatibility distance}) is the metric used to quantify the similarity between two genomes for speciation. As defined in Equation~\ref{equation:neat_ncl_genetic_distance}, this calculation incorporates (potentially with different weightings) the number of excess and disjoint genes, alongside the average weight differences between matching genes. 

\paragraph{Elitism}
It is a mechanism to ensure the best individuals within a species (more common)/population (less common) can survive to the next generation. Normally, the next generation is created via mutation and crossover. However, due to randomness, some best genomes might be lost. Elitism protects the ``elite'' genomes by carrying them over unchanged while the remaining genomes have to compete for reproduction. Elitism can be thought as a special treatment to the ``elite'' genomes and the number of ``elite'' genomes allowed to enter the next generation without competition can be set by developers. 

\paragraph{Stagnation}
It is the situation where a species stops improving in their fitness levels over a particular number of generations. A species is known as \textit{stagnant} in that case. Stagnant species are eventually removed from the population to focus resources on more promising species. 

\paragraph{Mutation}
It is a process of introducing random changes to genomes. This can be in the form of adjusting connection weight, adding/removing a connection, adding/removing a neuron. Structural mutations allow networks to become more complex over time. 

\paragraph{Fitness}
It is a score used to measure how well a phenome performs on a given task. A higher fitness means the phenome is fitter. The score is then used to select which genomes to survive in the next evolution epoch and which genomes to reproduce to give stronger offsprings. 

\paragraph{Fitness sharing}
It is a technique designed to maintain population diversity and prevent premature convergence in Genetic Algorithms. It allows different network topologies to evolve and co-exist without being immediately outcompeted by a single, temporarily superior species. It achieves this by reducing the fitness of individuals that are similar to each other, effectively forcing the population to spread out and explore different regions of the search space. 

\paragraph{Sharing factor}
It is a number that quantifies how crowded a genome's niche is. The higher the sharing factor, the more genomes are similar to the current genome. It is computed based on similarity to other genomes via a distance metric, which is can be based on genotypic similarity (e.g. Hamming distance for binary representations) or phenotypic similarity (e.g. Euclidean distance in the solution space). Such calculation is used in fitness sharing to compute the adjusted, final fitness of a phenome by the formula $adjusted\_fitness = \frac{raw\_fitness}{sharing\_factor}$ before selection for the next evolution epoch. 

\paragraph{Niche radius}
A threshold value used to determine whether two phenomes are considered similar enough to share fitness. It is denoted as $\sigma$ and is also known as \textit{sharing radius}. It is usually a hyperparameter set by developers but in Section~\ref{sec:stage_3}, it is computed dynamically. 


\subsubsection{Steps of traditional NEAT}
Below contains the steps for the NEAT evolutionary algorithm: 

\begin{enumerate}
    \item \textbf{Initialisation} \\
    The process starts by generating an initial population of candidate solutions. The initial solutions are generated randomly and each solution is typically represented as a “genome”, which encodes its features and characteristics.

    \item \textbf{Evaluation} \\
    Each individual in the population is evaluated according to how well it solves the given problem. This is done using a fitness function which assigns a numerical score to each solution. The higher the fitness, the better the solution. When fitness sharing is enabled, the raw fitness is updated to take into account sharing factor and the adjusted fitness is used for the next step. 

    \item \textbf{Selection} \\
    Based on the calculated fitness, a subset of the population is selected to become the “parents” of the next generation. Individuals with higher fitness scores have a higher probability of being selected. This is to encourage the reproduction of stronger offsprings in the next step. 

    \item \textbf{Reproduction} \\
    The selected individuals then reproduce to create “offspring” for the next generation using two genetic operators, crossover and mutation. 

    \item \textbf{Replacement} \\
    The newly generated offspring replaces the current population (either entirely or partially), forming the new generation. If elitism is enabled, $n$ (specified by developers) ``elite'' genomes are copied directly to the next generation without modification. 

    \item \textbf{Termination} \\
    Steps 2-5 are repeated until one of the two conditions is met. When either a certain fitness threshold (specified by developers) is reached (i.e. a satisfactory solution is found) or the maximum number of generations is reached, the process of NEAT evolution terminates. 
\end{enumerate}


\subsection{Graph algorithms}


\subsubsection{Topological sort}
Topological sort is a linear ordering of nodes in a \textbf{Directed Acyclic Graph} (\textit{DAG}) where for every directed edge \texttt{(node\_start, node\_end)}, \texttt{node\_start} comes before \texttt{node\_end}. It is used to schedule tasks with dependencies. 
Below contains the pseudo-code of the algorithm: 

\begin{lstlisting}[language=Python, caption={Depth-first search used in topological sort}]
function depth_first_search(node: Node, nodes_in_topological_sort: [Node]):
    # Get the state of this node: 
    # 'N' - New node that has not been explored
    # 'O' - Open node whose successors are currently being explored
    # 'F' - Finished node whose all successors have been explored
    node_state = get_state(node)

    if node_state == 'N':
        # Mark this node as 'O'
        set_state(node, 'O')
    else if node_state == 'O':
        # Cycle is detected
        # Error is raised because a DAG is expected in order to find a topological sort
        raise Error("Detected cycle in the graph")
    else if node_state == 'F':
        # This node has been explored so we do nothing
        pass
    else:
        # Invalid node_state
        raise Error("Invalid node_state (expect 'N'/'O'/'F')")
    
        for successor in node.successors:
            # successor is a child node of the current node
            depth_first_search(successor, nodes_in_topological_sort)
        
        # All successors of the current node have been explored
        # Mark the current node as 'F'
        set_state(node, 'F')

        # Prepend the current node to the beginning of the list of nodes in topological sort
        nodes_in_topological_sort.prepend(node)
\end{lstlisting}

\begin{lstlisting}[language=Python, caption={Topological sort}]
function topological_sort(nodes: [Node]):
    # A list of nodes in topological sort
    nodes_in_topological_sort = []

    # Set all nodes to their initial states 'N'
    for n in nodes:
        set_state(n, \'N\')
    
    # Get all starting nodes in the DAG
    start_nodes = find_all_start_nodes(nodes)

    for start_node in start_nodes:
        # Error checking to make sure a start node retrieved from findl_all_start_nodes(nodes) has the correct state
        assert get_state(start_node) == 'N'

        depth_first_search(start_node, nodes_in_topological_sort)
    
    return nodes_in_topological_sort
\end{lstlisting}


\section{Project Settings}


\subsection{Datasets}
\label{subsec:project_settings_datasets}
The Boston Housing dataset was used, where the task was to predict the median house price given a number of demographic features. There were 506 examples, each containing 13 input variables (12 continuous, 1 binary) and 1 continuous output variable in the range $0$ to $50$. All input variables were linearly rescaled, independently of one another, to be in the range $[0, 1]$ and the output variable was linearly rescaled to $[-1, 1]$. The data were converted to batches of PyTorch tensors with a batch size of $32$ (except for NEAT ensemble models where a batch size of $506$ was used). 


\subsection{Performance Metric}
In this project, MSE served as the primary metric for evaluating model performance. While certain experiments, such as the NCL study in Section~\ref{sec:stage_3}, incorporated a penalty term into the training loss, MSE remained the consistent evaluation standard for several reasons: 

\begin{enumerate}
    \item \textbf{Task Alignment} \\
    The Boston Housing dataset presents a regression challenge, for which MSE is the standard performance metric. 

    \item \textbf{Error Weighting} \\
    By squaring the prediction errors, MSE heavily penalises large outliers, forcing the model to minimise significant predictive errors. 

    \item \textbf{Mathematical Compatibility} \\
    The differentiability of MSE is particularly advantageous for NCL. With an appropriately tuned $\lambda$, the MSE derivative can simplify the NCL loss function by canceling specific components of the penalty term, resulting in a simpler derivative representation (refer to Theorem~\ref{theorem:loss_function_ncl_lambda_1}). 
\end{enumerate}


\subsection{Generating Experiment Data}
To ensure flexibility, all experiments were conducted in two stages. The first stage generated experimental data and outputted a \texttt{csv} file while the second stage produced the corresponding graphs from this output. 

This design enabled graph generation without re-running experiments. 


\subsection{Model Architecture Representation}
MLP architecture was varied to introduce different model complexities. MLPs differed in architectures in terms of numbers of hidden nodes and hidden layers. A MLP could have $2$, $6$ or $12$ hidden nodes in each hidden layer. A MLP could have $1$, $3$ or $6$ hidden layers. If a MLP had more than 1 hidden layer, the number of hidden nodes in each hidden layer was fixed to $2$. The idea was to make sure the MLPs of different model architectures to be compared have the same possible total numbers of hidden nodes, which were $2$, $6$ or $12$ in the experiment. 

The model architecture was represented by a Python list. The table below contained all model architectures used in this project: 

\begin{table}[H]
    \centering
    \begin{tabular}{|c|c|}
        \hline
        Model architecture & Description \\
        \hline
        \small\texttt{[2]} & 1 hidden layer of 2 hidden nodes \\
        \hline
        \small\texttt{[2, 2, 2]} & 3 hidden layers of 2 hidden nodes \\
        \hline
        \small\texttt{[2, 2, 2, 2, 2, 2]} & 6 hidden layers of 2 hidden nodes \\
        \hline
        \small\texttt{[6]} & 1 hidden layer of 6 hidden nodes \\
        \hline
        \small\texttt{[12]} & 1 hidden layer of 12 hidden nodes \\
        \hline
    \end{tabular}
    \caption{Overview of the neural network model architectures used in the experiments}
    \label{table:background_model_architecture}
\end{table}


\subsection{Experiment Graphs}
Experiment graphs were generated using the \texttt{plot} function from the Python visualisation library \texttt{Matplotlib}. 


\subsection{Potential Improvements}
Experiments were not repeated to reduce computational cost. Although repetition could mitigate noise, this project prioritised exploring a wider range of ensemble behaviours. The resulting noise was considered acceptable as key patterns remained observable. 

Hyperparameter tuning was applied only to critical parameters (e.g. $\lambda$) while others used default values (e.g. $learning\_rate = 0.001$). Although further tuning could be performed, the limited number and lower impact of these parameters made this approach sufficient for the study. 


\section{Stage 1: Building a MLP with PyTorch}
\label{sec:stage_1}


\subsection{Motivations}
All subsequent experiments used the same dataset, with MLPs as base learners. This stage validated the correctness of the PyTorch MLP implementation and identified suitable hyperparameters (e.g. number of training epochs). 
It served as a preparatory step for more critical experiments later. 


\subsection{Design and Implementation}


\subsubsection{Model construction}
\textit{Object-Oriented Programming} was adopted to improve maintainability and reusability. A custom MLP class (\texttt{MLP} in \path{mlp.py}) was implemented by extending \texttt{torch.nn.Module}, the base class for PyTorch neural networks. 

Each layer was fully connected, implemented using paired lists of \texttt{torch.nn.Linear} layers and activation functions (e.g. \texttt{torch.nn.ReLU}), with equal length. Predictions were generated by sequentially applying \texttt{activation\_function(linear\_layer(...))} across all layers. 

The class also provided helper methods to construct MLPs from alternative specifications, such as layer sizes or NEAT-Python genome instances. 


\subsubsection{Experiment}
\paragraph{Varied hyperparameters}
Table~\ref{table:experiment_1_varied_hyperparameters} contains all varied hyperparameters in this experiment: 

\begin{table}[H]
    \centering
    \begin{tabular}{|c|c|p{5cm}|c|}
        \hline
        Hyperparameter & Type & Choices / Ranges & Description \\
        \hline
        \small\texttt{epoch\_nums} & Integer & $[10, 1000]\text{, step}=10$ & Number of training epochs \\
        \hline
        \small\texttt{hidden\_sizes} & Python list & \{\texttt{[2]}, \texttt{[2, 2, 2]}, \texttt{[2, 2, 2, 2, 2, 2]}, \texttt{[6]}, \texttt{[12]}\} & Model architecture \\
        \hline
    \end{tabular}
    \caption{Varied hyperparameters and their configurations}
    \label{table:experiment_1_varied_hyperparameters}
\end{table}

\paragraph{Fixed hyperparameters}
Table~\ref{table:experiment_1_fixed_hyperparameters} contains all fixed hyperparameters in this experiment: 

\begin{table}[H]
    \centering
    \begin{tabular}{|c|c|c|c|}
        \hline
        Hyperparameter & Type & Choices / Ranges & Description \\
        \hline
        \small\texttt{input\_size} & Integer & $13$ & Number of input features of the MLP \\
        \hline
        \small\texttt{output\_size} & Integer & $1$ & Number of labels of the MLP \\
        \hline
        \small\texttt{optimizer} & Categorical & \texttt{torch.optim.Adam} & Optimisation algorithm used for training \\
        \hline
        \small\texttt{learning\_rate} & Float & $0.0001$ & Learning rate of the optimisation algorithm \\
        \hline
        \small\texttt{dropout\_rate} & Float & $0$ & Dropout rate applied to the data tensor after the input layer \\
        \hline
        \small\texttt{activation} & Categorical & ReLU & Activation function used in each dense layer of the MLP \\
        \hline
        \small\texttt{batch\_size} & Integer & $32$ & Number of training samples passed to the model in one go \\
        \hline
    \end{tabular}
    \caption{Fixed hyperparameters and their configurations}
    \label{table:experiment_1_fixed_hyperparameters}
\end{table}

\paragraph{Metric}
Table~\ref{table:experiment_1_metric} contains all evaluation metrics collected in this experiment, which were stored in \path{experiment_runner_mlp.csv}: 

\begin{table}[H]
    \centering
    \begin{tabular}{|c|c|c|}
        \hline
        Metric & Type & Description \\
        \hline
        \small\texttt{loss\_test} & Float & MSE of MLP during testing \\
        \hline
        \small\texttt{loss\_train} & Float & MSE of MLP during training \\
        \hline
    \end{tabular}
    \caption{Evaluation metrics}
    \label{table:experiment_1_metric}
\end{table}

\paragraph{Further notes}
\texttt{torch.optim.Adam} was chosen for its rapid convergence. Alternative optimisers and hyperparameter settings (e.g. learning rate) could be explored with additional time. The use of ReLU is discussed in Section~\ref{subsubsec:dynamic_nn_on_gpu}. 

To exploit parallelism, all models were trained on CSF3 \textit{A100 GPUs}. For efficiency, each epoch computed training loss, then temporarily switched the model to evaluation mode to measure testing loss before resuming training. 

Trained models were saved in \path{data/experiment/mlp/model/}. 


\subsection{Expected Observations}


\subsubsection{Training loss is lower than testing loss}
An MLP iteratively optimises its parameters to minimise error on the training set. However, due to inherent stochastic variations between the training and testing distributions, performance typically degrades on unseen data. Consequently, the training loss is expected to be lower than the testing loss. 


\subsubsection{Rapid initial loss reduction with diminishing returns}
Initially, an unoptimised MLP typically exhibits high loss due to \textit{underfitting} as its initial parameters fail to capture the data's underlying structure. As optimisation progresses, the model's predictive accuracy increases, resulting in a sharp initial decline in loss. However, after a certain threshold of epochs, the rate of improvement diminishes. This plateau occurs because the model has captured the primary data patterns. Thus, subsequent training iterations yield negligible gains as the parameters converge toward a local or global optimum. 


\subsubsection{Increase in testing loss for over-training}
Overfitting occurs when a model begins to memorise noise and specific patterns in the training data rather than learning its underlying structure. Consequently, as the number of training epochs increases beyond an optimal point, it is expected that testing loss rises, signaling a decline in generalisation performance. 


\subsubsection{Model of higher complexity has a lower loss}
Increased model complexity enables the capture of intricate, non-linear relationships within the underlying data. Consequently, assuming all other hyperparameters (e.g. training epochs) remain constant, complex models are expected to achieve higher predictive accuracy and lower loss than simpler models. 


\subsection{Evaluation and Empirical Comparisons}


\subsubsection{How does testing loss behave differently from training loss and in what way do they behave the same? }
\paragraph{Experiments and Observations}
Across all experiments, the testing loss (blue line) exhibited significantly higher fluctuations than the training loss (orange line), especially in Figure~\ref{fig:mlp_mse_against_epoch_num_10_1000_fixed_learning_rate_00001_hidden_size_2_2_2} and Figure~\ref{fig:mlp_mse_against_epoch_num_10_1000_fixed_learning_rate_00001_hidden_size_2_2_2_2_2_2}. These fluctuations likely stem from the model encountering unseen data structures during random shuffling, leading to irregular predictive errors. As the number of epochs increased, these fluctuations diminished, suggesting that the model eventually captured the underlying data distribution more consistently. 

While the training loss decreased continuously, the testing loss began to increase as the number of epochs approached 1000. Such increase is a classic indicator of \textit{overfitting}, where the model prioritises memorising training noise over generalising universal patterns. This phenomenon was amplified in the deeper architecture of Figure~\ref{fig:mlp_mse_against_epoch_num_10_1000_fixed_learning_rate_00001_hidden_size_2_2_2_2_2_2} as the increased capacity of the model facilitated more aggressive memorisation of the training set compared to the simpler model in Figure~\ref{fig:mlp_mse_against_epoch_num_10_1000_fixed_learning_rate_00001_hidden_size_2_2_2}. 

\paragraph{Conclusions from Experiment}
Testing loss always fluctuates more than the training loss. However, when number of training epochs increases, the fluctuations of training loss and testing loss both happen less often. Lastly, testing loss increases when number of training epochs becomes too big but training loss does not. 


\subsubsection{How does number of training epochs affect loss? }
\paragraph{Experiments and Observations}
The experimental results revealed three distinct phases of model training. 

Initially, both training and testing losses decreased drastically (as seen before epoch 50 in Figure~\ref{fig:mlp_mse_against_epoch_num_10_1000_fixed_learning_rate_00001_hidden_size_6}, Figure~\ref{fig:mlp_mse_against_epoch_num_10_1000_fixed_learning_rate_00001_hidden_size_12} and Figure~\ref{fig:mlp_mse_against_epoch_num_10_1000_fixed_learning_rate_00001_hidden_size_2}). In addition to capturing the data's underlying structure, this rapid improvement occurred because the high initial MSE generated large gradients for backpropagation, leading to significant weight adjustments as the model captured primary data patterns. 

As training progressed (epochs 100-400), the rate of improvement diminishes. The loss curves transitioned from near-vertical drops to a gentler gradient, eventually becoming almost horizontal. This indicated \textit{performance saturation}, where the model had learnt the majority of the underlying data structure and further weight updates yielded diminishing returns. 

In the final phase (epochs 900-1000), particularly in Figure~\ref{fig:mlp_mse_against_epoch_num_10_1000_fixed_learning_rate_00001_hidden_size_2_2_2} and Figure~\ref{fig:mlp_mse_against_epoch_num_10_1000_fixed_learning_rate_00001_hidden_size_2_2_2_2_2_2}, the training loss stabilised while the testing loss began to rise. This shift signified that the model has reached its generalisation limit and had begun to overfit, where additional training cycles degraded test performance without improving training accuracy. 

\paragraph{Conclusions from Experiment}
As number of training epochs is increased, both training and testing losses decrease drastically initially. After that, the losses continued to decrease but at a slower rate as number of training epochs was increased continuously. When number of training epochs becomes too large, any further increase leads to no decrease in training loss and increase in testing loss. 


\subsubsection{How does model complexity affect loss? }
\paragraph{Experiments and Observations}
Increasing model complexity through higher hidden node counts consistently reduced both training and testing losses since a greater model capacity allowed for more precise function approximations. For instance, the 2-node model in Figure~\ref{fig:mlp_mse_against_epoch_num_10_1000_fixed_learning_rate_00001_hidden_size_2} began with losses near $0.5$ at 10 epochs, converging to approximately $0.1$ by epoch 500. In contrast, the 6-node model in Figure~\ref{fig:mlp_mse_against_epoch_num_10_1000_fixed_learning_rate_00001_hidden_size_6} demonstrated superior initial performance with a loss of $0.27$ while the 12-node model in Figure~\ref{fig:mlp_mse_against_epoch_num_10_1000_fixed_learning_rate_00001_hidden_size_12} started even lower at $0.23$. Across these configurations, while higher complexity accelerated convergence and lowered the error, all models eventually approached a similar performance plateau at $0.5$, suggesting that the dataset's inherent noise limited further optimisation regardless of network size. 

\paragraph{Conclusions from Experiment}
An increase in model complexity leads to a decrease in loss at early stage if other hyperparameters remain unchanged. However, model complexity does not affect the loss after model is converged. 


\subsubsection{Which model architecture is better? }
\paragraph{Experiments and Observations}
A comparison between models with an identical total node count revealed that \textbf{network depth} significantly worsen stability and performance more than \textbf{network width}. Specifically, architectures with a single, wider hidden layer (e.g. Figure~\ref{fig:mlp_mse_against_epoch_num_10_1000_fixed_learning_rate_00001_hidden_size_6} and Figure~\ref{fig:mlp_mse_against_epoch_num_10_1000_fixed_learning_rate_00001_hidden_size_12}) achieved consistently lower testing errors of approximately $0.05$. In contrast, deeper architectures with the same total node count (Figure~\ref{fig:mlp_mse_against_epoch_num_10_1000_fixed_learning_rate_00001_hidden_size_2_2_2} and Figure~\ref{fig:mlp_mse_against_epoch_num_10_1000_fixed_learning_rate_00001_hidden_size_2_2_2_2_2_2}) exhibited higher testing errors (around $0.2$) and greater fluctuations. These volatilities suggested that deeper, narrower layers were more susceptible to ``severely wrong predictions'' when encountering stochastic noise or unseen data structures in shuffled batches. Consequently, for this specific dataset and task, increasing layer width rather than depth yielded a lower-bias model with more stable convergence (lower variance). 

\paragraph{Conclusions from Experiment}
Models that have 1 hidden layer of more hidden nodes are more suitable to solve the current regression task compared to models that have many hidden layers of fewer hidden nodes. Furthermore, one should aim at increasing number of hidden nodes in the same hidden layer instead of number of hidden layers to increase the model complexity. 


% ================================ Images ================================


% 2 Image
\begin{figure}[H]
    \centering

    \begin{minipage}[t]{0.48\textwidth}
        \vspace{0pt}
        \centering
        \includegraphics[width=\linewidth]{img/mlp/mlp_mse_against_epoch_num_10_1000_fixed_learning_rate_00001_hidden_size_6.png}
        \caption{Impact of number of training epochs on training and testing MSEs when the model has a learning rate of 0.0001 and 1 hidden layer of 6 hidden nodes}
        \label{fig:mlp_mse_against_epoch_num_10_1000_fixed_learning_rate_00001_hidden_size_6}
    \end{minipage}
    \hfill
    \begin{minipage}[t]{0.48\textwidth}
        \vspace{0pt}
        \centering
        \includegraphics[width=\linewidth]{img/mlp/mlp_mse_against_epoch_num_10_1000_fixed_learning_rate_00001_hidden_size_12.png}
        \caption{Impact of number of training epochs on training and testing MSEs when the model has a learning rate of 0.0001 and 1 hidden layer of 12 hidden nodes}
        \label{fig:mlp_mse_against_epoch_num_10_1000_fixed_learning_rate_00001_hidden_size_12}
    \end{minipage}
\end{figure}


% 2 Image
\begin{figure}[H]
    \centering

    \begin{minipage}[t]{0.48\textwidth}
        \vspace{0pt}
        \centering
        \includegraphics[width=\linewidth]{img/mlp/mlp_mse_against_epoch_num_10_1000_fixed_learning_rate_00001_hidden_size_2_2_2.png}
        \caption{Impact of number of training epochs on training and testing MSEs when the model has a learning rate of 0.0001 and 3 hidden layers, each of 2 hidden nodes}
        \label{fig:mlp_mse_against_epoch_num_10_1000_fixed_learning_rate_00001_hidden_size_2_2_2}
    \end{minipage}
    \hfill
    \begin{minipage}[t]{0.48\textwidth}
        \vspace{0pt}
        \centering
        \includegraphics[width=\linewidth]{img/mlp/mlp_mse_against_epoch_num_10_1000_fixed_learning_rate_00001_hidden_size_2_2_2_2_2_2.png}
        \caption{Impact of number of training epochs on training and testing MSEs when the model has a learning rate of 0.0001 and 6 hidden layers, each of 2 hidden nodes}
        \label{fig:mlp_mse_against_epoch_num_10_1000_fixed_learning_rate_00001_hidden_size_2_2_2_2_2_2}
    \end{minipage}
\end{figure}


% 1 Image
\begin{figure}[H]
    \centering
    \includegraphics[width=0.48\textwidth]{img/mlp/mlp_mse_against_epoch_num_10_1000_fixed_learning_rate_00001_hidden_size_2.png}
    \caption{Impact of number of training epochs on training and testing MSEs when the model has a learning rate of 0.0001 and 1 hidden layer of 2 hidden nodes}
    \label{fig:mlp_mse_against_epoch_num_10_1000_fixed_learning_rate_00001_hidden_size_2}
\end{figure}


\section{Stage 2: Building an Ensemble Model}
\label{sec:stage_2}


\subsection{Motivations}
In Section~\ref{sec:stage_1}, a single MLP was used to solve the problem. However, it was unreliable because a MLP has its limitations and thus, might not be able to completely learn a dataset despite the number of training epochs. Most real-life important decision making is made by a committee of people instead of one individual, e.g. Chairman. Can we apply this idea to machine learning? If there are multiple models contributing to the same prediction, they can average out their individual mistakes. Moreover, with appropriate aggregation method (more studies on aggregation method were explored in Section~\ref{sec:stage_5}), the final prediction is also less prone to noise and more reliable. Thus, in this section, the effect of having predicting with multiple MLPs was explored. 


\subsection{Design and Implementation}


\subsubsection{Ensemble population risk}
The population risk of the ensemble is less than or equal to the average population risk, given by the formula below: 

\begin{theorem}\label{theorem:ensemble_population_risk_le_average_population_risk}
    When MSE is used, 
    \begin{align*}
        R(\bar{y}) \le \frac{1}{m} \sum_{i=1}^{m} R(\hat{y}_i)
    \end{align*}
\end{theorem}
Refer to Appendix~\ref{appendix:ensemble_population_risk_le_average_population_risk} for its proof. 
\\
Note that population risk is not something measurable as it requires infinite amount of data. However, Theorem~\ref{theorem:ensemble_population_risk_le_average_population_risk} still provides a good motivation for ensemble learning. 


\subsubsection{Model construction}
\paragraph{Base learner}
Each base learner within the ensemble (\texttt{TraditionalBaseLearner} in \path{traditional_base_learner.py}) was derived from the core MLP class established in Section~\ref{sec:stage_1}, utilising the same loss function while incorporating new attributes for distributed execution, e.g. \texttt{rank} and \texttt{device\_rank}. To leverage the dual-GPU architecture of CSF3, all base learneres were partitioned across two devices. Each learner was uniquely identified by a \texttt{rank:device\_index} tuple, where \texttt{rank} specified the global GPU ID and \texttt{device\_index} denoted the learner's local position on that hardware. To maximise computational throughput, \texttt{torch.cuda.stream} was implemented, enabling concurrent execution of multiple base learners on a single GPU rather than sequential processing. 

\paragraph{Ensemble model}
The ensemble model (\texttt{TraditionalEnsembleNet} in \path{traditional_ensemble_net.py}) functioned as a central coordinator for multiple \texttt{TraditionalBaseLearner} instances. To manage these learners across a distributed environment, the system leveraged PyTorch modules \textit{DistributedDataParallel} (DDP) and \textit{Multiprocessing} for process management and inter-device communication. 
The execution workflow utilised \texttt{torch.multiprocessing.spawn} to initialise individual learner processes while 
\\
\texttt{torch.distributed.init\_process\_group} enabled communication across GPUs. This infrastructure enabled \texttt{torch.distributed.all\_gather} to synchronise and collect predictions from every base learner despite their host GPU. These predictions were managed via a \texttt{torch.multiprocessing.Manager().Queue()}, providing a thread-safe shared memory space for data aggregation. 
Finally, the ensemble generated an ensemble output using an aggregator class \texttt{ArithmeticMeanVoter} in this stage, which computed the arithmetic mean of all base learner predictions. 

Steps of ensemble model prediction: 

\begin{enumerate}
    \item \textbf{Getting input tensor} \\
    An input tensor is passed to the ensemble model. 

    \item \textbf{Creating base learners} \\
    If processes of base learners are not created, they are created in all available GPUs using 
    \\
    \texttt{torch.multiprocessing.spawn} and \texttt{torch.distributed.init\_process\_group} etc. Otherwise, the existing base learners are used. 

    \item \textbf{Simultaneous individual predictions} \\
    The input tensor is broadcasted to every base learner and each base learner calcualtes their predictions simultaneously. 

    \item \textbf{Aggregating individual predictions} \\
    All base learners are blocked before their next prediction until all base learners have completed their predictions. This can be achieved using \texttt{torch.distributed.all\_gather} and \texttt{torch.cuda.synchronize()}. All predictions are then passed to CPU for aggregation. If the ensemble model is in training mode, all base learners are then updated using backpropagation. 

    \item \textbf{Output predictions} \\
    The aggregated prediction is returned. 

    \item \textbf{Termination} \\
    Steps 1-5 are repeated until all input data have been read. 
\end{enumerate}


\subsubsection{Experiment}
Experiments had been conducted to compare the performance of 2 models. In the traditional ensemble model, all base learners were trained independently. It was essentially an arithmetic-mean-aggregator of all base learners, i.e. $\hat{y}_{ensemble} = \frac{1}{m} \sum_{i=1}^{m} \hat{y}_i$. The control MLP had the same configuration, e.g. activation function, as each base learner in the ensemble model. This gave the control MLP the same model complexity as the base learners for fair comparisons. It was to serve as a reference point of the ensemble performance. 

\paragraph{Varied hyperparameters}
Table~\ref{table:experiment_2_varied_hyperparameters} contains all varied hyperparameters in this experiment: 

\begin{table}[H]
    \centering
    \begin{tabular}{|c|c|c|c|}
        \hline
        Hyperparameter & Type & Choices / Ranges & Description \\
        \hline
        \small\texttt{ensemble\_sizes} & Integer & $[2, 16]\text{, step}=1$ & $m$ \\
        \hline
        \small\texttt{epoch\_nums} & Integer & \{$200$, $400$\} & Number of training epochs \\
        \hline
        \small\texttt{hidden\_sizes} & Python list & \{\texttt{[2]}, \texttt{[2, 2, 2]}, \texttt{[2, 2, 2, 2, 2, 2]}, \texttt{[6]}, \texttt{[12]}\} & Model architecture \\
        \hline
    \end{tabular}
    \caption{Varied hyperparameters and their configurations}
    \label{table:experiment_2_varied_hyperparameters}
\end{table}

\paragraph{Fixed hyperparameters}
Table~\ref{table:experiment_2_fixed_hyperparameters} contains all fixed hyperparameters in this experiment: 

\begin{table}[H]
    \centering
    \begin{tabular}{|c|c|c|p{5cm}|}
        \hline
        Hyperparameter & Type & Choices / Ranges & Description \\
        \hline
        \small\texttt{input\_size} & Integer & $13$ & Number of input features of the MLP \\
        \hline
        \small\texttt{output\_size} & Integer & $1$ & Number of labels of the MLP \\
        \hline
        \small\texttt{optimizer} & Categorical & \texttt{torch.optim.Adam} & Optimisation algorithm used for training \\
        \hline
        \small\texttt{learning\_rate} & Float & $0.001$ & Learning rate of the optimisation algorithm \\
        \hline
        \small\texttt{dropout\_rate} & Float & $0.5$ & Dropout rate applied to the data tensor after the input layer \\
        \hline
        \small\texttt{activation} & Categorical & ReLU & Activation function used in each dense layer of the MLP \\
        \hline
        \small\texttt{voter} & Categorical & \texttt{arithmetic\_mean} & Type of aggregator used to combine base learners' predictions \\
        \hline
        \small\texttt{batch\_size} & Integer & $32$ & Number of training samples passed to the model in one go \\
        \hline
    \end{tabular}
    \caption{Fixed hyperparameters and their configurations}
    \label{table:experiment_2_fixed_hyperparameters}
\end{table}

\paragraph{Metric}
Table~\ref{table:experiment_2_metric} contains all evaluation metrics collected in this experiment, which were stored in \path{experiment_runner_static_ncl_arithmetic_mean.csv}: 

\begin{table}[H]
    \centering
    \begin{tabular}{|c|c|c|}
        \hline
        Metric & Type & Description \\
        \hline
        \small\texttt{loss\_test} & Float & MSE of ensemble model during testing \\
        \hline
        \small\texttt{diversity\_coefficient\_test} & Float & $\rho$ during testing \\
        \hline
        \small\texttt{loss\_test\_base\_learner} & Float & Average MSE of base learner during testing \\
        \hline
    \end{tabular}
    \caption{Evaluation metrics}
    \label{table:experiment_2_metric}
\end{table}

\paragraph{Further notes}
To ensure comparison fairness, each ensemble model was evaluated exclusively against a control MLP that shared an identical architecture and training duration (epoch count). This ensured that any observed performance gains were due to the use of ensembling rather than individual base learner superiority. Furthermore, to maintain experimental focus, standard variance-reduction techniques such as \textit{bagging} were omitted from the training phase. Following the experiments, all base learner states were serialised and saved to \texttt{data/experiment/traditional/model/}. Refer to Appendix~\ref{appendix:diversity_coefficient} for the complete formula of the diversity coefficient $\rho$. 


\subsection{Expected Observations}


\subsubsection{Similar loss between ensemble base learners and the control MLP}
Each base learner within the ensemble shares identical hyperparameters with the control MLP, including layer depth and hidden node count. The only stochastic variable between the two is the specific data partitioning used during training. Consequently, the average MSE across all ensemble base learners is expected to converge to a value similar to the MSE of the standalone control MLP. 


\subsubsection{Ensemble MSE is always not higher than the average MSE of all base learners}
When MSE is used, it is mathematically proven that the ensemble loss has a maximum value of the average loss of all base learners. It is known as \textbf{ambiguity decomposition}. 

\begin{theorem}\label{theorem:ambiguity_decomposition}
    When MSE is used as the loss function, 
    \begin{align*}
        \underbrace{(\bar{y} - y)^2}_{\text{ensemble loss}} = \underbrace{\frac{1}{m} \sum_{i=1}^{m} (\hat{y}_i - y)^2}_{\text{average loss}} - \underbrace{\frac{1}{m} \sum_{i=1}^{m} (\hat{y}_i - \bar{y})^2}_{ambiguity}
    \end{align*}
\end{theorem}
Refer to Appendix~\ref{appendix:ambiguity_decomposition} for its proof. 
\\
\\
Since the ambiguity term is always non-negative, the ensemble loss is guaranteed to be not higher than the average loss at any time. For them to be equal, the ambiguity term must be $0$, which requires every base learner to have the same predictions, i.e. $\hat{y}_i = \bar{y}$ for every $i$th base learner. This ensemble model is no different that 1 single MLP. This is not likely to happen because every base learner is trained with a separate randomly-shuffled dataset so they have different weights even if other hyperparameters are set the same. As $m$ increases, the chance for $\hat{y}_i = \bar{y}}$ for $i \in \{1, 2, ..., m\}$ converges to 0. Thus, it is expected that the ensemble MSE is always lower than the average MSE of base learners in this experiment. 


\subsubsection{Increase in ensemble size does not lead to decrease in ensemble MSE}
Intuitively, increasing ensemble size should give a lower ensemble MSE. According to Theorem~\ref{theorem:ambiguity_decomposition}, ensemble loss is equal to average loss minus ambiguity. Furthermore, this is supported by the theorem \textbf{Bias-variance-diversity Decomposition}. 

\begin{theorem}\label{theorem:bias_variance_diversity_decomposition}
    For a set of models combined as an arithmetic mean $\hat{y}_{ensemble} = \frac{1}{m} \sum_{i=1}^{m} \hat{y}_i$, the expected risk of the ensemble admits the following decomposition (assume squared loss is used): 
    \begin{align*}
        \underbrace{\mathbb{E}_{\mathcal{S}}\left[\mathbb{E}_{\mathbf{x}y}\left[(y - \bar{y})^2\right]\right]}_{\text{expected ensemble risk}} = 
        \\
        &\phantom{=} \underbrace{\mathbb{E}_{\mathbf{x}y}\left[(y - \hat{y^*})^2\right]}_{\text{noise}} + \underbrace{\frac{1}{m} \sum_{i=1}^{m} \mathbb{E}_{\mathbf{x}y}\left[\left(\hat{y^*} - \mathbb{E}_{\mathcal{S}}[\hat{y}_i]\right)^2\right]}_{\text{average bias}} + \underbrace{\frac{1}{m} \sum_{i=1}^{m} \mathbb{E}_{\mathbf{x}y}\left[\mathbb{E}_{\mathcal{S}}\left[(\hat{y}_i - \mathbb{E}_{\mathcal{S}}[\hat{y}_i])^2\right]\right]}_{\text{average variance}}
        \\
        &\phantom{=}  - \underbrace{\mathbb{E}_{\mathbf{x}y}\left[\mathbb{E}_{\mathcal{S}}\left[\frac{1}{m} \sum_{i=1}^{m} (\hat{y}_i - \bar{y})^2\right]\right]}_{\text{diversity}}
    \end{align*}
\end{theorem}
Refer to Appendix~\ref{appendix:bias_variance_diversity_decomposition} for its proof. 
\\
\\
Note that the diversity term in Theorem~\ref{theorem:bias_variance_diversity_decomposition} can be thought as the expected value of the ambiguity term in Theorem~\ref{theorem:ambiguity_decomposition} over all training datasets $\mathcal{S}$ and all input data $\mathbf{x}$. Thus, even if average loss remains unchanged, when there are more base learners, it is more likely that the total distance between all base learners and the mean prediction increases because the chance that base learners make more diverse predictions due to different randomly-shuffled training datasets increases, leading to an increase in ambiguity and a decrease in ensemble loss. The expected risk of the ensemble is also expected to decrease when diversity increases. Thus, it is expected that, \textbf{intuitively}, increasing the ensemble size should reduce the ensemble MSE. 
\\
However, this is not expected in the current experiment settings. At this stage, base learners are not forced to be diverse. For instance, techniques like \textbf{bagging} have not been used. Instead, they are simply trained with different randomly-shuffled training datasets and they are not forced to learn different aspects of training data. In other words, the distance between each base learner and the mean prediction is approximately the same and increasing number of base learners $m$ does not increase the average of such distance (unless the new base learner makes an outlier prediction but this behaviour completely depends on the randomness of the training dataset). The ambiguity/diversity term can be thought as the average distance between base learner and the mean prediction. Therefore, the ambiguity/diversity term should remain unchanged when the ensemble size $m$ increases and so, ensemble MSE should remain approximately the same. The effect of diversity will be studied in Section~\ref{sec:stage_3}. 


\subsubsection{Better ensemble performance when base learners have better performance}
Increased model complexity, which is achieved through a higher number of hidden nodes or layers, typically improves performance by enabling the network to capture intricate, non-linear relationships within the data. While this progression is bounded by the risk of overfitting, complex models generally offer more accurate function approximations compared to their simpler counterparts. The current ensemble architecture utilises a simple arithmetic-mean aggregator. Base learners learn similar aspects of the training data due to the lack of NCL. Thus, the ensemble performance depends directly on the individual capabilities of its base learners. Consequently, an ensemble composed of these more sophisticated base learners is expected to yield higher predictive accuracy, effectively functioning as a ``committee of experts'' with greater collective specialised knowledge. 


\subsection{Evaluation and Empirical Comparisons}


\subsubsection{Ensemble base learners are just like the control MLP}
\paragraph{Experiments and Observations}
The average MSE of the ensemble base learners was similar to that of the control MLP as both shared identical architectures and hyperparameter configurations. The marginal variations in performance were primarily attributed to the stochastic nature of the training data partitioning for each model. This consistency was particularly visible in Figure~\ref{fig:traditional_mse_against_ensemble_size_2_16_fixed_architecture_2_2_2_2_2_2_fixed_epoch_num_200}, Figure~\ref{fig:traditional_mse_against_ensemble_size_2_16_fixed_architecture_2_2_2_2_2_2_fixed_epoch_num_400} and Figure~\ref{fig:traditional_mse_against_ensemble_size_2_16_fixed_architecture_6_fixed_epoch_num_400}. Even in instances where the control MLP's performance fell slightly outside the immediate range of the ensemble's average MSE, the absolute differences remained negligible. For example, in Figure~\ref{fig:traditional_mse_against_ensemble_size_2_16_fixed_architecture_12_fixed_epoch_num_400}, the control MLP achieved an MSE of approximately $0.04$ ($\text{RMSE} \approx 0.2$). In comparison, the ensemble base learners averaged an MSE of $0.06$ ($\text{RMSE} \approx 0.25$). Given that the target values were scaled within the range of $[-1, 1]$, the difference between an error of $0.20$ and $0.25$ was minimal, confirming that the base learners and the control model maintained approximately equivalent predictive capabilities. 

\paragraph{Conclusions from Experiment}
Despite the ensemble performance, each ensemble base learner has a similar performance as the control MLP. 


\subsubsection{Ensemble MSE is lower than the average MSE mostly and is never higher than the average MSE}
\paragraph{Experiments and Observations}
Consistent with the theoretical proofs in Theorem~\ref{theorem:ambiguity_decomposition} and Theorem~\ref{theorem:bias_variance_diversity_decomposition}, the experimental results demonstrated that ensemble loss remained consistently below the average loss of its corresponding base learners. Across all figures, the ensemble MSE (blue line) was lower than the average MSE (orange line), confirming the performance benefits of aggregation. In Figure~\ref{fig:traditional_mse_against_ensemble_size_2_16_fixed_architecture_2_2_2_fixed_epoch_num_400}, the blue data points were strictly below the orange points across the entire range of ensemble sizes. While certain configurations showed a convergence between the two metrics (such as in Figure~\ref{fig:traditional_mse_against_ensemble_size_2_16_fixed_architecture_2_2_2_2_2_2_fixed_epoch_num_400} at sizes $m \in \{4, 5, 6, 10, 15\}$), closer inspection revealed that the average MSE consistently maintained a higher value than the ensemble MSE. These observations empirically validated that the ensemble loss was bounded by the average loss of the base learners and never exceeded it. 

\paragraph{Conclusions from Experiment}
Ensemble MSE is lower than the average MSE most of the time. Furthermore, ensemble MSE is never higher than the average MSE. 


\subsubsection{How does ensemble size affect the diversity of predictions of base learners? }
\paragraph{Experiments and Observations}
Experimental results indicated that increasing the ensemble size did not significantly impact the diversity coefficient $\rho$. As illustrated in Figure~\ref{fig:traditional_diversity_coefficient_against_ensemble_size_2_16_fixed_epoch_num_400}, most traditional ensemble architectures maintained a stable $\rho$ value regardless of the number of base learners. The notable exception was the architecture consisting of three hidden layers with two nodes each (\texttt{[2, 2, 2]}). The high fluctuation in its diversity coefficient was consistent with the findings in Section~\ref{sec:stage_1} (Figure~\ref{fig:mlp_mse_against_epoch_num_10_1000_fixed_learning_rate_00001_hidden_size_2_2_2}), which identified this specific configuration as unsuitable for the task due to significant predictive instability. When multiple fluctuating base learners were combined, their poor individual predictive behaviors resulted in an unexpectedly large and unstable ensemble diversity coefficient. Similarly, the \texttt{[2, 2, 2, 2, 2, 2]} architecture (purple line) exhibited fluctuations of higher magnitude and a higher $\rho$ than simpler models (e.g. \texttt{[2]}, \texttt{[6]} etc.), further suggesting that deeper, narrower networks introduced stochastic noise that increased prediction diversity. 

\paragraph{Conclusions from Experiment}
When the right model architecture is picked, changing the ensemble size does not affect the diversity/ambiguity of predictions of base learners in this experiment due to the way these models are trained. 


\subsubsection{How does the complexity of base learner affect the ensemble performance? }
\paragraph{Experiments and Observations}
Ensemble performance was primarily driven by the individual predictive power of its base learners, which scaled with both architectural complexity and training duration. Comparisons across Figure~\ref{fig:traditional_mse_against_ensemble_size_2_16_fixed_architecture_2_fixed_epoch_num_200}, Figure~\ref{fig:traditional_mse_against_ensemble_size_2_16_fixed_architecture_6_fixed_epoch_num_200} and Figure~\ref{fig:traditional_mse_against_ensemble_size_2_16_fixed_architecture_12_fixed_epoch_num_200} confirmed that increasing the number of hidden nodes reduced the ensemble MSE. This improvement was a direct consequence of the more accurate function approximation capabilities of more complex base learners, as evidenced by the lower training and testing MSEs observed in Section~\ref{sec:stage_1} (Figure~\ref{fig:mlp_mse_against_epoch_num_10_1000_fixed_learning_rate_00001_hidden_size_6}, Figure~\ref{fig:mlp_mse_against_epoch_num_10_1000_fixed_learning_rate_00001_hidden_size_12} and Figure~\ref{fig:mlp_mse_against_epoch_num_10_1000_fixed_learning_rate_00001_hidden_size_2}). 

Extending the training duration similarly enhanced ensemble accuracy. For a single-layer architecture with two hidden nodes, increasing the training period from 200 to 400 epochs reduced the ensemble MSE from approximately $0.135$ to $0.125$ (Figure~\ref{fig:traditional_mse_against_ensemble_size_2_16_fixed_architecture_2_fixed_epoch_num_200} and Figure~\ref{fig:traditional_mse_against_ensemble_size_2_16_fixed_architecture_2_fixed_epoch_num_400}). While this trend remained visible in deeper architectures, the results were significantly more volatile. For instance, the three-layer \texttt{[2, 2, 2]} configuration, where MSE dropped from $0.175$ to $0.15$ over the same increase in epoch (Figure~\ref{fig:traditional_mse_against_ensemble_size_2_16_fixed_architecture_2_2_2_fixed_epoch_num_200} and Figure~\ref{fig:traditional_mse_against_ensemble_size_2_16_fixed_architecture_2_2_2_fixed_epoch_num_400}). The same phenomenon could also be seen in Figure~\ref{fig:traditional_mse_against_ensemble_size_2_16_fixed_architecture_2_2_2_2_2_2_fixed_epoch_num_200} and Figure~\ref{fig:traditional_mse_against_ensemble_size_2_16_fixed_architecture_2_2_2_2_2_2_fixed_epoch_num_400}. This persistent instability further validated the conclusion from Section~\ref{sec:stage_1}: for this specific dataset, increasing model complexity via layer width is fundamentally more effective and stable than increasing depth with a minimal node count per layer. 

\paragraph{Conclusions from Experiment}
In general, the higher the model complexity of base learner, the better the ensemble performance. However, this can concluded in a broader aspect. The better the performance of the base learners, the better the ensemble performance. This includes increasing number of training epochs and using the right model architecture. 


% ================================ Images ================================


% 2 Image
\begin{figure}[H]
    \centering

    \begin{minipage}[t]{0.48\textwidth}
        \vspace{0pt}
        \centering
        \includegraphics[width=\linewidth]{img/traditional/traditional_mse_against_ensemble_size_2_16_fixed_architecture_2_fixed_epoch_num_200.png}
        \caption{Impact of ensemble size on testing MSEs when the model with 200 epochs has a learning rate of 0.001 and 1 hidden layer of 2 hidden nodes}
        \label{fig:traditional_mse_against_ensemble_size_2_16_fixed_architecture_2_fixed_epoch_num_200}
    \end{minipage}
    \hfill
    \begin{minipage}[t]{0.48\textwidth}
        \vspace{0pt}
        \centering
        \includegraphics[width=\linewidth]{img/traditional/traditional_mse_against_ensemble_size_2_16_fixed_architecture_2_fixed_epoch_num_400.png}
        \caption{Impact of ensemble size on testing MSEs when the model trained with 400 epochs has a learning rate of 0.001 and 1 hidden layer of 2 hidden nodes}
        \label{fig:traditional_mse_against_ensemble_size_2_16_fixed_architecture_2_fixed_epoch_num_400}
    \end{minipage}
\end{figure}


% 2 Image
\begin{figure}[H]
    \centering

    \begin{minipage}[t]{0.48\textwidth}
        \vspace{0pt}
        \centering
        \includegraphics[width=\linewidth]{img/traditional/traditional_mse_against_ensemble_size_2_16_fixed_architecture_2_2_2_fixed_epoch_num_200.png}
        \caption{Impact of ensemble size on testing MSEs when the model with 200 epochs has a learning rate of 0.001 and 3 hidden layers of 2 hidden nodes}
        \label{fig:traditional_mse_against_ensemble_size_2_16_fixed_architecture_2_2_2_fixed_epoch_num_200}
    \end{minipage}
    \hfill
    \begin{minipage}[t]{0.48\textwidth}
        \vspace{0pt}
        \centering
        \includegraphics[width=\linewidth]{img/traditional/traditional_mse_against_ensemble_size_2_16_fixed_architecture_2_2_2_fixed_epoch_num_400.png}
        \caption{Impact of ensemble size on testing MSEs when the model trained with 400 epochs has a learning rate of 0.001 and 3 hidden layers of 2 hidden nodes}
        \label{fig:traditional_mse_against_ensemble_size_2_16_fixed_architecture_2_2_2_fixed_epoch_num_400}
    \end{minipage}
\end{figure}


% 2 Image
\begin{figure}[H]
    \centering

    \begin{minipage}[t]{0.48\textwidth}
        \vspace{0pt}
        \centering
        \includegraphics[width=\linewidth]{img/traditional/traditional_mse_against_ensemble_size_2_16_fixed_architecture_2_2_2_2_2_2_fixed_epoch_num_200.png}
        \caption{Impact of ensemble size on testing MSEs when the model with 200 epochs has a learning rate of 0.001 and 6 hidden layers of 2 hidden nodes}
        \label{fig:traditional_mse_against_ensemble_size_2_16_fixed_architecture_2_2_2_2_2_2_fixed_epoch_num_200}
    \end{minipage}
    \hfill
    \begin{minipage}[t]{0.48\textwidth}
        \vspace{0pt}
        \centering
        \includegraphics[width=\linewidth]{img/traditional/traditional_mse_against_ensemble_size_2_16_fixed_architecture_2_2_2_2_2_2_fixed_epoch_num_400.png}
        \caption{Impact of ensemble size on testing MSEs when the model trained with 400 epochs has a learning rate of 0.001 and 6 hidden layers of 2 hidden nodes}
        \label{fig:traditional_mse_against_ensemble_size_2_16_fixed_architecture_2_2_2_2_2_2_fixed_epoch_num_400}
    \end{minipage}
\end{figure}


% 2 Image
\begin{figure}[H]
    \centering

    \begin{minipage}[t]{0.48\textwidth}
        \vspace{0pt}
        \centering
        \includegraphics[width=\linewidth]{img/traditional/traditional_mse_against_ensemble_size_2_16_fixed_architecture_6_fixed_epoch_num_200.png}
        \caption{Impact of ensemble size on testing MSEs when the model with 200 epochs has a learning rate of 0.001 and 1 hidden layer of 6 hidden nodes}
        \label{fig:traditional_mse_against_ensemble_size_2_16_fixed_architecture_6_fixed_epoch_num_200}
    \end{minipage}
    \hfill
    \begin{minipage}[t]{0.48\textwidth}
        \vspace{0pt}
        \centering
        \includegraphics[width=\linewidth]{img/traditional/traditional_mse_against_ensemble_size_2_16_fixed_architecture_6_fixed_epoch_num_400.png}
        \caption{Impact of ensemble size on testing MSEs when the model trained with 400 epochs has a learning rate of 0.001 and 1 hidden layer of 6 hidden nodes}
        \label{fig:traditional_mse_against_ensemble_size_2_16_fixed_architecture_6_fixed_epoch_num_400}
    \end{minipage}
\end{figure}


% 2 Image
\begin{figure}[H]
    \centering

    \begin{minipage}[t]{0.48\textwidth}
        \vspace{0pt}
        \centering
        \includegraphics[width=\linewidth]{img/traditional/traditional_mse_against_ensemble_size_2_16_fixed_architecture_12_fixed_epoch_num_200.png}
        \caption{Impact of ensemble size on testing MSEs when the model with 200 epochs has a learning rate of 0.001 and 1 hidden layer of 12 hidden nodes}
        \label{fig:traditional_mse_against_ensemble_size_2_16_fixed_architecture_12_fixed_epoch_num_200}
    \end{minipage}
    \hfill
    \begin{minipage}[t]{0.48\textwidth}
        \vspace{0pt}
        \centering
        \includegraphics[width=\linewidth]{img/traditional/traditional_mse_against_ensemble_size_2_16_fixed_architecture_12_fixed_epoch_num_400.png}
        \caption{Impact of ensemble size on testing MSEs when the model trained with 400 epochs has a learning rate of 0.001 and 1 hidden layer of 12 hidden nodes}
        \label{fig:traditional_mse_against_ensemble_size_2_16_fixed_architecture_12_fixed_epoch_num_400}
    \end{minipage}
\end{figure}


% 1 Image
\begin{figure}[H]
    \centering
    \includegraphics[width=0.48\textwidth]{img/traditional/traditional_diversity_coefficient_against_ensemble_size_2_16_fixed_epoch_num_400.png}
    \caption{Impact of ensemble size on the diversity coefficient when all models were trained with 400 epochs}
    \label{fig:traditional_diversity_coefficient_against_ensemble_size_2_16_fixed_epoch_num_400}
\end{figure}


\section{Stage 3: Building a NCL Ensemble Model}
\label{sec:stage_3}


\subsection{Motivations}
Section~\ref{sec:stage_2} showed that diversity aids ensemble performance according to Theorem~\ref{theorem:ambiguity_decomposition} and Theorem~\ref{theorem:bias_variance_diversity_decomposition}. Section~\ref{subsec:background_ncl} also stated that techniques like bagging and boosting encourage diversity passively, giving an improvement in ensembel performance compared to a single MLP. Can we take the initiative to encourage diversity of base learners actively? Is there a way to penalise base learners for being similar to one another? This stage aims at finding a way to encourage diversity. However, when the attention on diversity is too much, the ensemble model suffers instead of benefits from diversity. Thus, one should strive for a balance between individual accuracy and diversity to maximise ensemble performance. In this stage, NCL is explored. 


\subsection{Design and Implementation}


\subsubsection{Reference experiment}
The experiment on NCL ensemble model was based on the NCL experiments on the Boston Housing dataset conducted in Brown et al.\ \cite{brown2005managing}. It aimed at reproducing similar experiment results to study the effect of correlation penalty coefficient $\lambda$ (represented as \textit{Gamma value} $\gamma$ in the paper Brown et al.\ \cite{brown2005managing}) on NCL ensemble loss. To extend the experiment, 3 extensions had been added to the NCL experiment in this work. In order to better compare the performance of NCL ensemble model and the original performance, MSE of a control MLP had been added. Although varying $\lambda$ did not affect the performance of the control MLP and the traditional ensemble model (from Section~\ref{sec:stage_2}), having a control MLP and a control traditional ensemble model enable the visualisation of the performance improvement of NCL ensemble models at different values of $\lambda$. Furthermore, as mentioned in Brown et al.\ \cite{brown2005managing}, the upper bound of $\lambda$ could be larger than $1$ if the ensemble size was small enough (e.g. $\lambda = 2$). Thus, the NCL experiment extended the range of $\lambda$ from $[0, 1]$ to $[0, 1.5]$, allowing us to study the NCL ensemble performance for $\lambda$ slightly beyond $1$. Lastly, there were 2 ways of varying the model complexity. In Brown et al.\ \cite{brown2005managing}, model complexity was increased by adding more hidden nodes to the neural network but each neural network always had 1 hidden layer. In this project, model could get more complex by either increasing the number of hidden nodes (then all models always had 1 hidden layer) or the number of hidden layers (then always 2 hidden nodes for each hidden layer). 


\subsubsection{Datasets}
The dataset was shuffled and 20\% of data was kept as the testing set and the remaining 80\% of data was for training and validation. 80\% of the remaining data was for training and the remaining was for validation. With the Boston Housing dataset, this equated to 303 for training, 101 for validation and 102 for testing. The validation data was used to perform early stopping by the following procedure: 

\begin{itemize}
    \item Train the model while noting the validation error every 50 epochs. 
    \item If the validation error has risen in comparison to 500 epochs ago, terminate training and reset weights to the best point within that 500-epoch window (at a resolution of 50 epochs) according to the validation data. 
\end{itemize}

In the current experiments, such reset never happened. 


\subsubsection{Model construction}
The architectural framework remained consistent with Section~\ref{sec:stage_2}, with the primary modification residing in the base learners' loss function. To explicitly encourage model divergence, the standard MSE was augmented with a correlation-penalty term that discouraged similarity between individual learners. A balancing coefficient, $\lambda$ (also known as \textit{correlation penalty coefficient}), was introduced to manage the trade-off between individual accuracy and ensemble diversity. 

The complete mathematical derivation of this NCL objective is provided in Appendix~\ref{appendix:loss_function_ncl}. During backpropagation, the optimisation process simultaneously minimised two key components: 

\begin{enumerate}
    \item \textbf{Individual Error} \\
    Minimising $(\hat{y}_i - y)^2$ ensures each base learner maintains high predictive accuracy relative to the ground truth. 

    \item \textbf{Collective Correlation} \\
    Minimising the penalty term $(\hat{y}_i - \bar{y})^2$ forces the individual learner's prediction to deviate from the ensemble mean $\bar{y}$ and therefore, maximising functional diversity across the population. 
\end{enumerate}


\subsubsection{Experiment}
Experiments had been conducted to mainly compare the performance of a NCL ensemble model and that of a traditional ensemble model (the performance of a MLP was however also included for reference). The experiment settings were identical to those in Section~\ref{sec:stage_2} except a NCL ensemble model was added to the experiment too. In the NCL ensemble model, all base learners were trained using NCL where a correlation-penalty term was added to the loss functions of base learners. 

\paragraph{Varied hyperparameters}
Table~\ref{table:experiment_3_varied_hyperparameters} contains all varied hyperparameters in this experiment: 

\begin{table}[H]
    \centering
    \begin{tabular}{|c|c|p{5cm}|c|}
        \hline
        Hyperparameter & Type & Choices / Ranges & Description \\
        \hline
        \small\texttt{correlation\_penalty\_coefficients} & Float & $[0, 1.5]\text{, step}=0.1$ & $\lambda$ \\
        \hline
        \small\texttt{ensemble\_sizes} & Integer & $[2, 16]\text{, step}=1$ & $m$ \\
        \hline
        \small\texttt{hidden\_sizes} & Python list & \{\texttt{[2]}, \texttt{[2, 2, 2]}, \texttt{[2, 2, 2, 2, 2, 2]}, \texttt{[6]}, \texttt{[12]}\} & Model architecture \\
        \hline
    \end{tabular}
    \caption{Varied hyperparameters and their configurations}
    \label{table:experiment_3_varied_hyperparameters}
\end{table}

\paragraph{Fixed hyperparameters}
Table~\ref{table:experiment_3_fixed_hyperparameters} contains all fixed hyperparameters in this experiment: 

\begin{table}[H]
    \centering
    \begin{tabular}{|c|c|c|p{5cm}|}
        \hline
        Hyperparameter & Type & Choices / Ranges & Description \\
        \hline
        \small\texttt{input\_size} & Integer & $13$ & Number of input features of the MLP \\
        \hline
        \small\texttt{output\_size} & Integer & $1$ & Number of labels of the MLP \\
        \hline
        \small\texttt{optimizer} & Categorical & \texttt{torch.optim.Adam} & Optimisation algorithm used for training \\
        \hline
        \small\texttt{learning\_rate} & Float & $0.001$ & Learning rate of the optimisation algorithm \\
        \hline
        \small\texttt{dropout\_rate} & Float & $0.5$ & Dropout rate applied to the data tensor after the input layer \\
        \hline
        \small\texttt{activation} & Categorical & ReLU & Activation function used in each dense layer of the MLP \\
        \hline
        \small\texttt{epoch\_nums} & Integer & $400$ & Number of training epochs \\
        \hline
        \small\texttt{voter} & Categorical & \texttt{arithmetic\_mean} & Type of aggregator used to combine base learners' predictions \\
        \hline
        \small\texttt{batch\_size} & Integer & $32$ & Number of training samples passed to the model in one go \\
        \hline
    \end{tabular}
    \caption{Fixed hyperparameters and their configurations}
    \label{table:experiment_3_fixed_hyperparameters}
\end{table}

\paragraph{Metric}
Table~\ref{table:experiment_3_metric} contains all evaluation metrics collected in this experiment, which were stored in \path{experiment_runner_static_ncl_arithmetic_mean.csv}: 

\begin{table}[H]
    \centering
    \begin{tabular}{|c|c|c|}
        \hline
        Metric & Type & Description \\
        \hline
        \small\texttt{loss\_test} & Float & MSE of ensemble model during testing \\
        \hline
        \small\texttt{diversity\_coefficient\_test} & Float & $\rho$ during testing \\
        \hline
    \end{tabular}
    \caption{Evaluation metrics}
    \label{table:experiment_3_metric}
\end{table}

\paragraph{Further notes}
To ensure comparison fairness, each NCL ensemble model was evaluated exclusively against a control MLP that shared an identical architecture and training duration (epoch count). This ensured that any observed performance gains were due to the diversity, inducing penalty term rather than individual architectural advantages. 
Upon completion of the experiments, all constituent base learner states were serialised and saved to \texttt{data/experiment/static-ncl/model/}. 
For the complete formula of the diversity coefficient $\rho$ used to evaluate these models, refer to Appendix~\ref{appendix:diversity_coefficient}. 


\subsubsection{Viewing Experiments Better}
There are 2 categories for the figures below. One category provides comparisons of the impacts of correlation penalty coefficient on MSE and diversity coefficient across particular ensemble models with a \textbf{fixed model architecture}. The other category provides comparisons of the impacts of correlation penalty coefficient on MSE and diversity coefficient across particular ensemble models with a \textbf{fixed ensemble size}. Furthermore, when $\lambda$ exceeds an upper bound, ensemble MSE and diversity coefficient increase exponentially. If $\lambda$ in a graph has range $[0, 1.5]$, it will be difficult to compare the MSE and diversity coefficient for smaller values of $\lambda$. Thus, $\lambda$ in graphs either have range $[0, 1]$ or $[0.8, 1.5]$ so that MSE and diversity coefficient comparisons can be viewed more clearly. 


\subsection{Expected Observations}


\subsubsection{Ensemble performance changes as diversity of base learners changes}
The complete loss function of NCL is given below: 

\begin{theorem}\label{theorem:loss_function_ncl}
    \begin{align*}
        \underbrace{\mathcal{L}_i}_{\text{ith base learner loss}} = \underbrace{(\hat{y}_i - y)^2}_{\text{squared loss of 1 data point}} - \lambda \underbrace{(\hat{y}_i - \bar{y})^2}_{\text{correlation penalty function}}
    \end{align*}
\end{theorem}
Note that this is similar to the loss function used for the ensemble model in Section~\ref{sec:stage_2} except it has a penalty term called \textit{correlation penalty function}. It has a high value when base learners are positively correlated, where base learners' predictions are close to one another and have values close to $\bar{y}$. Positively correlated base learners have a small value of $(\hat{y}_i - \bar{y})^2$. Minimising this loss function can be thought as minimising $(\hat{y}_i - y)^2$ and maximising $(\hat{y}_i - \bar{y})^2$. Minimising MSE promotes accuracy of the current base learner while maximising the correlation penalty function penalises the current base learner from being similar to other base learners. The value of $\lambda$ is usually set between $[0, 1]$ but as mentioned in Brown et al.\ \cite{brown2005managing}, its upper bound, given by the formula in Appendix~\ref{appendix:lambda_upper_bound}, can sometimes be larger than $1$, e.g. when $m = 2$. It controls the balance between individual base learner accuracy and diversity between all base learners. Initially, when $\lambda = 0$, the loss function becomes $(\hat{y}_i - y)^2$. This loss function is for 1 data point. In the experiment, $(\hat{y}_i - y)^2$ is averaged across all data points, making it the same as the MSE used in Section~\ref{sec:stage_2}. This is no different from training all base learners independently and thus, it is expected that when $\lambda = 0$, a NCL ensemble model should have similar performance as the ensemble model in Section~\ref{sec:stage_2}. When $\lambda$ is increased, base learners are forced to be diverse. According to Theorem~\ref{theorem:ambiguity_decomposition} and Theorem~\ref{theorem:bias_variance_diversity_decomposition}, when the ambiguity/diversity term increases, the ensemble loss/risk should decrease. Thus, it is expected that the ensemble model has an improvement in performance when $\lambda$ is increased initially. According to Theorem~\ref{theorem:loss_function_ncl_lambda_1}, a special case happens when $\lambda = 1$. 

\begin{theorem}\label{theorem:loss_function_ncl_lambda_1}
    For ith base learner, when MSE is used and the ensemble model is an arithmetic-mean-aggregator, the average loss of the ith base learner across all $n$ data points is given below: 
    \begin{align*}
        \underbrace{\mathcal{L}_i}_{\text{ith base learner loss}} = \frac{1}{n} \sum^{n}_{i = 1} \underbrace{(\bar{y} - y)^2}_{\text{squared loss of 1 data point}}
    \end{align*}
\end{theorem}
Refer to Appendix~\ref{appendix:loss_function_ncl_lambda_1} for its proof. 
\\
Note that the right-hand side of the equation does not contain any terms that are dependent on ith base learner. This means when $\lambda = 1$, the ensemble model can be treated as 1 single estimator, i.e. $\hat{y}_{ensemble} = \bar{y}$. Thus, it is expected that the ensemble model should have a similar performance as a single MLP when $\lambda = 1$. Moreover, if $\lambda$ continues to increase, the focus on diversity should take over the focus on accuracy. This is similar to how a parliament is not efficient when every committee member is disagreeing only for the sake of it (without concerning correctness). Therefore, it is expected that the ensemble model should have a performance recession when $\lambda$ becomes too large, i.e. larger than $\lambda_{upper\_bound}$. 


\subsubsection{Higher diversity of base learners when NCL is used}
The correlation penalty function, as defined in Theorem~\ref{theorem:loss_function_ncl}, systematically penalises base learners for exhibiting high similarity and positive correlation. By prioritising the development of negatively-correlated learners, the loss function encourages the population to explore a broader range of the hypothesis space, resulting in more diverse individual predictions. Consequently, $\rho$ is expected to scale proportionally with the penalty coefficient $\lambda$. As the value of $\lambda$ increases, the optimisation process places greater weight on the penalty term $(\hat{y}_i - \bar{y})^2$, forcing base learners to deviate further from the ensemble mean. This theoretical expectation suggests that higher $\lambda$ values will directly result in a more heterogeneous committee of learners, maximising the \textit{ambiguity} term in Theorem~\ref{theorem:ambiguity_decomposition}. 


\subsubsection{Simpler base learners are better than more complex base learners}
The introduction of the NCL penalty alters the traditional relationship between model complexity and performance. While Section~\ref{sec:stage_2} established that increased complexity typically improves individual accuracy, the inclusion of a diversity constraint leads a trade-off during optimisation. For high-complexity base learners, the pressure to minimise the penalty term $(\hat{y}_i - \bar{y})^2$ can become counterproductive. As these ``senior'' models already possess high individual accuracies, forcing them to deviate from one another often requires them to sacrifice their optimal function approximations. This results in a ``committee'' where only a few members remain near the ground truth while others are pushed to make wrong predictions to satisfy the diversity requirement, ultimately degrading ensemble performance. 
Conversely, simpler base learners (analogous to junior developers) initially possess higher individual bias and lower accuracy. When the NCL penalty encourages these models to explore different aspects of the data structure, it effectively prevents them from all converging on the same suboptimal local minima. This forced exploration increases the probability that specific learners will randomly capture unique features of the underlying data, allowing the aggregated ensemble prediction to outperform any single constituent. This phenomenon suggests that diversity-focused training is most beneficial when the base learners lack sufficient individual capacity as it compensates for high individual bias through a broader, non-redundant search of the solution space. 


\subsection{Evaluation and Empirical Comparisons}


\subsubsection{What is the effect of correlation penalty coefficient?}
\paragraph{Experiments and Observations}
The experimental results demonstrated a distinct non-linear relationship between the correlation penalty coefficient $\lambda$ and ensemble performance. When $\lambda = 0$, the NCL ensemble behaved identically to a traditional ensemble of independent learners as there was no penalty for model similarity. However, as $\lambda$ increased toward a threshold of approximately $0.7$, a visible reduction in ensemble MSE occurred (as seen in Figure~\ref{fig:static_ncl_mse_against_correlation_penalty_coefficient_0_1_fixed_architecture_2}, Figure~\ref{fig:static_ncl_mse_against_correlation_penalty_coefficient_0_1_fixed_architecture_6}, Figure~\ref{fig:static_ncl_mse_against_correlation_penalty_coefficient_0_1_fixed_architecture_12} Figure~\ref{fig:static_ncl_mse_against_correlation_penalty_coefficient_0_1_fixed_ensemble_size_6_architecture_6_comparator_traditional}). This improvement coincided with a slight increase in the diversity coefficient $\rho$ (Figure~\ref{fig:static_ncl_diversity_coefficient_against_correlation_penalty_coefficient_0_1_fixed_architecture_2}, Figure~\ref{fig:static_ncl_diversity_coefficient_against_correlation_penalty_coefficient_0_1_fixed_architecture_6}, Figure~\ref{fig:static_ncl_diversity_coefficient_against_correlation_penalty_coefficient_0_1_fixed_architecture_12} and Figure~\ref{fig:static_ncl_diversity_coefficient_against_correlation_penalty_coefficient_0_1_fixed_ensemble_size_6_architecture_6}), suggesting that a moderate degree of functional diversity effectively reduced the collective error. 

Beyond the optimal threshold ($\lambda \approx 0.6$ to $0.8$), the ensemble became over-diversified. Figure~\ref{fig:static_ncl_diversity_coefficient_against_correlation_penalty_coefficient_0_1_fixed_ensemble_size_6_architecture_6} and Figure~\ref{fig:static_ncl_diversity_coefficient_against_correlation_penalty_coefficient_0_1_fixed_ensemble_size_6_architecture_2_2_2} showed an exponential spike in $\rho$, which triggered a corresponding surge in ensemble MSE. While deeper architectures (Figure~\ref{fig:static_ncl_mse_against_correlation_penalty_coefficient_0_1_fixed_architecture_2_2_2} and Figure~\ref{fig:static_ncl_mse_against_correlation_penalty_coefficient_0_1_fixed_architecture_2_2_2_2_2_2}) exhibited more stochastic fluctuations, they followed the same trend: performance remained stable until $\lambda \approx 0.9$, after which MSEs began to peak (e.g. traditional ensemble models of architecture \texttt{[6]}). This suggested that forcing excessive divergence compelled base learners to prioritise being ``different'' over being ``accurate'' ultimately degrading the ensemble's predictive ability. 

Interestingly, the data validated the theoretical behavior of NCL when $\lambda = 1$. In Figure~\ref{fig:static_ncl_mse_against_correlation_penalty_coefficient_0_1_fixed_architecture_2} and Figure~\ref{fig:static_ncl_mse_against_correlation_penalty_coefficient_0_1_fixed_architecture_6}, the ensemble MSE at $\lambda = 1$ converged toward the MSE of the single control MLP. This empirically supported the principle that when $\lambda = 1$, the ensemble mathematically collapsed into a single estimator. The minor difference between the ensemble and control MSEs at this point were likely due to the variance introduced by different random seeds and data shuffling during the training process of individual models. 

\paragraph{Conclusions from Experiment}
The correlation penalty coefficient $\lambda$ affects the diversity of predictions of base learners. Increasing it causes base learners to be diverse, which usually gives a more accurate ensemble prediction. Once the coefficient reaches its upper bound, base learners drastically become very diverse that the accuracy of ensemble prediction decreases. Thus, the value of $\lambda$ should be right below its upper bound to maximise ensemble prediction accuracy. In this case, setting $\lambda = 0.7$ seems to give a good ensemble performance. 


\subsubsection{What is the effect of network complexity?}
\paragraph{Experiments and Observations}
A critical finding in the NCL experiments was the transition of ``stronger'' base learners from complex to simpler architectures as the penalty coefficient $\lambda$ increased. In Section~\ref{sec:stage_2}, more complex models consistently dominated. However, the introduction of the diversity constraint $p_i$ reversed this trend. As shown in Figure~\ref{fig:static_ncl_mse_against_correlation_penalty_coefficient_0_1_fixed_ensemble_size_6_architecture_6_comparator_traditional}, while complex 12-node base learners were superior at low penalty values ($\lambda < 0.6$), they suffered a significant performance recession as $\lambda$ approached $0.8$. At this higher threshold, the 2-node ensemble, which was previously the worst performer, became the most accurate while the 12-node ensemble became the least effective. This phenomenon suggested a high-efficiency advantage: with NCL and a correctly tuned $\lambda$, competitive performance can be achieved using simpler architectures, significantly reducing the computational resources and time required for training. This phenomenon is supported by existing literature. Brown et al. \cite{brown2005managing} observed a similar shift in their Figure~14, where simpler ensembles (2--6 base learners) outperformed larger ones (10--12 learners) once the penalty coefficient $\gamma$ exceeded $0.5$. Likewise, Eastwood et al. \cite{eastwood2006lambda} demonstrated that on synthetic datasets (the bottom-right graph of Figure~6), the optimal architecture shifted from 20 hidden nodes to 10 nodes as $\lambda$ increased from $0$ to $1$. 

Furthermore, the data revealed an inverse relationship between base learner complexity and the optimal correlation penalty coefficient ($\lambda^*$). As architectural complexity increased, $\lambda^*$ decreased: 

\begin{enumerate}
    \item \textbf{2-node model (Simple)} \\
    $\lambda^* \approx 0.8$ (from Figure~\ref{fig:static_ncl_mse_against_correlation_penalty_coefficient_0_1_fixed_architecture_2})

    \item \textbf{6-node model (Medium)} \\
    $\lambda^* \approx 0.7$ (from Figure~\ref{fig:static_ncl_mse_against_correlation_penalty_coefficient_0_1_fixed_architecture_6})

    \item \textbf{12-node model (Complex)} \\
    $\lambda^* \approx 0.5$ (from Figure~\ref{fig:static_ncl_mse_against_correlation_penalty_coefficient_0_1_fixed_architecture_12})
\end{enumerate}

This suggested that complex models, which were already highly specialised, were more sensitive to the NCL penalty whereas simpler models had more ``room'' to explore diverse search spaces without sacrificing their fundamental predictive capabilities. 

\paragraph{Conclusions from Experiment}
As we increase $\lambda$, we should shift from using complex base learners to simple base learners in order to get a better ensemble performance. Ensemble model of simpler base learners have a higher optimal correlation penalty coefficient. 


\subsubsection{What is the effect of ensemble size?}
\paragraph{Experiments and Observations}
A larger ensemble size appeared to dampen the stochastic fluctuations typically associated with NCL training. As observed in Figure~\ref{fig:static_ncl_mse_against_correlation_penalty_coefficient_0_1_fixed_architecture_2}, Figure~\ref{fig:static_ncl_mse_against_correlation_penalty_coefficient_0_1_fixed_architecture_2_2_2} and Figure~\ref{fig:static_ncl_mse_against_correlation_penalty_coefficient_0_1_fixed_architecture_2_2_2_2_2_2}, ensemble models with only 2 or 6 base learners exhibited high sensitivity to changes in $\lambda$, leading to volatile MSE curves. In contrast, ensembles with 12 base learners showed a more consistent and smooth performance trajectory as $\lambda$ increased. This implies that a higher ensemble size $m$ provided a ``buffer'' against the erratic individual weight adjustments forced by the diversity penalty $p_i$, leading to more reliable collective convergence. 

The same magnitude change in $\lambda$ scaled the diversity of the base learners' predictions more when the ensemble size was larger. While initial peaks in diversity appeared across all architectures once $\lambda > 0.6$, a dramatic exponential increase occurred after $\lambda = 0.9$ (Figure~\ref{fig:static_ncl_diversity_coefficient_against_correlation_penalty_coefficient_0_1_fixed_architecture_2}, Figure~\ref{fig:static_ncl_diversity_coefficient_against_correlation_penalty_coefficient_0_1_fixed_architecture_6} and Figure~\ref{fig:static_ncl_diversity_coefficient_against_correlation_penalty_coefficient_0_1_fixed_architecture_12}). From Figure~\ref{fig:static_ncl_diversity_coefficient_against_correlation_penalty_coefficient_0_1_fixed_architecture_2}, Figure~\ref{fig:static_ncl_diversity_coefficient_against_correlation_penalty_coefficient_0_1_fixed_architecture_12} and Figure~\ref{fig:static_ncl_diversity_coefficient_against_correlation_penalty_coefficient_0_1_fixed_architecture_2_2_2_2_2_2}, for any fixed increase in $\lambda$, the 12-base-learner ensemble consistently produced a significantly larger surge in $\rho$ compared to the 2-base-learner or 6-base-learner configurations. Mathematically, this aligned with the diversity coefficient formula in Appendix~\ref{appendix:diversity_coefficient}, where a larger ensemble size amplified the average distance between individual base learner predictions ($\hat{y}_i$) and the ensemble mean ($\bar{y}$). Essentially, in a larger population, the NCL penalty term $(\hat{y}_i - \bar{y})^2$ exerted a more powerful repulsive force as each learner must become more different to maintain uniqueness. This suggested that as ensemble size grew, the system became increasingly sensitive to the diversity penalty. 

\paragraph{Conclusions from Experiment}
When the ensemble size is larger, increasing $\lambda$ improves the ensemble performance (before reaching the upper bound of $\lambda$) more consistently. Furthermore, increasing $\lambda$ leads to a greater diversity of base learners when the ensemble size is larger. 


\subsubsection{Should we always maximise diversity of base learners?}
\paragraph{Experiments and Observations}
When base learners were trained independently ($\lambda = 0$), architectural complexity is the primary determinant of whether an ensemble outperformed a single MLP. As illustrated in Figure~\ref{fig:static_ncl_mse_against_correlation_penalty_coefficient_0_1_fixed_ensemble_size_6_architecture_6_comparator_mlp}, ensembles of complex base learners naturally surpassed the performance of a control MLP of moderate complexity. In contrast, ensembles composed of simple base learners initially performed worse than the control MLP. Without a diversity constraint, these simple learners converged toward similar, inaccurate predictions. Since they were functionally redundant and individually limited, their averaged prediction remained inferior to a single, more sophisticated model. 

As $\lambda$ increased, the introduction of penalty term allowed both simple and complex ensembles to improve and eventually outperformed the control MLP. However, this benefit was highly sensitive to the magnitude of the penalty. When diversity was over-emphasised ($\lambda > 0.6$), the performance of complex ensembles degraded rapidly, eventually falling behind the accuracy of a single simple MLP. Interestingly, as shown in Figure~\ref{fig:static_ncl_mse_against_correlation_penalty_coefficient_0_1_fixed_ensemble_size_6_architecture_6_comparator_mlp}, simple ensembles had similar performance as their corresponding control MLPs up to $\lambda = 1$. However, Figure~\ref{fig:static_ncl_mse_against_correlation_penalty_coefficient_08_15_fixed_ensemble_size_6_architecture_6_comparator_mlp} demonstrated that once $\lambda$ exceeded $1$, even these simple ensembles were outperformed by a single MLP. This occurred because the NCL objective began to prioritise the maximisation of the distance between predictions $(\hat{y}_i - \bar{y})^2$ over the minimisation of the individual error $(\hat{y}_i - y)^2$. At this stage of over-diversification, base learners actively avoided the ground truth in favor of maintaining uniqueness, resulting in a collective mean that significantly deviated from the true value. 

\paragraph{Conclusions from Experiment}
A higher diversity does not necessarily guarantee an improvement in performance and model selection should not only aim at maximising diversity when constructing an ensemble model. Instead, a balance should be striven so that base learners are somewhat diverse but not too diverse. 


\subsubsection{Exponential ensemble performance recession when $\lambda$ exceeds its upper bound}
\paragraph{Experiments and Observations}
The relationship between the ensemble size ($m$) and the penalty coefficient ($\lambda$) is governed by the theoretical limits derived in Appendix~\ref{appendix:lambda_upper_bound}. According to the formula, the upper bound for the penalty coefficient decreases as the population grows: 

\begin{enumerate}
    \item \textbf{$m = 2$} \\
    $\lambda_{upper\_bound} = 2.0$

    \item \textbf{$m = 6$} \\
    $\lambda_{upper\_bound} = 0.72$

    \item \textbf{$m = 12$} \\
    $\lambda_{upper\_bound} \approx 0.595$
\end{enumerate}

Experimental observations in Figure~\ref{fig:static_ncl_mse_against_correlation_penalty_coefficient_08_15_fixed_architecture_2}, Figure~\ref{fig:static_ncl_mse_against_correlation_penalty_coefficient_08_15_fixed_architecture_6} and Figure~\ref{fig:static_ncl_mse_against_correlation_penalty_coefficient_08_15_fixed_architecture_12} confirmed that when $\lambda$ exceeded these thresholds (particularly as it moved beyond $1.0$), the ensemble MSE surged by approximately 30 times. This catastrophic performance degradation was directly linked to the exponential increase in diversity of the base learners. As evidenced by Figure~\ref{fig:static_ncl_diversity_coefficient_against_correlation_penalty_coefficient_08_15_fixed_architecture_2}, Figure~\ref{fig:static_ncl_diversity_coefficient_against_correlation_penalty_coefficient_08_15_fixed_architecture_6} and Figure~\ref{fig:static_ncl_diversity_coefficient_against_correlation_penalty_coefficient_08_15_fixed_architecture_12}, the diversity coefficient $\rho$ underwent an exponential increase once these bounds were breached. This indicated that the repulsive force of the NCL penalty became the dominant component of the loss function, overshadowing the accuracy term. This resulted in a set of predictions that were too scattered to yield a meaningful or accurate ensemble mean. Thus, $\lambda$ should be kept below its upper bound in order to utilise the benefits of NCL. 

\paragraph{Conclusions from Experiment}
When $\lambda$ exceeds its upper bound given by the formula in Appendix~\ref{appendix:lambda_upper_bound}, ensemble MSE increases exponentially due to exponential increase in diversity coefficient. The ensemble model becomes useless due to its inaccuracies. Thus, $\lambda$ should be kept below its upper bound in order to utilise the benefits of NCL. 


\subsubsection{Descent in diversity coefficient}
\paragraph{Experiments and Observations}
An unexpected phenomenon observed in the experimental data was the sudden decrease of diversity coefficient $\rho$ when the correlation penalty coefficient $\lambda$ reached extreme values. As illustrated in Figure~\ref{fig:static_ncl_diversity_coefficient_against_correlation_penalty_coefficient_08_15_fixed_architecture_2}, Figure~\ref{fig:static_ncl_diversity_coefficient_against_correlation_penalty_coefficient_08_15_fixed_architecture_6} and Figure~\ref{fig:static_ncl_diversity_coefficient_against_correlation_penalty_coefficient_08_15_fixed_architecture_12}, once $\lambda$ exceeded a critical threshold, the diversity coefficient $\rho$ decreased drastically, dropping from peaks near $2.5$ to values as low as $0.25$. The magnitude and timing of this descent were strongly moderated by the ensemble size $m$: 

\begin{enumerate}
    \item \textbf{Sensitivity To Ensemble Size} \\
    The drop was more significant in larger ensembles. When $m = 12$, the decrease of diversity was more severe than in smaller ensemble sizes. 

    \item \textbf{Shifted Threshold} \\
    Larger ensembles appeared more resilient to this collapse, as the ``breaking point'' occurred at higher $\lambda$ values. For instance, in Figure~\ref{fig:static_ncl_diversity_coefficient_against_correlation_penalty_coefficient_08_15_fixed_architecture_2}, the diversity decrease triggered at $\lambda = 1.2$ for an ensemble size of $6$ but was delayed until $\lambda = 1.4$ for an ensemble size of $12$. 
\end{enumerate}

\paragraph{Conclusions from Experiment}
When $\lambda$ becomes extremely large, $\rho$ decreases drastically. When ensemble size is larger, such drop occurs at a larger magnitude and at a larger value of $\lambda$. 


% ================================ Images ================================


% 2 Image
\begin{figure}[H]
    \centering

    \begin{minipage}[t]{0.48\textwidth}
        \vspace{0pt}
        \centering
        \includegraphics[width=\linewidth]{img/static-ncl/static_ncl_mse_against_correlation_penalty_coefficient_0_1_fixed_architecture_2.png}
        \caption{Impact of correlation penalty coefficient (in the range $[0, 1]$) on MSE. Base learners had 1 hidden layer of 2 hidden nodes and were trained with 400 epochs. They increased their model complexities by increasing number of hidden nodes in the hidden layer. }
        \label{fig:static_ncl_mse_against_correlation_penalty_coefficient_0_1_fixed_architecture_2}
    \end{minipage}
    \hfill
    \begin{minipage}[t]{0.48\textwidth}
        \vspace{0pt}
        \centering
        \includegraphics[width=\linewidth]{img/static-ncl/static_ncl_mse_against_correlation_penalty_coefficient_08_15_fixed_architecture_2.png}
        \caption{Impact of correlation penalty coefficient (in the range $[0.8, 1.5]$) on MSE. Base learners had 1 hidden layer of 2 hidden nodes and were trained with 400 epochs. They increased their model complexities by increasing number of hidden nodes in the hidden layer. }
        \label{fig:static_ncl_mse_against_correlation_penalty_coefficient_08_15_fixed_architecture_2}
    \end{minipage}
\end{figure}


% 2 Image
\begin{figure}[H]
    \centering

    \begin{minipage}[t]{0.48\textwidth}
        \vspace{0pt}
        \centering
        \includegraphics[width=\linewidth]{img/static-ncl/static_ncl_mse_against_correlation_penalty_coefficient_0_1_fixed_architecture_6.png}
        \caption{Impact of correlation penalty coefficient (in the range $[0, 1]$) on MSE. Base learners had 1 hidden layer of 6 hidden nodes and were trained with 400 epochs. They increased their model complexities by increasing number of hidden nodes in the hidden layer. }
        \label{fig:static_ncl_mse_against_correlation_penalty_coefficient_0_1_fixed_architecture_6}
    \end{minipage}
    \hfill
    \begin{minipage}[t]{0.48\textwidth}
        \vspace{0pt}
        \centering
        \includegraphics[width=\linewidth]{img/static-ncl/static_ncl_mse_against_correlation_penalty_coefficient_08_15_fixed_architecture_6.png}
        \caption{Impact of correlation penalty coefficient (in the range $[0.8, 1.5]$) on MSE. Base learners had 1 hidden layer of 6 hidden nodes and were trained with 400 epochs. They increased their model complexities by increasing number of hidden nodes in the hidden layer. }
        \label{fig:static_ncl_mse_against_correlation_penalty_coefficient_08_15_fixed_architecture_6}
    \end{minipage}
\end{figure}


% 2 Image
\begin{figure}[H]
    \centering

    \begin{minipage}[t]{0.48\textwidth}
        \vspace{0pt}
        \centering
        \includegraphics[width=\linewidth]{img/static-ncl/static_ncl_mse_against_correlation_penalty_coefficient_0_1_fixed_architecture_12.png}
        \caption{Impact of correlation penalty coefficient (in the range $[0, 1]$) on MSE. Base learners had 1 hidden layer of 12 hidden nodes and were trained with 400 epochs. They increased their model complexities by increasing number of hidden nodes in the hidden layer. }
        \label{fig:static_ncl_mse_against_correlation_penalty_coefficient_0_1_fixed_architecture_12}
    \end{minipage}
    \hfill
    \begin{minipage}[t]{0.48\textwidth}
        \vspace{0pt}
        \centering
        \includegraphics[width=\linewidth]{img/static-ncl/static_ncl_mse_against_correlation_penalty_coefficient_08_15_fixed_architecture_12.png}
        \caption{Impact of correlation penalty coefficient (in the range $[0.8, 1.5]$) on MSE. Base learners had 1 hidden layer of 12 hidden nodes and were trained with 400 epochs. They increased their model complexities by increasing number of hidden nodes in the hidden layer. }
        \label{fig:static_ncl_mse_against_correlation_penalty_coefficient_08_15_fixed_architecture_12}
    \end{minipage}
\end{figure}


% 2 Image
\begin{figure}[H]
    \centering

    \begin{minipage}[t]{0.48\textwidth}
        \vspace{0pt}
        \centering
        \includegraphics[width=\linewidth]{img/static-ncl/static_ncl_mse_against_correlation_penalty_coefficient_0_1_fixed_architecture_2_2_2.png}
        \caption{Impact of correlation penalty coefficient (in the range $[0, 1]$) on MSE. Base learners had 3 hidden layers of 2 hidden nodes and were trained with 400 epochs. They increased their model complexities by increasing number of hidden layers (each layer always has 2 hidden nodes). }
        \label{fig:static_ncl_mse_against_correlation_penalty_coefficient_0_1_fixed_architecture_2_2_2}
    \end{minipage}
    \hfill
    \begin{minipage}[t]{0.48\textwidth}
        \vspace{0pt}
        \centering
        \includegraphics[width=\linewidth]{img/static-ncl/static_ncl_mse_against_correlation_penalty_coefficient_0_1_fixed_architecture_2_2_2_2_2_2.png}
        \caption{Impact of correlation penalty coefficient (in the range $[0, 1]$) on MSE. Base learners had 6 hidden layers of 2 hidden nodes and were trained with 400 epochs. They increased their model complexities by increasing number of hidden layers (each layer always has 2 hidden nodes). }
        \label{fig:static_ncl_mse_against_correlation_penalty_coefficient_0_1_fixed_architecture_2_2_2_2_2_2}
    \end{minipage}
\end{figure}


% 2 Image
\begin{figure}[H]
    \centering

    \begin{minipage}[t]{0.48\textwidth}
        \vspace{0pt}
        \centering
        \includegraphics[width=\linewidth]{img/static-ncl/static_ncl_diversity_coefficient_against_correlation_penalty_coefficient_0_1_fixed_architecture_2.png}
        \caption{Impact of correlation penalty coefficient (in the range $[0, 1]$) on diversity coefficient. Base learners had 1 hidden layer of 2 hidden nodes and were trained with 400 epochs. They increased their model complexities by increasing number of hidden nodes in the hidden layer. }
        \label{fig:static_ncl_diversity_coefficient_against_correlation_penalty_coefficient_0_1_fixed_architecture_2}
    \end{minipage}
    \hfill
    \begin{minipage}[t]{0.48\textwidth}
        \vspace{0pt}
        \centering
        \includegraphics[width=\linewidth]{img/static-ncl/static_ncl_diversity_coefficient_against_correlation_penalty_coefficient_08_15_fixed_architecture_2.png}
        \caption{Impact of correlation penalty coefficient (in the range $[0.8, 1.5]$) on diversity coefficient. Base learners had 1 hidden layer of 2 hidden nodes and were trained with 400 epochs. They increased their model complexities by increasing number of hidden nodes in the hidden layer. }
        \label{fig:static_ncl_diversity_coefficient_against_correlation_penalty_coefficient_08_15_fixed_architecture_2}
    \end{minipage}
\end{figure}


% 2 Image
\begin{figure}[H]
    \centering

    \begin{minipage}[t]{0.48\textwidth}
        \vspace{0pt}
        \centering
        \includegraphics[width=\linewidth]{img/static-ncl/static_ncl_diversity_coefficient_against_correlation_penalty_coefficient_0_1_fixed_architecture_6.png}
        \caption{Impact of correlation penalty coefficient (in the range $[0, 1]$) on diversity coefficient. Base learners had 1 hidden layer of 6 hidden nodes and were trained with 400 epochs. They increased their model complexities by increasing number of hidden nodes in the hidden layer. }
        \label{fig:static_ncl_diversity_coefficient_against_correlation_penalty_coefficient_0_1_fixed_architecture_6}
    \end{minipage}
    \hfill
    \begin{minipage}[t]{0.48\textwidth}
        \vspace{0pt}
        \centering
        \includegraphics[width=\linewidth]{img/static-ncl/static_ncl_diversity_coefficient_against_correlation_penalty_coefficient_08_15_fixed_architecture_6.png}
        \caption{Impact of correlation penalty coefficient (in the range $[0.8, 1.5]$) on diversity coefficient. Base learners had 1 hidden layer of 6 hidden nodes and were trained with 400 epochs. They increased their model complexities by increasing number of hidden nodes in the hidden layer. }
        \label{fig:static_ncl_diversity_coefficient_against_correlation_penalty_coefficient_08_15_fixed_architecture_6}
    \end{minipage}
\end{figure}


% 2 Image
\begin{figure}[H]
    \centering

    \begin{minipage}[t]{0.48\textwidth}
        \vspace{0pt}
        \centering
        \includegraphics[width=\linewidth]{img/static-ncl/static_ncl_diversity_coefficient_against_correlation_penalty_coefficient_0_1_fixed_architecture_12.png}
        \caption{Impact of correlation penalty coefficient (in the range $[0, 1]$) on diversity coefficient. Base learners had 1 hidden layer of 12 hidden nodes and were trained with 400 epochs. They increased their model complexities by increasing number of hidden nodes in the hidden layer. }
        \label{fig:static_ncl_diversity_coefficient_against_correlation_penalty_coefficient_0_1_fixed_architecture_12}
    \end{minipage}
    \hfill
    \begin{minipage}[t]{0.48\textwidth}
        \vspace{0pt}
        \centering
        \includegraphics[width=\linewidth]{img/static-ncl/static_ncl_diversity_coefficient_against_correlation_penalty_coefficient_08_15_fixed_architecture_12.png}
        \caption{Impact of correlation penalty coefficient (in the range $[0.8, 1.5]$) on diversity coefficient. Base learners had 1 hidden layer of 12 hidden nodes and were trained with 400 epochs. They increased their model complexities by increasing number of hidden nodes in the hidden layer. }
        \label{fig:static_ncl_diversity_coefficient_against_correlation_penalty_coefficient_08_15_fixed_architecture_12}
    \end{minipage}
\end{figure}


% 2 Image
\begin{figure}[H]
    \centering

    \begin{minipage}[t]{0.48\textwidth}
        \vspace{0pt}
        \centering
        \includegraphics[width=\linewidth]{img/static-ncl/static_ncl_diversity_coefficient_against_correlation_penalty_coefficient_0_1_fixed_architecture_2_2_2.png}
        \caption{Impact of correlation penalty coefficient (in the range $[0, 1]$) on diversity coefficient. Base learners had 3 hidden layers of 2 hidden nodes and were trained with 400 epochs. They increased their model complexities by increasing number of hidden layers (each layer always has 2 hidden nodes). }
        \label{fig:static_ncl_diversity_coefficient_against_correlation_penalty_coefficient_0_1_fixed_architecture_2_2_2}
    \end{minipage}
    \hfill
    \begin{minipage}[t]{0.48\textwidth}
        \vspace{0pt}
        \centering
        \includegraphics[width=\linewidth]{img/static-ncl/static_ncl_diversity_coefficient_against_correlation_penalty_coefficient_0_1_fixed_architecture_2_2_2_2_2_2.png}
        \caption{Impact of correlation penalty coefficient (in the range $[0, 1]$) on diversity coefficient. Base learners had 6 hidden layers of 2 hidden nodes and were trained with 400 epochs. They increased their model complexities by increasing number of hidden layers (each layer always has 2 hidden nodes). }
        \label{fig:static_ncl_diversity_coefficient_against_correlation_penalty_coefficient_0_1_fixed_architecture_2_2_2_2_2_2}
    \end{minipage}
\end{figure}


% 2 Image
\begin{figure}[H]
    \centering

    \begin{minipage}[t]{0.48\textwidth}
        \vspace{0pt}
        \centering
        \includegraphics[width=\linewidth]{img/static-ncl/static_ncl_mse_against_correlation_penalty_coefficient_0_1_fixed_ensemble_size_6_architecture_6_comparator_traditional.png}
        \caption{Impact of correlation penalty coefficient (in the range $[0, 1]$) on diversity coefficient. Each ensemble model has 6 base learners. Base learners had 1 hidden layer of 6 hidden nodes and were trained with 400 epochs. They increased their model complexities by increasing number of hidden nodes in the hidden layer. }
        \label{fig:static_ncl_mse_against_correlation_penalty_coefficient_0_1_fixed_ensemble_size_6_architecture_6_comparator_traditional}
    \end{minipage}
    \hfill
    \begin{minipage}[t]{0.48\textwidth}
        \vspace{0pt}
        \centering
        \includegraphics[width=\linewidth]{img/static-ncl/static_ncl_mse_against_correlation_penalty_coefficient_08_15_fixed_ensemble_size_6_architecture_6_comparator_traditional.png}
        \caption{Impact of correlation penalty coefficient (in the range $[0.8, 1.5]$) on diversity coefficient. Each ensemble model has 6 base learners. Base learners had 1 hidden layer of 6 hidden nodes and were trained with 400 epochs. They increased their model complexities by increasing number of hidden nodes in the hidden layer. }
        \label{fig:static_ncl_mse_against_correlation_penalty_coefficient_08_15_fixed_ensemble_size_6_architecture_6_comparator_traditional}
    \end{minipage}
\end{figure}


% 2 Image
\begin{figure}[H]
    \centering

    \begin{minipage}[t]{0.48\textwidth}
        \vspace{0pt}
        \centering
        \includegraphics[width=\linewidth]{img/static-ncl/static_ncl_mse_against_correlation_penalty_coefficient_0_1_fixed_ensemble_size_6_architecture_2_2_2_comparator_traditional.png}
        \caption{Impact of correlation penalty coefficient (in the range $[0, 1]$) on MSE. Each ensemble model always has 6 base learners. Base learners had 6 hidden layers of 2 hidden nodes and were trained with 400 epochs. They increased their model complexities by increasing number of hidden layers (each layer always has 2 hidden nodes). }
        \label{fig:static_ncl_mse_against_correlation_penalty_coefficient_0_1_fixed_ensemble_size_6_architecture_2_2_2_comparator_traditional}
    \end{minipage}
    \hfill
    \begin{minipage}[t]{0.48\textwidth}
        \vspace{0pt}
        \centering
        \includegraphics[width=\linewidth]{img/static-ncl/static_ncl_mse_against_correlation_penalty_coefficient_08_15_fixed_ensemble_size_6_architecture_2_2_2_comparator_traditional.png}
        \caption{Impact of correlation penalty coefficient (in the range $[0.8, 1.5]$) on MSE. Each ensemble model always has 6 base learners. Base learners had 6 hidden layers of 2 hidden nodes and were trained with 400 epochs. They increased their model complexities by increasing number of hidden layers (each layer always has 2 hidden nodes). }
        \label{fig:static_ncl_mse_against_correlation_penalty_coefficient_08_15_fixed_ensemble_size_6_architecture_2_2_2_comparator_traditional}
    \end{minipage}
\end{figure}


% 2 Image
\begin{figure}[H]
    \centering

    \begin{minipage}[t]{0.48\textwidth}
        \vspace{0pt}
        \centering
        \includegraphics[width=\linewidth]{img/static-ncl/static_ncl_mse_against_correlation_penalty_coefficient_0_1_fixed_ensemble_size_6_architecture_6_comparator_mlp.png}
        \caption{Impact of correlation penalty coefficient (in the range $[0, 1]$) on diversity coefficient. Each ensemble model has 6 base learners. Base learners had 1 hidden layer of 6 hidden nodes and were trained with 400 epochs. They increased their model complexities by increasing number of hidden nodes in the hidden layer. }
        \label{fig:static_ncl_mse_against_correlation_penalty_coefficient_0_1_fixed_ensemble_size_6_architecture_6_comparator_mlp}
    \end{minipage}
    \hfill
    \begin{minipage}[t]{0.48\textwidth}
        \vspace{0pt}
        \centering
        \includegraphics[width=\linewidth]{img/static-ncl/static_ncl_mse_against_correlation_penalty_coefficient_08_15_fixed_ensemble_size_6_architecture_6_comparator_mlp.png}
        \caption{Impact of correlation penalty coefficient (in the range $[0.8, 1.5]$) on diversity coefficient. Each ensemble model has 6 base learners. Base learners had 1 hidden layer of 6 hidden nodes and were trained with 400 epochs. They increased their model complexities by increasing number of hidden nodes in the hidden layer. }
        \label{fig:static_ncl_mse_against_correlation_penalty_coefficient_08_15_fixed_ensemble_size_6_architecture_6_comparator_mlp}
    \end{minipage}
\end{figure}


% 2 Image
\begin{figure}[H]
    \centering

    \begin{minipage}[t]{0.48\textwidth}
        \vspace{0pt}
        \centering
        \includegraphics[width=\linewidth]{img/static-ncl/static_ncl_mse_against_correlation_penalty_coefficient_0_1_fixed_ensemble_size_6_architecture_2_2_2_comparator_mlp.png}
        \caption{Impact of correlation penalty coefficient (in the range $[0, 1]$) on MSE. Each ensemble model always has 6 base learners. Base learners had 6 hidden layers of 2 hidden nodes and were trained with 400 epochs. They increased their model complexities by increasing number of hidden layers (each layer always has 2 hidden nodes). }
        \label{fig:static_ncl_mse_against_correlation_penalty_coefficient_0_1_fixed_ensemble_size_6_architecture_2_2_2_comparator_mlp}
    \end{minipage}
    \hfill
    \begin{minipage}[t]{0.48\textwidth}
        \vspace{0pt}
        \centering
        \includegraphics[width=\linewidth]{img/static-ncl/static_ncl_mse_against_correlation_penalty_coefficient_08_15_fixed_ensemble_size_6_architecture_2_2_2_comparator_mlp.png}
        \caption{Impact of correlation penalty coefficient (in the range $[0.8, 1.5]$) on MSE. Each ensemble model always has 6 base learners. Base learners had 6 hidden layers of 2 hidden nodes and were trained with 400 epochs. They increased their model complexities by increasing number of hidden layers (each layer always has 2 hidden nodes). }
        \label{fig:static_ncl_mse_against_correlation_penalty_coefficient_08_15_fixed_ensemble_size_6_architecture_2_2_2_comparator_mlp}
    \end{minipage}
\end{figure}


% 2 Image
\begin{figure}[H]
    \centering

    \begin{minipage}[t]{0.48\textwidth}
        \vspace{0pt}
        \centering
        \includegraphics[width=\linewidth]{img/static-ncl/static_ncl_diversity_coefficient_against_correlation_penalty_coefficient_0_1_fixed_ensemble_size_6_architecture_6.png}
        \caption{Impact of correlation penalty coefficient (in the range $[0, 1]$) on diversity coefficient. Each ensemble model has 6 base learners. Base learners had 1 hidden layer of 6 hidden nodes and were trained with 400 epochs. They increased their model complexities by increasing number of hidden nodes in the hidden layer. }
        \label{fig:static_ncl_diversity_coefficient_against_correlation_penalty_coefficient_0_1_fixed_ensemble_size_6_architecture_6}
    \end{minipage}
    \hfill
    \begin{minipage}[t]{0.48\textwidth}
        \vspace{0pt}
        \centering
        \includegraphics[width=\linewidth]{img/static-ncl/static_ncl_diversity_coefficient_against_correlation_penalty_coefficient_08_15_fixed_ensemble_size_6_architecture_6.png}
        \caption{Impact of correlation penalty coefficient (in the range $[0.8, 1.5]$) on diversity coefficient. Each ensemble model has 6 base learners. Base learners had 1 hidden layer of 6 hidden nodes and were trained with 400 epochs. They increased their model complexities by increasing number of hidden nodes in the hidden layer. }
        \label{fig:static_ncl_diversity_coefficient_against_correlation_penalty_coefficient_08_15_fixed_ensemble_size_6_architecture_6}
    \end{minipage}
\end{figure}


% 2 Image
\begin{figure}[H]
    \centering

    \begin{minipage}[t]{0.48\textwidth}
        \vspace{0pt}
        \centering
        \includegraphics[width=\linewidth]{img/static-ncl/static_ncl_diversity_coefficient_against_correlation_penalty_coefficient_0_1_fixed_ensemble_size_6_architecture_2_2_2.png}
        \caption{Impact of correlation penalty coefficient (in the range $[0, 1]$) on diversity coefficient. Each ensemble model always has 6 base learners. Base learners had 6 hidden layers of 2 hidden nodes and were trained with 400 epochs. They increased their model complexities by increasing number of hidden layers (each layer always has 2 hidden nodes). }
        \label{fig:static_ncl_diversity_coefficient_against_correlation_penalty_coefficient_0_1_fixed_ensemble_size_6_architecture_2_2_2}
    \end{minipage}
    \hfill
    \begin{minipage}[t]{0.48\textwidth}
        \vspace{0pt}
        \centering
        \includegraphics[width=\linewidth]{img/static-ncl/static_ncl_diversity_coefficient_against_correlation_penalty_coefficient_08_15_fixed_ensemble_size_6_architecture_2_2_2.png}
        \caption{Impact of correlation penalty coefficient (in the range $[0.8, 1.5]$) on diversity coefficient. Each ensemble model always has 6 base learners. Base learners had 6 hidden layers of 2 hidden nodes and were trained with 400 epochs. They increased their model complexities by increasing number of hidden layers (each layer always has 2 hidden nodes). }
        \label{fig:static_ncl_diversity_coefficient_against_correlation_penalty_coefficient_08_15_fixed_ensemble_size_6_architecture_2_2_2}
    \end{minipage}
\end{figure}


\section{Stage 4: Combining NCL with NEAT Adaptively}
\label{sec:stage_4}


\subsection{Motivations}
In Section~\ref{sec:stage_3}, base learners were explicitly forced to be diverse using NCL. The assumption that all base learners have the same architecture, i.e. MLPs with the same number of hidden nodes and topologies, had been made. Base learners were different because of their weights only. What if the base learners have diverse architectures? How will this affect the ensemble performance? What is the best architecture for this problem? Furthermore, as shown in Section~\ref{sec:stage_3}, different ensemble models of different base learners had different $\lambda^*$. Which $\lambda$ value should be used without trying all possible values? In this stage, NEAT was combined with NCL. There were a few aims. The first goal was to create ensemble models composed of base learners with diverse architectures and weights. Furthermore, one should be able to determine the best base learner architecture through evolution. Lastly, one should be able to determine $\lambda$ dynamically instead of fixing it before training. 


\subsection{Design and Implementation}


\subsubsection{Reference experiment}
The experiment on NEAT-NCL ensemble model was based on the NEAT-NCL experiments conducted in Sheng et al.\ \cite{sheng2017niching}. It aimed at reproducing similar experiment results to study the effect of NEAT with NCL ensemble loss. To extend the experiment, 3 modifications had been added to the NEAT-NCL experiment in this work. In order to better compare the performance of NEAT-NCL ensemble model and the previous ensemble performance, MSE of a control MLP, a traditional ensemble model (from Section~\ref{sec:stage_2}) and a NCL ensemble model (from Section~\ref{sec:stage_3}) had been added. Furthermore, the original experiment contained graphs of MSE, number of hidden nodes and population diversity $D$. To allow deeper analysis, more graphs had been added to this experiment, including raw fitness, adjusted fitness, sharing factor, $\sigma$, $\lambda$. Lastly, the formula of genetic distance for calculating the population diversity was modified. The modified formula has different weights for different types of genes and does not treat all connection weights equally. 


\subsubsection{NEAT with NCL algorithm}
The traditional NEAT algorithm (mentioned in Section~\ref{subsec:background_neat}) was modified. Between \textit{reproduction} and \textit{replacement}, a step of \textit{backpropagation with NCL} was added. Then, the fitness of each base learner was updated before \textit{replacement}. 

Below contains the steps of the algorithm. 

\begin{enumerate}
    \item \textbf{Initialisation} \\
    An initial population of $m$ neural network base learners is generated. Each base learner is encoded using graph encoding represented by a collection of nodes and edges. The number of hidden neurons, connections and connection weights in the base learners are initialised randomly. Based on the configurations, the initial number of connections and their weights can be partially controlled (more explanations in Section~\ref{subsubsec:stage_3_model_construction}). 

    \item \textbf{Evaluation} \\
    The raw fitness of each base learner in the population is computed using Equation~\ref{equation:neat_ncl_raw_fitness}. 
    To maintain population diversity and prevent premature convergence, fitness sharing is implemented. In traditional fitness sharing, the niche radius is pre-computed. In most real applications, there is no prior knowledge about fitness landscape available and it is difficult to set the niche radius appropriately. In this experiment, the niche radius $\sigma$ is computed dynamically using Equation~\ref{equation:neat_ncl_niche_radius}. To compute the genetic distance between 2 base learners, Equation~\ref{equation:neat_ncl_genetic_distance} is used. Below contains the pseudo-code used to update the raw fitness of 1 base learner by taking into account sharing factor during fitness sharing: 

    \begin{lstlisting}[language=Python, caption={Modified fitness sharing}]
    # The sharing factor of current_base_learner
    sharing_factor = 0
    # Calculate the dynamic niche radius based on the current ensemble size m
    sigma = get_dynamic_niche_radius(m)
    # A scaling factor pre-computed by developers
    alpha = 1

    for other_base_learner in population where other_base_learner is not current_base_learner:
        d = get_genetic_distance(current_base_learner, other_base_learner)

        if d < sigma:
            # current_base_learner is genetically too close to other_base_learner
            sharing_factor += 1 - (d / sigma) ** alpha
        
    # sharing_factor is ready

    # Update the raw fitness of the current base learner for fitness sharing
    current_base_learner.fitness = current_base_learner.fitness / sharing_factor
    \end{lstlisting}
    The fitness of every base learner is updated using the above algorithm. The updated fitness is called \textit{adjusted fitness}. 

    \item \textbf{Selection} \\
    Based on the adjusted fitness, stronger parent genomes are selected. Crossover among dissimilar individuals often produces low-performance offsprings. Thus, the selection process is modified. Specifically, during selection, parent genomes are ranked in descending order of their adjusted fitness values. $P_1$ and $P_2$, which have the 2 highest fitness values, will be mated. Then, $P_3$ and $P_4$ are mated and so on. When there are not enough offspring genomes produced after reaching the bottom of the population, the pointers go back and $P_1$ and $P_2$ are mated again and so on. 

    \item \textbf{Reproduction} \\
    Crossover and mutation on selected paired parent genomes are applied to generate stronger intermediate offsprings. 

    \item \textbf{Backpropagation with NCL} \\
    Backpropagation with NCL is applied to intermediate offsprings to update their connection weights. Each offspring has the loss function specified in Appendix~\ref{appendix:loss_function_ncl}. An appropriate value of $\lambda$ is required for such loss function, which is dynamically computed using Equation~\ref{equation:neat_ncl_dynamic_lambda}. Equation~\ref{equation:neat_ncl_dynamic_lambda} requires population diversity $D$, which is dynamically computed (refer to Section~\ref{subsubsec:neat_ncl_population_diversity} for the calculation steps). 
    Based on Equation~\ref{equation:neat_ncl_dynamic_lambda}, $\lambda$ is set to a very large number initially (because the initial population has little diversity, i.e. small $D$) and it is decreased gradually until it reaches its optimal value. In other words, diversity of base learners is prioritised at the beginning but accuracy receives more attention as there are more diverse genomes in the population. 

    \item \textbf{Evaluation of offsprings} \\
    The fitness of each offspring base learner is computed. The same procedure from Step 2 are applied to the new genomes created. This means the current population contains both offspring and parent genomes. 

    \item \textbf{Replacement} \\
    To maintain the size of the population to $m$, all genomes are ranked based on their values of adjusted fitness. The top $m$ genomes are selected to survive in the next generation while the remaining genomes are eliminated. Since elitism is enabled, certain ``elite'' genomes of each species are exceptionally copied directly to the next generation. 

    \item \textbf{Termination} \\
    In the current configurations, steps 2-7 are repeated until the maximum number of generations is reached. The fitness of base learners is not used to determine termination and is merely for selection. 
\end{enumerate}

A \textbf{subpopulation} is defined as a group of genomes with the same number of hidden nodes. 


\subsubsection{Population diversity}
\label{subsubsec:neat_ncl_population_diversity}
The population diversity of a population is computed using the steps below: 

\begin{enumerate}
    \item \textbf{Output distance} \\
    For each 2 genomes within a subpopulation, calculate their output distance $output\_distance_{ij}$ using Equation~\ref{equation:neat_ncl_output_distance}. 

    \item \textbf{Genetic distance} \\
    For each 2 genomes within a subpopulation, calculate their genetic distance $genetic\_distance_{ij}$ using Equation~\ref{equation:neat_ncl_genetic_distance}. 

    \item \textbf{Diversity contribution} \\
    For each 2 genomes within a subpopulation, calculate their contribution to the subpopulation diversity using Equation~\ref{equation:neat_ncl_diversity_contribution}. 

    \item \textbf{Subpopulation diversity} \\
    For each subpopulation of the whole population, calculate the subpopulation diversity using Equation~\ref{equation:neat_ncl_subpopulation_diversity}. 

    \item \textbf{Population diversity} \\
    Calculate the population diversity using Equation~\ref{equation:neat_ncl_population_diversity}. 
\end{enumerate}


\subsubsection{Output distance}
The output distance between 2 neural networks is computed using the equation below: 

\begin{equation}
    \label{equation:neat_ncl_output_distance}
    output\_distance_{ij} = \frac{1}{N} \sum_{n = 1}^N \left| \hat{y}_{i,n} - \hat{y}_{j,n} \right|
\end{equation}

where $N$ is the number of training data. $\hat{y}_{i,n}$ and $\hat{y}_{j,n}$ are the predictions of the $i$th and $j$th networks, on the $n$th training data respectively. 


\subsubsection{Genetic distance}
The genetic distance (i.e. compatibility distance) between 2 neural networks is computed using the equation below: 

\begin{equation}
    \label{equation:neat_ncl_genetic_distance}
    \delta = \frac{c_1 N_{excess}}{N} + \frac{c_2 N_{disjoint}}{N} + c_3 \bar{W}
\end{equation}

where $\delta$ is the genetic distance, $N_{excess}$ is the number of excess genes, $N_{disjoint}$ is the number of disjoint genes, $N$ is the number of genes in the larger genome (acting as a normalisation factor), $\bar{W}$ is the average weight difference of matching genes, $c_1$, $c_2$ and $c_3$ are coefficients controlling the importance of each term. 


\subsubsection{Diversity contribution}
The diversity contribution between 2 neural networks within a subpopulation is computed using the equation below: 
    
\begin{equation}
    \label{equation:neat_ncl_diversity_contribution}
    subpopulation\_diversity_{ij} = \frac{fitness_i + fitness_{j}}{2 \times fitness_{average}} \times genetic\_distance_{ij} \times output\_distance_{ij}
\end{equation}

where $f_{average}$ is the average fitness of the subpopulation. $fitness_i$ and $fitness_j$ are the fitness of the $i$th and $j$th networks respectively. 


\subsubsection{Subpopulation diversity}
The subpopulation diversity is computed using the equation below: 

\begin{equation}
    \label{equation:neat_ncl_subpopulation_diversity}
    subpopulation\_diversity = \frac{2}{k(k - 1)} \sum_{i = 1}^{k - 1} \sum_{j = i + 1}^k subpopulation\_diversity_{ij}
\end{equation}

where $k$ is the number of neural networks in the subpopulation. The sum of $k$ of all subpopulations gives $m$. Note that this formula is calculating the average subpopulation diversity. 


\subsubsection{Population diversity}
The population diversity is computed using the equation below: 
    
\begin{equation}
    \label{equation:neat_ncl_population_diversity}
    population\_diversity = \sum_{i = 1}^c subpopulation\_diversity_{i}
\end{equation}

where $c$ is the number of subpopulations. Population diversity can be treated as sum of all subpopulation diversities. 


\subsubsection{Dynamic niche radius}
The dynamic niche radius for fitness sharing is computed using the equation below: 

\begin{equation}
    \label{equation:neat_ncl_niche_radius}
    \sigma = \frac{1}{2 m^2} \sum_{i = 1}^{m} \sum_{j = 1}^{m} distance_{ij}
\end{equation}

where $m$ is the ensemble size, i.e. population size. $distance_{ij}$ is the genetic distance between the $i$th and $j$ base learners. 


\subsubsection{Raw fitness}
The raw fitness of each base learner in the population is computed using the equation below: 

\begin{equation}
    \label{equation:neat_ncl_raw_fitness}
    fitness = (\hat{y}_i - y)^2 - (\hat{y}_i - \bar{y})^2 + \sum_{ij} \frac{weight_{ij}^2}{1 + weight_{ij}^2}
\end{equation}

where $\hat{y}_i$ is the prediction made by the $i$th base learner, $\bar{y}$ is the mean prediction of all base learners, $y$ is the true value. The third term is a weight decay, which is used to estimate the complexity of the base learner's architecture and improve the generalisation of base learners. $weight_{ij}$ denotes the connection weight from Neuron $i$ to Neuron $j$. 


\subsubsection{Dynamic Correlation Penalty Coefficient}
The value of $\lambda$ is computed dynamically using the equation below: 

\begin{equation}
    \label{equation:neat_ncl_dynamic_lambda}
    \lambda = \lambda_{max} - \frac{D}{D_{max}} \times (\lambda_{max} - \lambda_{min})
\end{equation}

where $\lambda_{max}$ and $\lambda_{min}$ correspond to the upper and lower bounds values of $\lambda$ set by developers. Thus, $\lambda$ varies from $\lambda_{min}$ to $\lambda_{max}$. $D$ is the population diversity computed in the current population. $D_{max}$ is the maximum population diversity recorded in the population. 


\subsubsection{Model construction}
\label{subsubsec:stage_3_model_construction}
The NEAT-NCL codebase was constructed upon the \textit{NEAT-Python} library, extending its core evolutionary mechanisms, e.g. selection and crossover, to support NCL. Key modifications included customised selection, fitness sharing and the integration of backpropagation with NCL post-reproduction. 

\paragraph{Data Management and Algorithmic Support}
To handle the high computational overhead of iterating through genome combinations, a specialised \texttt{Table} data structure in \path{data_container.py} was implemented. This structure facilitates efficient storage and retrieval of population matrices, including: 

\begin{enumerate}
    \item \textbf{Genetic and Output Distances} \\
    Measuring structural and functional divergence. 

    \item \textbf{Subpopulation Diversity} \\
    Tracking the contribution of individual genomes to niche variety. 

    \item \textbf{Fitness Sharing} \\
    Scaling individual fitness based on the proximity of peers. 
\end{enumerate}

Algorithmic logic and statistical formulae, e.g. the sharing factor and $\sigma$ calculations, were encapsulated in \path{neat_ncl_trainer.py} and \path{maths_helper.py}. The integrity of these calculations was verified through extensive unit testing in \path{data_container_test.py} and \path{graph_helper_test.py}. 

\paragraph{Evolutionary Configuration}
The execution is governed by \texttt{config-housing-neat-ncl.ini}, which defines the structural and behavioural attributes of the population: 

\begin{enumerate}
    \item \textbf{Weight Initialisation} \\
    Connections use a Gaussian distribution (\lstinline|weight_init_mean = 0.0|, \lstinline|weight_init_stdev = 1.0|). 

    \item \textbf{Topology Evolution} \\
    Genomes can be initialised as \lstinline|unconnected| (no initial connections), \lstinline|full_direct| (complete input-to-output mapping) or \lstinline|partial_direct 0.5| (a $50\%$ random subset of a full-direct mapping). Complexity is added incrementally through connection mutation. 

    \item \textbf{Population Dynamics} \\
    An ensemble size of $150$ can be maintained using \lstinline|pop_size = 150| with \lstinline|elitism = 2| to preserve the best-performing genomes across generations. Species stability is ensured via \lstinline|min_species_size = 2|. 

    \item \textbf{Network Constraints} \\
    The \lstinline|feed_forward = True| setting is mandatory to support the topological evaluation algorithm described in Section~\ref{subsubsec:dynamic_nn_on_gpu}. 
\end{enumerate}

\paragraph{Structural Representation}
The transition from genotype to phenotype was managed by \texttt{NEATNCLBaseLearner} (\path{neat_ncl_base_learner.py}), which instantiated individual phenomes. Below contains some of its important components: 

\begin{enumerate}
    \item \textbf{Nodes and Edges} \\
    Neurons are represented by \texttt{Node} objects (\path{node.py}), identified by IDs where input neurons are strictly negative ($node\_id_{input} \in \{-1, -2, ...\}$). Weighted connections are managed via \texttt{Edge} objects (\path{edge.py}). 

    \item \textbf{Forward Pass} \\
    The \texttt{torch.nn.Module} prediction method \texttt{forward} has been re-engineered to follow a strict topological evaluation order, ensuring that signal propagation respects the dynamic nature of the evolved network. 
\end{enumerate}

The entire population was orchestrated by \texttt{NEATNCLEnsembleNet} (\path{neat_ncl_ensemble_net.py}). Its primary method, \texttt{evolve}, executed the evolutionary loop for a specified number of generations using a dictionary of parameters, e.g. $\lambda_{min}$, $\lambda_{max}$ and training partitions, to evaluate and refine genome fitness through NCL. 


\subsubsection{Dynamic neural network on GPU}
\label{subsubsec:dynamic_nn_on_gpu}
GPUs are built for massively parallel, uniform workloads. In NEAT, neural networks are evolved and their topologies change over each generation. Each NEAT individual in the population has a different architecture so uniform batch of matrix multiplications and high-throughput SIMD execution cannot be executed on GPU. 
However, instead of running networks on CPU, which is much slower, an algorithm has been developed to convert a dynamic neural network to a GPU-runnable representation. The goal of the algorithm is to create a different GPU-runnable neural network such that it is equivalent to the original dynamic neural network. The new regular neural network should have the correct sequential order of data flow between neurons and the correct connection weights applied to give the same final prediction as the original network. 
Each neural network can be treated as a DAG with nodes representing neurons and edges representing weighted connections. Essentially, the algorithm converts a DAG to a graph of nodes (\textit{fully-connected graph}) such that each node is connected to every node in the next layer, i.e. a fully-connected layer. 
Below contains the steps of the algorithm: 

\begin{enumerate}
    \item \textbf{Topological sort} \\
    The topological sort of nodes is calculated. This creates a correct sequential order of data flow through nodes. 

    \item \textbf{Set up input neurons in the first layer} \\
    For a fully-connected layer, all input neurons must be in the first layer. New nodes are created in the first layer and connections with weight~$1$ are added to allow direct data flow from the first layer to every input neuron in the topological sort. 

    \item \textbf{Create connections} \\
    This step ensures connection weights are applied in the correct order and with the right value according to the original DAG. For every edge \texttt{(start\_node\_id, end\_node\_id)} in the DAG, the connection weight ($w$), IDs of the start and end nodes (\texttt{start\_node\_id} and \texttt{end\_node\_id}) are retrieved. The algorithm attempts to form a linking between the corresponding 2 nodes in the fully-connected graph (\texttt{start\_node} and \texttt{end\_node}). It follows 2 simple rules: 
    
    \begin{enumerate}
        \item \textbf{2 nodes are directly next to each other} \\
        Add a direct connection with weight~$w$ between the 2 nodes

        \item \textbf{2 nodes are not directly next to each other} \\
        For each layer between the 2 nodes, add a node at the bottom of the layer. Add connections connecting \texttt{start\_node} and \texttt{end\_node} through the newly added nodes. All connections have weights $1$ except the last connection, which has weight~$w$. 
    \end{enumerate}

    \item \textbf{Add non-functional connections} \\
    At this state, an equivalent minimum graph of nodes and edges is constructed successfully and data flows with the correct sequential orders and values. Then, a connection with weight~$0$ is added to connect every pair of nodes which does not have an existing connection yet. This ensures the final graph represents a fully-connected network and can be run on GPU using PyTorch. 

    \item \textbf{Create instances of \texttt{nn.LinearLayer} based on neuron bias and connection weight} \\
    Using the generated fully-connected graph, instances of \texttt{nn.LinearLayer} are created. An instance of \texttt{nn.LinearLayer} is represented by a matrix. They are initialised with weight~$0$ and bias~$0$ initially. All non-zero connection weights and biases are then used to update the \texttt{nn.LinearLayer} matrix at the correct positions. 
\end{enumerate}

\paragraph{Assumption}
In order for this algorithm to work, 2 assumptions have been made. 
There must not be recurrent connections, i.e. the network must be \textbf{feed-forward}. 
Furthermore, the activation function used must be \textbf{idempotent} because the algorithm expects the values of input data remain unchanged when passing every node until they reach the last connection with weight~$w$. The activation function \texttt{torch.nn.ReLU} used in this experiment was acceptable as long as the input data are \textit{non-negative}. Since this algorithm was not used in the experiment (because it does not support backpropagation) so the use of \texttt{torch.nn.ReLU} in base learners was accepted. Another acceptable activation function for this algorithm is the \textit{identity function}. 

\paragraph{Limitation}
The algorithm is \textbf{irreversible} as it can only convert a DAG to a fully-connected graph but not the other way around. Thus, after backpropagation with NCL, the weights of the fully-connected graph are updated and the DAG of the genome will no longer be updated and cannot be used for crossover for the next generation. Therefore, in this project, this algorithm is only used when the phenomes are executed in inference mode but not training mode where weights are updated via backpropagation. For training the networks and evolutionary algorithm, networks are run on CPU and this algorithm is not used, e.g. during the experiments. 
ReLU can be used in this algorithm only if every node has a non-negative value. 

\paragraph{Testing}
Unit tests in \path{algorithm_test.py} and \path{mlp_test.py} are created to test the correctness of the algorithm implementation. The model created by the algorithm is compared against that created using 
\\
\texttt{neat.nn.FeedForwardNetwork.create(genome, config)} from NEAT-Python to validate the algorithm correctness. 

\paragraph{Example}
Figure~\ref{fig:neat_ncl_dynamic_nn_on_gpu_example_1} shows a dynamic neural network. The number in each neuron is the node ID and the number next to each connection is its weight. The network has 3 input neurons with IDs $-1$, $-2$ and $-3$ respectively. The output neuron with ID $1$ is the output neuron. The topological sort of nodes is $[-3, -2, -1, 2, 3, 1]$, giving Figure~\ref{fig:neat_ncl_dynamic_nn_on_gpu_example_2}. Then, new neurons are added so that all input neurons are connected to the first layer directly, giving Figure~\ref{fig:neat_ncl_dynamic_nn_on_gpu_example_3}. New connections and nodes are then added sequentially so that each connection in the original DAG can be represented, as shown in Figure~\ref{fig:neat_ncl_dynamic_nn_on_gpu_example_4} to Figure~\ref{fig:neat_ncl_dynamic_nn_on_gpu_example_10} accordingly. The network in Figure~\ref{fig:neat_ncl_dynamic_nn_on_gpu_example_10} is an equivalent minimum graph of the original DAG. Lastly, connections with weight~$0$ are added (represented by light grey lines) to make the network fully-connected, giving the final graph shown in Figure~\ref{fig:neat_ncl_dynamic_nn_on_gpu_example_11}. Based on the generated fully-connected graph, 5 matrices of \texttt{nn.LinearLayer} are initialised with $0$ weights and biases. Specific cells (of corresponding positions) in the matrices are then updated due to the 7 non-zero connection weights and 1 non-zero bias. 


% 1 Image
\begin{figure}[H]
    \centering
    \includegraphics[width=0.64\textwidth]{img/neat-ncl/dynamic-nn-on-gpu/neat_ncl_dynamic_nn_on_gpu_example_1.png}
    \caption{An example NEAT dynamic neural network with connection weights and node bias labelled}
    \label{fig:neat_ncl_dynamic_nn_on_gpu_example_1}
\end{figure}


% 2 Image
\begin{figure}[H]
    \centering

    \begin{minipage}[t]{0.48\textwidth}
        \vspace{0pt}
        \centering
        \includegraphics[width=\linewidth]{img/neat-ncl/dynamic-nn-on-gpu/neat_ncl_dynamic_nn_on_gpu_example_2.png}
        \caption{Compute the toplogical sort of nodes (read from left to right). }
        \label{fig:neat_ncl_dynamic_nn_on_gpu_example_2}
    \end{minipage}
    \hfill
    \begin{minipage}[t]{0.48\textwidth}
        \vspace{0pt}
        \centering
        \includegraphics[width=\linewidth]{img/neat-ncl/dynamic-nn-on-gpu/neat_ncl_dynamic_nn_on_gpu_example_3.png}
        \caption{Connect all input neurons to the first layer directly. }
        \label{fig:neat_ncl_dynamic_nn_on_gpu_example_3}
    \end{minipage}
\end{figure}


% 2 Image
\begin{figure}[H]
    \centering

    \begin{minipage}[t]{0.48\textwidth}
        \vspace{0pt}
        \centering
        \includegraphics[width=\linewidth]{img/neat-ncl/dynamic-nn-on-gpu/neat_ncl_dynamic_nn_on_gpu_example_4.png}
        \caption{New connections and nodes are added to represent the connection \texttt{(-1, 2)} in the original DAG. }
        \label{fig:neat_ncl_dynamic_nn_on_gpu_example_4}
    \end{minipage}
    \hfill
    \begin{minipage}[t]{0.48\textwidth}
        \vspace{0pt}
        \centering
        \includegraphics[width=\linewidth]{img/neat-ncl/dynamic-nn-on-gpu/neat_ncl_dynamic_nn_on_gpu_example_5.png}
        \caption{New connections and nodes are added to represent the connection \texttt{(2, 1)} in the original DAG. }
        \label{fig:neat_ncl_dynamic_nn_on_gpu_example_5}
    \end{minipage}
\end{figure}


% 2 Image
\begin{figure}[H]
    \centering

    \begin{minipage}[t]{0.48\textwidth}
        \vspace{0pt}
        \centering
        \includegraphics[width=\linewidth]{img/neat-ncl/dynamic-nn-on-gpu/neat_ncl_dynamic_nn_on_gpu_example_6.png}
        \caption{New connections and nodes are added to represent the connection \texttt{(-2, 2)} in the original DAG. }
        \label{fig:neat_ncl_dynamic_nn_on_gpu_example_6}
    \end{minipage}
    \hfill
    \begin{minipage}[t]{0.48\textwidth}
        \vspace{0pt}
        \centering
        \includegraphics[width=\linewidth]{img/neat-ncl/dynamic-nn-on-gpu/neat_ncl_dynamic_nn_on_gpu_example_7.png}
        \caption{New connections and nodes are added to represent the connection \texttt{(2, 3)} in the original DAG. }
        \label{fig:neat_ncl_dynamic_nn_on_gpu_example_7}
    \end{minipage}
\end{figure}


% 2 Image
\begin{figure}[H]
    \centering

    \begin{minipage}[t]{0.48\textwidth}
        \vspace{0pt}
        \centering
        \includegraphics[width=\linewidth]{img/neat-ncl/dynamic-nn-on-gpu/neat_ncl_dynamic_nn_on_gpu_example_8.png}
        \caption{New connections and nodes are added to represent the connection \texttt{(3, 1)} in the original DAG. }
        \label{fig:neat_ncl_dynamic_nn_on_gpu_example_8}
    \end{minipage}
    \hfill
    \begin{minipage}[t]{0.48\textwidth}
        \vspace{0pt}
        \centering
        \includegraphics[width=\linewidth]{img/neat-ncl/dynamic-nn-on-gpu/neat_ncl_dynamic_nn_on_gpu_example_9.png}
        \caption{New connections and nodes are added to represent the connection \texttt{(-3, 2)} in the original DAG. }
        \label{fig:neat_ncl_dynamic_nn_on_gpu_example_9}
    \end{minipage}
\end{figure}


% 2 Image
\begin{figure}[H]
    \centering

    \begin{minipage}[t]{0.48\textwidth}
        \vspace{0pt}
        \centering
        \includegraphics[width=\linewidth]{img/neat-ncl/dynamic-nn-on-gpu/neat_ncl_dynamic_nn_on_gpu_example_10.png}
        \caption{New connections and nodes are added to represent the connection \texttt{(-3, 1)} in the original DAG. The current graph is already an equivalent minimum graph of the original DAG but it is not enough to create the corresponding PyTorch MLP yet. }
        \label{fig:neat_ncl_dynamic_nn_on_gpu_example_10}
    \end{minipage}
    \hfill
    \begin{minipage}[t]{0.48\textwidth}
        \vspace{0pt}
        \centering
        \includegraphics[width=\linewidth]{img/neat-ncl/dynamic-nn-on-gpu/neat_ncl_dynamic_nn_on_gpu_example_11.png}
        \caption{Add connections with weight~$0$ to make the network fully-connected. The translucent connections have weight~$0$ while the solid connections have weights other than $0$. }
        \label{fig:neat_ncl_dynamic_nn_on_gpu_example_11}
    \end{minipage}
\end{figure}


\subsubsection{Experiment}
\paragraph{Varied hyperparameters}
Table~\ref{table:experiment_4_varied_hyperparameters} contains all varied hyperparameters in this experiment: 

\begin{table}[H]
    \centering
    \begin{tabular}{|c|c|c|c|}
        \hline
        Hyperparameter & Type & Choices / Ranges & Description \\
        \hline
        \small\texttt{ensemble\_sizes} & Integer & $[50, 250]\text{, step}=50$ & $m$ \\
        \hline
    \end{tabular}
    \caption{Varied hyperparameters and their configurations}
    \label{table:experiment_4_varied_hyperparameters}
\end{table}

\paragraph{Fixed hyperparameters}
Table~\ref{table:experiment_4_fixed_hyperparameters} contains all fixed hyperparameters in this experiment: 

\begin{table}[H]
    \centering
    \begin{tabular}{|c|c|c|p{5cm}|}
        \hline
        Hyperparameter & Type & Choices / Ranges & Description \\
        \hline
        \small\texttt{input\_size} & Integer & $13$ & Number of input features of the MLP \\
        \hline
        \small\texttt{output\_size} & Integer & $1$ & Number of labels of the MLP \\
        \hline
        \small\texttt{optimizer} & Categorical & \texttt{torch.optim.Adam} & Optimisation algorithm used for training \\
        \hline
        \small\texttt{learning\_rate} & Float & $0.001$ & Learning rate of the optimisation algorithm \\
        \hline
        \small\texttt{activation} & Categorical & ReLU & Activation function used in each dense layer of the MLP \\
        \hline
        \small\texttt{evolution\_epoch} & Integer & \{$500$\} & Number of evolutionary generations during training \\
        \hline
        \small\texttt{voter} & Categorical & \texttt{arithmetic\_mean} & Type of aggregator used to combine base learners' predictions \\
        \hline
        \small\texttt{min\_lambda} & Float & $0.1$ & $\lambda_{min}$ \\
        \hline
        \small\texttt{max\_lambda} & Float & $0.9$ & $\lambda_{max}$ \\
        \hline
        \small\texttt{evolution\_epoch\_step\_size} & Integer & $10$ & Number of evolutionary generations to collect all metrics once \\
        \hline
        \small\texttt{batch\_size} & Integer & $506$ & Number of training samples passed to phenomes in one go \\
        \hline
    \end{tabular}
    \caption{Fixed hyperparameters and their configurations}
    \label{table:experiment_4_fixed_hyperparameters}
\end{table}

\paragraph{Metric}
Table~\ref{table:experiment_4_metric} contains all evaluation metrics collected in this experiment, which were stored in \path{experiment_runner_neat_ncl_unconnected_arithmetic_mean.csv}, \path{experiment_runner_neat_ncl_partial_direct_arithmetic_mean.csv} and \path{experiment_runner_neat_ncl_full_direct_arithmetic_mean.csv}: 

\begin{table}[H]
    \centering
    \begin{tabular}{|c|c|p{5cm}|}
        \hline
        Metric & Type & Description \\
        \hline
        \small\texttt{loss\_test} & Float & MSE of ensemble model during testing \\
        \hline
        \small\texttt{average\_active\_hidden\_node\_num} & Float & Average number of active hidden nodes in a phenome during testing \\
        \hline
        \small\texttt{diversity\_coefficient\_test} & Float & $\rho$ during testing \\
        \hline
        \small\texttt{population\_diversity\_test} & Float & $D$ during testing \\
        \hline
        \small\texttt{max\_population\_diversity\_test} & Float & $D_{max}$ during testing \\
        \hline
        \small\texttt{correlation\_penalty\_coefficient\_test} & Float & $\lambda$ during testing \\
        \hline
        \small\texttt{subpopulation\_num} & Integer & Number of subpopulations during testing \\
        \hline
        \small\texttt{average\_niche\_radius} & Float & $\sigma$ during testing \\
        \hline
        \small\texttt{average\_sharing\_factor} & Float & Average sharing factor of all phenomes during testing \\
        \hline
        \small\texttt{average\_raw\_fitness} & Float & Average raw fitness of all phenomes during testing \\
        \hline
        \small\texttt{average\_adjusted\_fitness} & Float & Average adjusted fitness (after fitness sharing) of all phenomes during testing \\
        \hline
    \end{tabular}
    \caption{Evaluation metrics}
    \label{table:experiment_4_metric}
\end{table}

\paragraph{Further notes}
For raw fitness and adjusted fitness, the larger their values, the better the genome is. 
Refer to Appendix~\ref{appendix:diversity_coefficient} for the complete formula of $\rho$. 


\subsection{Expected Observations}


\subsubsection{Increase in performance over evolution generations}
Consistent with the principles of biological evolution, NEAT-NCL iteratively refines the population through selective pressure and genetic variation. By selecting high-fitness individuals for reproduction and applying stochastic mutations (e.g. topological structural changes or weight adjustments), the algorithm searches for superior solutions. Consequently, the average individual fitness is expected to increase over successive generations. As these base learners become more specialised and accurate, the collective ensemble MSE should therefore decrease, reflecting the successful adaptation of the population to the underlying data structure. 


\subsubsection{Linear relationship between $\lambda$ and $-D$}
The implementation of a dynamic correlation penalty, $\lambda$, serves as a regulatory mechanism to balance the trade-off between individual accuracy and collective diversity throughout the evolutionary process. By scaling $\lambda$ inversely with the current population diversity ($D$), the optimisation objective adaptively shifts its priority. In the early stages of evolution, when randomly initialised base learners exhibit low functional diversity, a high $\lambda$ value (approaching $\lambda_{\max}$) is applied to force the exploration of the hypothesis space. As NCL effectively drives the population toward a more heterogeneous state, the increase in $D$ triggers a reduction in $\lambda$, allowing the base learners to switch their remaining capacity towards minimising individual MSE. Mathematically, this relationship is defined by the formula below: 

$$
\lambda = \lambda_{\max} - \frac{D}{D_{\max}} \times (\lambda_{\max} - \lambda_{\min})
$$

This formula ensures that the penalty coefficient is strictly bounded between $[\lambda_{\min}, \lambda_{\max}]$. $\lambda$ is a linear function of $-D$. Consequently, the graph of $\lambda$ against $D$ is expected to yield a negatively-sloped straight line. This dynamic scaling prevents the ``over-diversification'' observed in NCL experiments in Seciton~\ref{sec:stage_3} as the pressure to deviate from the ensemble mean is automatically relieved once $D$ is approaching a sufficient population diversity threshold $D_{\max}$. 


\subsubsection{Different values of optimal $\lambda$ for different NEAT-NCL ensemble models}
The dynamic calculation of $\lambda$ functions as an automated tuning mechanism that adapts to the specific constraints of each ensemble configuration. It ensures that the NCL pressure is always relative to the current structural capacity of the ensemble regardless of whether the model is a small committee of complex nodes or a large population of simple ones. According to Section~\ref{subsubsec:neat_ncl_population_diversity}, the calculation of $D$ depends on population-dependent variables like genetic distance and $m$, which vary among different NEAT-NCL ensemble models. Thus, each ensemble model should have different values of $D$ and $D_{max}$. Since $\lambda_{max}$ and $\lambda_{min}$ are fixed, it is expected that different ensemble models have different values of $\lambda^*$ due to their difference in configurations. 


\subsubsection{More diversity seen in NEAT-NCL base learners}
NEAT-NCL creates a more diverse group of learners because it combines the correlation penalty with structural evolution. In a simple average ensemble, models are trained separately and often become redundant because they lack a reason to be different. Standard NCL ensemble models improve on this by using a penalty term to stay unique but they are still restricted by a fixed number of layers and nodes. NEAT-NCL base learners are the most diverse because they vary in both their connection weights and their physical shapes. Since each model can evolve its own unique number of hidden nodes and specific connections, the population explores many different ways to solve the same problem. This mix of different ``brains'' and different ``structures'' makes NEAT-NCL much more varied than models that only change their weights. 


\subsection{Evaluation and Empirical Comparisons}


\subsubsection{How does each attribute change throughout generations of evolution? }
\paragraph{Experiments and Observations}
As base learners evolved, their raw fitness increased because stronger models were selected to reproduce. Figure~\ref{fig:neat_ncl_unconnected_average_raw_fitness_against_generation_index_0_480} shows that fitness improved from $-0.5$ to $-0.35$ within the first 20 generations before reaching a saturation point. Smaller ensembles achieved higher raw fitness but showed more fluctuation, whereas larger ensembles were more stable but had lower individual fitness. This lower fitness in larger groups occurred because NCL forced them to be more diverse, as shown in Figure~\ref{fig:static_ncl_diversity_coefficient_against_correlation_penalty_coefficient_0_1_fixed_architecture_2}. Similar trends appeared in Figure~\ref{fig:neat_ncl_partial_direct_average_raw_fitness_against_generation_index_0_480} and Figure~\ref{fig:neat_ncl_full_direct_average_raw_fitness_against_generation_index_0_480}, regardless of how connections were initialised. 

Adjusted fitness followed a similar but smoother path, as seen in Figure~\ref{fig:neat_ncl_unconnected_average_adjusted_fitness_against_generation_index_0_480}, Figure~\ref{fig:neat_ncl_partial_direct_average_adjusted_fitness_against_generation_index_0_480} and Figure~\ref{fig:neat_ncl_full_direct_average_adjusted_fitness_against_generation_index_0_480}. Since adjusted fitness is the raw fitness divided by a sharing factor ($adjusted\_fitness = \frac{raw\_fitness}{sharing\_factor}$), the variations in raw performance were dampened. 

The niche radius ($\sigma$) also increased initially. Figure~\ref{fig:neat_ncl_unconnected_average_niche_radius_against_generation_index_0_480} shows it rising from $0.13$ to $0.2$ in the first 10 generations. The formula of $\sigma$ is given in Equation~\ref{equation:neat_ncl_niche_radius}. Initially, randomly initialised models were similar but NCL pushed them apart, increasing the average distance between them, i.e. $d_{ij}$, which explained the initial increase in $\sigma$. While ensemble size did not change the radius value itself, larger ensembles produced more stable results because it is harder for a large group to reach an extreme average. This increase in $\sigma$ also caused the sharing factor to rise and saturate after 20 generations, as shown in Figure~\ref{fig:neat_ncl_unconnected_average_sharing_factor_against_generation_index_0_480}. Larger ensembles had higher sharing factors because more models were likely to be within the niche radius of one another. 

While number of evolutions did not directly affect $\lambda$, ensemble size affected the stability and value of $\lambda$. Figure~\ref{fig:neat_ncl_unconnected_correlation_penalty_coefficient_against_generation_index_0_480} showed that smaller ensembles had the most variation in $\lambda$ while large ensembles were the steadiest. This mirrored the population diversity $D$ seen in Figure~\ref{fig:neat_ncl_unconnected_population_diversity_against_generation_index_0_480}. Larger ensembles generally used a smaller $\lambda$. For example, an ensemble model of size 50 used $\lambda \approx 0.8$ while an ensemble model of size 250 used $\lambda \approx 0.5$. This is because larger ensembles had a smaller $D_{max}$, as seen in Figure~\ref{fig:neat_ncl_unconnected_max_population_diversity_against_generation_index_0_480}. Note that some $D_{max}$ in Figure~\ref{fig:neat_ncl_unconnected_max_population_diversity_against_generation_index_0_480} could not be found in Figure~\ref{fig:neat_ncl_unconnected_correlation_penalty_coefficient_against_generation_index_0_480}. This was expected because Figure~\ref{fig:neat_ncl_unconnected_correlation_penalty_coefficient_against_generation_index_0_480} recorded $D$ only every 10 generations, which might have missed the generations at which the current $D_{max}$ occured. 

Finally, ensemble MSE dropped sharply from $0.65$ to $0.4$ in the first 10 generations, as illustrated in Figure~\ref{fig:neat_ncl_unconnected_mse_against_generation_index_0_480_fixed_ensemble_size_50}, Figure~\ref{fig:neat_ncl_unconnected_mse_against_generation_index_0_480_fixed_ensemble_size_100}, Figure~\ref{fig:neat_ncl_unconnected_mse_against_generation_index_0_480_fixed_ensemble_size_150}, Figure~\ref{fig:neat_ncl_unconnected_mse_against_generation_index_0_480_fixed_ensemble_size_200} and Figure~\ref{fig:neat_ncl_unconnected_mse_against_generation_index_0_480_fixed_ensemble_size_250}. Figure~\ref{fig:neat_ncl_unconnected_average_raw_fitness_against_generation_index_0_480} and Figure~\ref{fig:neat_ncl_unconnected_average_adjusted_fitness_against_generation_index_0_480} confirmed this was due to the rapid gain in complexity and thus, fitness during early evolution. However, performance stopped improving after about 30 generations. This same pattern of quick early gains followed by a plateau was also observed in studies by Sheng et al. \cite{sheng2017niching}. In Graph~a and Graph~e of Figure~4 of their experiments, testing errors dropped drastically at the first 20 generations but stopped afterwards. 

\paragraph{Conclusions from Experiment}
As the number of evolution generations increases, ensemble MSE decreases due to increase in raw fitness and adjusted fitness of base learners. Increase in number of evolution generations also leads to increases in niche radius and sharing factor but it does not affect $\lambda$ directly. 


\subsubsection{What is the impact of niche radius? }
\paragraph{Experiments and Observations}
In general, increasing the niche radius did not significantly change the sharing factor. Figure~\ref{fig:neat_ncl_unconnected_average_sharing_factor_against_average_niche_radius} showed that for any given ensemble size, the sharing factor stayed roughly the same even as $\sigma$ increased. For example, with an ensemble size of $50$, the sharing factor consistently stayed around $10$.

However, increasing niche radius $\sigma$ did lead to a slight drop in adjusted fitness, meaning performance got slightly worse. This effect was more noticeable in smaller ensembles. As seen in Figure~\ref{fig:neat_ncl_unconnected_average_adjusted_fitness_against_average_niche_radius}, for a group of $50$, adjusted fitness fell from $-0.3$ to $-0.5$ as $\sigma$ grew from $0.15$ to $0.25$. In contrast, a group of $250$ had a much smaller drop, moving from $-0.35$ to $-0.45$ over a narrower range of $\sigma$ ($0.18$ to $0.22$).

Additionally, increasing niche radius $\sigma$ slightly improved population diversity $D$. This trend was also more prominent in smaller ensembles. Figure~\ref{fig:neat_ncl_unconnected_population_diversity_against_average_niche_radius} indicated that for an ensemble size of $50$, diversity increased from near $0$ to $0.1$ as $\sigma$ moved from $0.05$ to $0.25$. For an ensemble size of $250$, the increase was much smaller, only rising from $0.05$ to $0.08$ when $\sigma$ increased from $0.18$ to $0.22$. 

\paragraph{Conclusions from Experiment}
Increasing niche radius has no direct effect on the sharing factor. However, increasing niche radius leads to a decrease in adjusted fitness and increase in population diversity. 


\subsubsection{What is the impact of ensemble size? }
\paragraph{Experiments and Observations}
A larger ensemble size significantly reduced fluctuations in model attributes. As shown in Figure~\ref{fig:neat_ncl_unconnected_average_raw_fitness_against_generation_index_0_480}, Figure~\ref{fig:neat_ncl_unconnected_average_adjusted_fitness_against_generation_index_0_480}, Figure~\ref{fig:neat_ncl_unconnected_average_niche_radius_against_generation_index_0_480}, Figure~\ref{fig:neat_ncl_unconnected_population_diversity_against_generation_index_0_480} and Figure~\ref{fig:neat_ncl_unconnected_correlation_penalty_coefficient_against_generation_index_0_480}, larger ensembles (size 250) stayed much steadier than smaller ones (size 50). This behaviour followed the \textbf{Law of Large Numbers}, which states that as a sample size grows, the average becomes more consistent and stays closer to the true mean. This is also why the data points for larger ensembles appeared more concentrated in Figure~\ref{fig:neat_ncl_unconnected_average_sharing_factor_against_average_niche_radius}, Figure~\ref{fig:neat_ncl_unconnected_average_adjusted_fitness_against_average_niche_radius} and Figure~\ref{fig:neat_ncl_unconnected_population_diversity_against_average_niche_radius}. 

A larger ensemble size led to a larger initial increase in sharing factor as shown in Figure~\ref{fig:neat_ncl_unconnected_average_sharing_factor_against_generation_index_0_480}. However, the sharing factor fluctuated with the same extent despite ensemble size. On the contrary, a larger ensemble size led to a smaller initial increase in $D_{max}$. Figure~\ref{fig:neat_ncl_unconnected_max_population_diversity_against_generation_index_0_480} showed that $D_{max}$ surged from $0$ to $0.2$ in the beginning generations when ensemble size $m$ was $50$. Figure~\ref{fig:neat_ncl_unconnected_max_population_diversity_against_generation_index_0_480} showed that $D_{max}$ only surged from $0$ to $0.1$ when ensemble size $m$ was $250$. 

While larger ensembles were more stable, they generally had lower individual base learner fitness. Figure~\ref{fig:neat_ncl_unconnected_average_raw_fitness_against_generation_index_0_480} and Figure~\ref{fig:neat_ncl_unconnected_average_adjusted_fitness_against_generation_index_0_480} showed that base learners in a group of $50$ performed better than those in a group of $250$. This is because NCL forced larger groups to be more diverse, which pulled individual models away from the best possible accuracy to ensure the whole group stayed unique. 

As seen in Figure~\ref{fig:neat_ncl_unconnected_max_population_diversity_against_generation_index_0_480}, $D_{max}$ was smaller for larger ensembles, dropping from $0.3$ for a size of $50$ to $0.1$ for a size of $250$. 

The sharing factor was also much higher and more concentrated (explained using LLN) in larger ensembles. Figure~\ref{fig:neat_ncl_unconnected_average_sharing_factor_against_generation_index_0_480}, Figure~\ref{fig:neat_ncl_unconnected_average_sharing_factor_against_average_niche_radius}, Figure~\ref{fig:neat_ncl_unconnected_average_adjusted_fitness_against_average_sharing_factor} and Figure~\ref{fig:neat_ncl_unconnected_population_diversity_against_average_sharing_factor} all highlighted this connection. $\sigma$ was dynamically computed as the average distance between all base learners and could be thought as a probability threshold that distinguished half the population of base learners. A formula $p \dot m$ is created, where $p$ is a constant. Since the sharing factor increased whenever models were close to each other (i.e. distance was smaller than $\sigma$), having more models naturally led to more increments. Essentially, the sharing factor scaled with the ensemble size ($m$), resulting in values around $35$ for the largest group and only $10$ for the smallest. 

Interestingly, increasing the ensemble size had no direct effect on diversity across base learners and their populations. Figure~\ref{fig:neat_ncl_unconnected_population_diversity_against_average_sharing_factor} showed that $D$ had value around $0.05$ despite ensemble size. From Figure~\ref{fig:neat_ncl_unconnected_mse_against_diversity_coefficient_fixed_ensemble_size_50}, Figure~\ref{fig:neat_ncl_unconnected_mse_against_diversity_coefficient_fixed_ensemble_size_100}, Figure~\ref{fig:neat_ncl_unconnected_mse_against_diversity_coefficient_fixed_ensemble_size_150}, Figure~\ref{fig:neat_ncl_unconnected_mse_against_diversity_coefficient_fixed_ensemble_size_200} and Figure~\ref{fig:neat_ncl_unconnected_mse_against_diversity_coefficient_fixed_ensemble_size_250}, increasing ensemble size did not change $\rho$. Both $D$ and $\rho$ measured the average distance. This suggested that the average distances across base learners and populations were not affected by ensemble size. 

\paragraph{Conclusions from Experiment}
A larger ensemble size resulted in fewer fluctuations in the average values of the ensemble models' attributes. It also leads to worse individual performance of base learners, a smaller $D_{max}$ and a larger sharing factor. Finally, it had no direct effect on diversity across base learners and their populations. 


\subsubsection{What is the impact of population diversity? }
\paragraph{Experiments and Observations}
When population diversity $D$ increased, the correlation penalty coefficient $\lambda$ decreased. Figure~\ref{fig:neat_ncl_unconnected_correlation_penalty_coefficient_against_population_diversity} showed multiple straight lines with negative slopes, confirming this inverse relationship between $D$ and $\lambda$. Since $D_{max}$ usually stayed constant, as seen in Figure~\ref{fig:neat_ncl_unconnected_max_population_diversity_against_generation_index_0_480}, the formula for $\lambda$ functioned like a linear equation ($\lambda = -(\frac{\lambda_{max} - \lambda_{min}}{D_{max}})D + \lambda_{max}$). The straight, negatively-sloped lines in the data proved that the dynamic $\lambda$ was implemented correctly. The reason different ensemble sizes had different slopes was that $D_{max}$ changed based on the number of learners (as illustrated by Figure~\ref{fig:neat_ncl_unconnected_max_population_diversity_against_generation_index_0_480}), which altered the steepness of the linear decrease. 

However, higher population diversity $D$ also led to a higher ensemble MSE, meaning performance actually got worse. Figure~\ref{fig:neat_ncl_unconnected_mse_against_population_diversity} illustrated this positive correlation between ensemble MSE and $D$. This suggests that while diversity was necessary, forcing the models to be too different from one another eventually damaged the accuracy of the group's collective prediction. 

\paragraph{Conclusions from Experiment}
Increasing population diversity leads to a linear decrease in $\lambda$ but an increase in ensemble MSE, i.e. worse ensemble performance. 


\subsubsection{Is NEAT-NCL the better architecture? }
\paragraph{Experiments and Observations}
NEAT-NCL performed significantly worse than other architectures like the MLP, the traditional ensemble model and the standard NCL ensemble. While other models reached a low MSE of approximately $0.05$, Figure~\ref{fig:neat_ncl_unconnected_mse_against_generation_index_0_480_fixed_ensemble_size_50}, Figure~\ref{fig:neat_ncl_unconnected_mse_against_generation_index_0_480_fixed_ensemble_size_100}, Figure~\ref{fig:neat_ncl_unconnected_mse_against_generation_index_0_480_fixed_ensemble_size_150}, Figure~\ref{fig:neat_ncl_unconnected_mse_against_generation_index_0_480_fixed_ensemble_size_200} and Figure~\ref{fig:neat_ncl_unconnected_mse_against_generation_index_0_480_fixed_ensemble_size_250} showed that the NEAT-NCL ensemble stabilised at a much higher MSE of $0.35$. 

This poor performance was caused by the extremely high diversity encouraged within the NEAT-NCL population. As shown in Figure~\ref{fig:neat_ncl_unconnected_mse_against_diversity_coefficient_fixed_ensemble_size_50}, Figure~\ref{fig:neat_ncl_unconnected_mse_against_diversity_coefficient_fixed_ensemble_size_100}, Figure~\ref{fig:neat_ncl_unconnected_mse_against_diversity_coefficient_fixed_ensemble_size_150}, Figure~\ref{fig:neat_ncl_unconnected_mse_against_diversity_coefficient_fixed_ensemble_size_200} and Figure~\ref{fig:neat_ncl_unconnected_mse_against_diversity_coefficient_fixed_ensemble_size_250}, NEAT-NCL had a diversity coefficient $\rho$ ranging from $0.001$ to $0.2$, which was significantly higher than the $0.001$ for the traditional ensemble model or the $0.003$ for the standard NCL ensemble model. 

There was a clear positive correlation between this high diversity and the error rate, as illustrated by Figure~\ref{fig:neat_ncl_unconnected_mse_against_diversity_coefficient_fixed_ensemble_size_50}, Figure~\ref{fig:neat_ncl_unconnected_mse_against_diversity_coefficient_fixed_ensemble_size_100}, Figure~\ref{fig:neat_ncl_unconnected_mse_against_diversity_coefficient_fixed_ensemble_size_150}, Figure~\ref{fig:neat_ncl_unconnected_mse_against_diversity_coefficient_fixed_ensemble_size_200} and Figure~\ref{fig:neat_ncl_unconnected_mse_against_diversity_coefficient_fixed_ensemble_size_250}. As the diversity coefficient increased, the ensemble MSE of NEAT-NCL also increased. This indicated that the models became so different from one another that they could no longer produce an accurate collective prediction, leading to the observed drop in overall performance. 

\paragraph{Conclusions from Experiment}
NEAT-NCL is not the best architecture for the current task. Its base learners are more diverse than base learners of other architectures, e.g. an arithmetic-mean-aggregator and a NCL ensemble model, leading to a lower ensemble MSE. This is supported by the positive correlation between MSE and $\rho$. 


\subsubsection{What is the best ensemble size for a NEAT-NCL ensemble model? }
\paragraph{Experiments and Observations}
Smaller NEAT-NCL ensemble sizes consistently showed better overall performance. According to Figure~\ref{fig:neat_ncl_unconnected_mse_against_generation_index_0_480_fixed_ensemble_size_50}, Figure~\ref{fig:neat_ncl_unconnected_mse_against_generation_index_0_480_fixed_ensemble_size_100}, Figure~\ref{fig:neat_ncl_unconnected_mse_against_generation_index_0_480_fixed_ensemble_size_150}, Figure~\ref{fig:neat_ncl_unconnected_mse_against_generation_index_0_480_fixed_ensemble_size_200} and Figure~\ref{fig:neat_ncl_unconnected_mse_against_generation_index_0_480_fixed_ensemble_size_250}, an ensemble size of $50$ achieved an MSE of approximately $0.3$, which was slightly lower than the $0.35$ seen in larger groups. This trend was clearly visible in Figure~\ref{fig:neat_ncl_unconnected_mse_against_population_diversity}, where the smallest ensemble size (the blue dots) maintained the lowest population diversity at around $0.05$, leading to a more accurate collective prediction. 

This superior performance was also driven by the higher quality of individual models within smaller groups. Figure~\ref{fig:neat_ncl_unconnected_average_raw_fitness_against_generation_index_0_480}, Figure~\ref{fig:neat_ncl_unconnected_average_adjusted_fitness_against_generation_index_0_480}, Figure~\ref{fig:neat_ncl_partial_direct_average_raw_fitness_against_generation_index_0_480}, Figure~\ref{fig:neat_ncl_partial_direct_average_adjusted_fitness_against_generation_index_0_480}, Figure~\ref{fig:neat_ncl_full_direct_average_raw_fitness_against_generation_index_0_480} and Figure~\ref{fig:neat_ncl_full_direct_average_adjusted_fitness_against_generation_index_0_480} all showed that base learners achieved their highest raw and adjusted fitness when the ensemble size was kept small. Essentially, by reducing the pressure to be extremely diverse, smaller ensembles allowed their individual members to focus more on accuracy, resulting in a better final result. 

\paragraph{Conclusions from Experiment}
Smallest NEAT-NCL ensemble models perform better and thus, the best ensemble size in this experiment is 50. Its better performance compared to other NEAT-NCL ensemble models is due to lower population diversity and higher fitness of base learners. 


\subsubsection{How many hidden nodes should be in a base learner for this problem? }
\paragraph{Experiments and Observations}
For this specific problem, the most effective base learner architecture proved to be one with no hidden nodes, where input nodes connected directly to the output. Figure~\ref{fig:neat_ncl_unconnected_average_active_hidden_node_num_against_generation_index_0_480} showed that the number of active hidden nodes remained at zero throughout the entire evolutionary process. This indicated that the strongest models and their offspring consistently thrived without adding internal complexity. 

To confirm this was not just a result of starting with empty networks, the experiment was repeated using different setups like \lstinline|partial_direct 0.5| and \lstinline|full_direct|. Even when models started with many connections, they eventually evolved back toward having few or no hidden nodes, as seen in Figure~\ref{fig:neat_ncl_partial_direct_average_active_hidden_node_num_against_generation_index_0_480} and Figure~\ref{fig:neat_ncl_full_direct_average_active_hidden_node_num_against_generation_index_0_480}. This matched the observation in Section~\ref{sec:stage_1} that adding hidden layers actually hurt performance. Both findings suggested that a simpler model structure was better for this task. 

\paragraph{Conclusions from Experiment}
A NEAT-NCL base learner should have 0 hidden nodes for this regression task. 


% ================================ Images ================================


% 2 Image
\begin{figure}[H]
    \centering

    \begin{minipage}[t]{0.48\textwidth}
        \vspace{0pt}
        \centering
        \includegraphics[width=\linewidth]{img/neat-ncl/unconnected/neat_ncl_average_raw_fitness_against_generation_index_0_480.png}
        \caption{Average raw fitness of a base learner throughout generations of evolution (\lstinline[language=ini]|initial_connection = unconnected|)}
        \label{fig:neat_ncl_unconnected_average_raw_fitness_against_generation_index_0_480}
    \end{minipage}
    \hfill
    \begin{minipage}[t]{0.48\textwidth}
        \vspace{0pt}
        \centering
        \includegraphics[width=\linewidth]{img/neat-ncl/unconnected/neat_ncl_average_adjusted_fitness_against_generation_index_0_480.png}
        \caption{Average adjusted fitness of a base learner throughout generations of evolution (\lstinline[language=ini]|initial_connection = unconnected|)}
        \label{fig:neat_ncl_unconnected_average_adjusted_fitness_against_generation_index_0_480}
    \end{minipage}
\end{figure}


% 2 Image
\begin{figure}[H]
    \centering

    \begin{minipage}[t]{0.48\textwidth}
        \vspace{0pt}
        \centering
        \includegraphics[width=\linewidth]{img/neat-ncl/partial-direct/neat_ncl_average_raw_fitness_against_generation_index_0_480.png}
        \caption{Average raw fitness of a base learner throughout generations of evolution (\lstinline[language=ini]|initial_connection = partial_direct 0.5|)}
        \label{fig:neat_ncl_partial_direct_average_raw_fitness_against_generation_index_0_480}
    \end{minipage}
    \hfill
    \begin{minipage}[t]{0.48\textwidth}
        \vspace{0pt}
        \centering
        \includegraphics[width=\linewidth]{img/neat-ncl/partial-direct/neat_ncl_average_adjusted_fitness_against_generation_index_0_480.png}
        \caption{Average adjusted fitness of a base learner throughout generations of evolution (\lstinline[language=ini]|initial_connection = partial_direct 0.5|)}
        \label{fig:neat_ncl_partial_direct_average_adjusted_fitness_against_generation_index_0_480}
    \end{minipage}
\end{figure}


% 2 Image
\begin{figure}[H]
    \centering

    \begin{minipage}[t]{0.48\textwidth}
        \vspace{0pt}
        \centering
        \includegraphics[width=\linewidth]{img/neat-ncl/full-direct/neat_ncl_average_raw_fitness_against_generation_index_0_480.png}
        \caption{Average raw fitness of a base learner throughout generations of evolution (\lstinline[language=ini]|initial_connection = full_direct|)}
        \label{fig:neat_ncl_full_direct_average_raw_fitness_against_generation_index_0_480}
    \end{minipage}
    \hfill
    \begin{minipage}[t]{0.48\textwidth}
        \vspace{0pt}
        \centering
        \includegraphics[width=\linewidth]{img/neat-ncl/full-direct/neat_ncl_average_adjusted_fitness_against_generation_index_0_480.png}
        \caption{Average adjusted fitness of a base learner throughout generations of evolution (\lstinline[language=ini]|initial_connection = full_direct|)}
        \label{fig:neat_ncl_full_direct_average_adjusted_fitness_against_generation_index_0_480}
    \end{minipage}
\end{figure}


% 2 Image
\begin{figure}[H]
    \centering

    \begin{minipage}[t]{0.48\textwidth}
        \vspace{0pt}
        \centering
        \includegraphics[width=\linewidth]{img/neat-ncl/unconnected/neat_ncl_average_niche_radius_against_generation_index_0_480.png}
        \caption{Average niche radius throughout generations of evolution (\lstinline[language=ini]|initial_connection = unconnected|)}
        \label{fig:neat_ncl_unconnected_average_niche_radius_against_generation_index_0_480}
    \end{minipage}
    \hfill
    \begin{minipage}[t]{0.48\textwidth}
        \vspace{0pt}
        \centering
        \includegraphics[width=\linewidth]{img/neat-ncl/unconnected/neat_ncl_average_sharing_factor_against_generation_index_0_480.png}
        \caption{Average sharing factor of a base learner throughout generations of evolution (\lstinline[language=ini]|initial_connection = unconnected|)}
        \label{fig:neat_ncl_unconnected_average_sharing_factor_against_generation_index_0_480}
    \end{minipage}
\end{figure}


% 2 Image
\begin{figure}[H]
    \centering

    \begin{minipage}[t]{0.48\textwidth}
        \vspace{0pt}
        \centering
        \includegraphics[width=\linewidth]{img/neat-ncl/unconnected/neat_ncl_population_diversity_against_generation_index_0_480.png}
        \caption{$D$ throughout generations of evolution (\lstinline[language=ini]|initial_connection = unconnected|)}
        \label{fig:neat_ncl_unconnected_population_diversity_against_generation_index_0_480}
    \end{minipage}
    \hfill
    \begin{minipage}[t]{0.48\textwidth}
        \vspace{0pt}
        \centering
        \includegraphics[width=\linewidth]{img/neat-ncl/unconnected/neat_ncl_max_population_diversity_against_generation_index_0_480.png}
        \caption{$D_{max}$ throughout generations of evolution (\lstinline[language=ini]|initial_connection = unconnected|)}
        \label{fig:neat_ncl_unconnected_max_population_diversity_against_generation_index_0_480}
    \end{minipage}
\end{figure}


% 2 Image
\begin{figure}[H]
    \centering

    \begin{minipage}[t]{0.48\textwidth}
        \vspace{0pt}
        \centering
        \includegraphics[width=\linewidth]{img/neat-ncl/unconnected/neat_ncl_correlation_penalty_coefficient_against_generation_index_0_480.png}
        \caption{$\lambda$ throughout generations of evolution (\lstinline[language=ini]|initial_connection = unconnected|)}
        \label{fig:neat_ncl_unconnected_correlation_penalty_coefficient_against_generation_index_0_480}
    \end{minipage}
    \hfill
    \begin{minipage}[t]{0.48\textwidth}
        \vspace{0pt}
        \centering
        \includegraphics[width=\linewidth]{img/neat-ncl/unconnected/neat_ncl_average_sharing_factor_against_average_niche_radius.png}
        \caption{Impact of niche radius on the sharing factor of a base learner (\lstinline[language=ini]|initial_connection = unconnected|)}
        \label{fig:neat_ncl_unconnected_average_sharing_factor_against_average_niche_radius}
    \end{minipage}
\end{figure}


% 2 Image
\begin{figure}[H]
    \centering

    \begin{minipage}[t]{0.48\textwidth}
        \vspace{0pt}
        \centering
        \includegraphics[width=\linewidth]{img/neat-ncl/unconnected/neat_ncl_average_adjusted_fitness_against_average_niche_radius.png}
        \caption{Impact of niche radius on the adjusted fitness of a base learner (\lstinline[language=ini]|initial_connection = unconnected|)}
        \label{fig:neat_ncl_unconnected_average_adjusted_fitness_against_average_niche_radius}
    \end{minipage}
    \hfill
    \begin{minipage}[t]{0.48\textwidth}
        \vspace{0pt}
        \centering
        \includegraphics[width=\linewidth]{img/neat-ncl/unconnected/neat_ncl_population_diversity_against_average_niche_radius.png}
        \caption{Impact of niche radius on population diversity (\lstinline[language=ini]|initial_connection = unconnected|)}
        \label{fig:neat_ncl_unconnected_population_diversity_against_average_niche_radius}
    \end{minipage}
\end{figure}


% 2 Image
\begin{figure}[H]
    \centering

    \begin{minipage}[t]{0.48\textwidth}
        \vspace{0pt}
        \centering
        \includegraphics[width=\linewidth]{img/neat-ncl/unconnected/neat_ncl_average_adjusted_fitness_against_average_sharing_factor.png}
        \caption{Impact of sharing factor on the adjusted fitness of a base learner (\lstinline[language=ini]|initial_connection = unconnected|)}
        \label{fig:neat_ncl_unconnected_average_adjusted_fitness_against_average_sharing_factor}
    \end{minipage}
    \hfill
    \begin{minipage}[t]{0.48\textwidth}
        \vspace{0pt}
        \centering
        \includegraphics[width=\linewidth]{img/neat-ncl/unconnected/neat_ncl_population_diversity_against_average_sharing_factor.png}
        \caption{Impact of sharing factor on population diversity (\lstinline[language=ini]|initial_connection = unconnected|)}
        \label{fig:neat_ncl_unconnected_population_diversity_against_average_sharing_factor}
    \end{minipage}
\end{figure}


% 2 Image
\begin{figure}[H]
    \centering

    \begin{minipage}[t]{0.48\textwidth}
        \vspace{0pt}
        \centering
        \includegraphics[width=\linewidth]{img/neat-ncl/unconnected/neat_ncl_correlation_penalty_coefficient_against_population_diversity_ratio.png}
        \caption{Impact of $\frac{D}{D_{max}}$ on $\lambda$ (\lstinline[language=ini]|initial_connection = unconnected|)}
        \label{fig:neat_ncl_unconnected_correlation_penalty_coefficient_against_population_diversity_ratio}
    \end{minipage}
    \hfill
    \begin{minipage}[t]{0.48\textwidth}
        \vspace{0pt}
        \centering
        \includegraphics[width=\linewidth]{img/neat-ncl/unconnected/neat_ncl_correlation_penalty_coefficient_against_population_diversity.png}
        \caption{Impact of $D$ on $\lambda$ (\lstinline[language=ini]|initial_connection = unconnected|)}
        \label{fig:neat_ncl_unconnected_correlation_penalty_coefficient_against_population_diversity}
    \end{minipage}
\end{figure}


% 2 Image
\begin{figure}[H]
    \centering

    \begin{minipage}[t]{0.48\textwidth}
        \vspace{0pt}
        \centering
        \includegraphics[width=\linewidth]{img/neat-ncl/unconnected/neat_ncl_mse_against_population_diversity.png}
        \caption{Comparison of ensemble MSE with population diversity (\lstinline[language=ini]|initial_connection = unconnected|)}
        \label{fig:neat_ncl_unconnected_mse_against_population_diversity}
    \end{minipage}
    \hfill
    \begin{minipage}[t]{0.48\textwidth}
        \vspace{0pt}
        \centering
        \includegraphics[width=\linewidth]{img/neat-ncl/unconnected/neat_ncl_average_active_hidden_node_num_against_generation_index_0_480.png}
        \caption{Average number of active hidden nodes of a base learner throughout generations of evolution (\lstinline[language=ini]|initial_connection = unconnected|)}
        \label{fig:neat_ncl_unconnected_average_active_hidden_node_num_against_generation_index_0_480}
    \end{minipage}
\end{figure}


% 2 Image
\begin{figure}[H]
    \centering

    \begin{minipage}[t]{0.48\textwidth}
        \vspace{0pt}
        \centering
        \includegraphics[width=\linewidth]{img/neat-ncl/partial-direct/neat_ncl_average_active_hidden_node_num_against_generation_index_0_480.png}
        \caption{Average number of active hidden nodes of a base learner throughout generations of evolution (\lstinline[language=ini]|initial_connection = partial_direct 0.5|)}
        \label{fig:neat_ncl_partial_direct_average_active_hidden_node_num_against_generation_index_0_480}
    \end{minipage}
    \hfill
    \begin{minipage}[t]{0.48\textwidth}
        \vspace{0pt}
        \centering
        \includegraphics[width=\linewidth]{img/neat-ncl/full-direct/neat_ncl_average_active_hidden_node_num_against_generation_index_0_480.png}
        \caption{Average number of active hidden nodes of a base learner throughout generations of evolution (\lstinline[language=ini]|initial_connection = full_direct|)}
        \label{fig:neat_ncl_full_direct_average_active_hidden_node_num_against_generation_index_0_480}
    \end{minipage}
\end{figure}


% 2 Image
\begin{figure}[H]
    \centering

    \begin{minipage}[t]{0.48\textwidth}
        \vspace{0pt}
        \centering
        \includegraphics[width=\linewidth]{img/neat-ncl/unconnected/neat_ncl_mse_against_diversity_coefficient_fixed_ensemble_size_50.png}
        \caption{Comparison of ensemble MSE with diversity coefficient when the ensemble size is 50 (\lstinline[language=ini]|initial_connection = unconnected|)}
        \label{fig:neat_ncl_unconnected_mse_against_diversity_coefficient_fixed_ensemble_size_50}
    \end{minipage}
    \hfill
    \begin{minipage}[t]{0.48\textwidth}
        \vspace{0pt}
        \centering
        \includegraphics[width=\linewidth]{img/neat-ncl/unconnected/neat_ncl_mse_against_generation_index_0_480_fixed_ensemble_size_50.png}
        \caption{Ensemble MSE throughout generations of evolution when the ensemble size is 50 (\lstinline[language=ini]|initial_connection = unconnected|)}
        \label{fig:neat_ncl_unconnected_mse_against_generation_index_0_480_fixed_ensemble_size_50}
    \end{minipage}
\end{figure}


% 2 Image
\begin{figure}[H]
    \centering

    \begin{minipage}[t]{0.48\textwidth}
        \vspace{0pt}
        \centering
        \includegraphics[width=\linewidth]{img/neat-ncl/unconnected/neat_ncl_mse_against_diversity_coefficient_fixed_ensemble_size_100.png}
        \caption{Comparison of ensemble MSE with diversity coefficient when the ensemble size is 100 (\lstinline[language=ini]|initial_connection = unconnected|)}
        \label{fig:neat_ncl_unconnected_mse_against_diversity_coefficient_fixed_ensemble_size_100}
    \end{minipage}
    \hfill
    \begin{minipage}[t]{0.48\textwidth}
        \vspace{0pt}
        \centering
        \includegraphics[width=\linewidth]{img/neat-ncl/unconnected/neat_ncl_mse_against_generation_index_0_480_fixed_ensemble_size_100.png}
        \caption{Ensemble MSE throughout generations of evolution when the ensemble size is 100 (\lstinline[language=ini]|initial_connection = unconnected|)}
        \label{fig:neat_ncl_unconnected_mse_against_generation_index_0_480_fixed_ensemble_size_100}
    \end{minipage}
\end{figure}


% 2 Image
\begin{figure}[H]
    \centering

    \begin{minipage}[t]{0.48\textwidth}
        \vspace{0pt}
        \centering
        \includegraphics[width=\linewidth]{img/neat-ncl/unconnected/neat_ncl_mse_against_diversity_coefficient_fixed_ensemble_size_150.png}
        \caption{Comparison of ensemble MSE with diversity coefficient when the ensemble size is 150 (\lstinline[language=ini]|initial_connection = unconnected|)}
        \label{fig:neat_ncl_unconnected_mse_against_diversity_coefficient_fixed_ensemble_size_150}
    \end{minipage}
    \hfill
    \begin{minipage}[t]{0.48\textwidth}
        \vspace{0pt}
        \centering
        \includegraphics[width=\linewidth]{img/neat-ncl/unconnected/neat_ncl_mse_against_generation_index_0_480_fixed_ensemble_size_150.png}
        \caption{Ensemble MSE throughout generations of evolution when the ensemble size is 150 (\lstinline[language=ini]|initial_connection = unconnected|)}
        \label{fig:neat_ncl_unconnected_mse_against_generation_index_0_480_fixed_ensemble_size_150}
    \end{minipage}
\end{figure}


% 2 Image
\begin{figure}[H]
    \centering

    \begin{minipage}[t]{0.48\textwidth}
        \vspace{0pt}
        \centering
        \includegraphics[width=\linewidth]{img/neat-ncl/unconnected/neat_ncl_mse_against_diversity_coefficient_fixed_ensemble_size_200.png}
        \caption{Comparison of ensemble MSE with diversity coefficient when the ensemble size is 200 (\lstinline[language=ini]|initial_connection = unconnected|)}
        \label{fig:neat_ncl_unconnected_mse_against_diversity_coefficient_fixed_ensemble_size_200}
    \end{minipage}
    \hfill
    \begin{minipage}[t]{0.48\textwidth}
        \vspace{0pt}
        \centering
        \includegraphics[width=\linewidth]{img/neat-ncl/unconnected/neat_ncl_mse_against_generation_index_0_480_fixed_ensemble_size_200.png}
        \caption{Ensemble MSE throughout generations of evolution when the ensemble size is 200 (\lstinline[language=ini]|initial_connection = unconnected|)}
        \label{fig:neat_ncl_unconnected_mse_against_generation_index_0_480_fixed_ensemble_size_200}
    \end{minipage}
\end{figure}


% 2 Image
\begin{figure}[H]
    \centering

    \begin{minipage}[t]{0.48\textwidth}
        \vspace{0pt}
        \centering
        \includegraphics[width=\linewidth]{img/neat-ncl/unconnected/neat_ncl_mse_against_diversity_coefficient_fixed_ensemble_size_250.png}
        \caption{Comparison of ensemble MSE with diversity coefficient when the ensemble size is 250 (\lstinline[language=ini]|initial_connection = unconnected|)}
        \label{fig:neat_ncl_unconnected_mse_against_diversity_coefficient_fixed_ensemble_size_250}
    \end{minipage}
    \hfill
    \begin{minipage}[t]{0.48\textwidth}
        \vspace{0pt}
        \centering
        \includegraphics[width=\linewidth]{img/neat-ncl/unconnected/neat_ncl_mse_against_generation_index_0_480_fixed_ensemble_size_250.png}
        \caption{Ensemble MSE throughout generations of evolution when the ensemble size is 250 (\lstinline[language=ini]|initial_connection = unconnected|)}
        \label{fig:neat_ncl_unconnected_mse_against_generation_index_0_480_fixed_ensemble_size_250}
    \end{minipage}
\end{figure}


\section{Stage 5: Using different aggregators}
\label{sec:stage_5}


\subsection{Motivations}
In Section~\ref{sec:stage_2}, Section~\ref{sec:stage_3} and Section~\ref{sec:stage_4}, an arithmetic-mean-aggregator was used to combine predictions from all base learners to give the final ensemble prediction. Every base learner was treated equally important because they have the same weight and contribution to the final ensemble prediction. However, how will the ensemble performance change if different base learners have different importance and contributions to the final ensemble prediction? Instead of calculating the average predictions, what if the ensemble prediction is the \textbf{median} of all predictions (which is more robust against data skewness)? In this stage, different aggregator types, i.e. an arithmetic-mean-aggregator, a median-aggregator and a neural-network-aggregator, were used and their ensemble performance was compared. This section also provided reference data to Section~\ref{sec:stage_6}, in which the ensemble performance of different architectures given a specific aggregator was studied. 


\subsection{Design and Implementation}


\subsubsection{Updated loss functions of base learners and aggregator in a NCL ensemble model}
When a neural-network-aggregator is used, the loss function of a base learners is updated using the formula below: 
\begin{equation}
    \label{equation:recursive_nn_loss_function_base_learner}
    \mathcal{L}_i = (\hat{y}_i - y)^2 - \lambda (\hat{y}_i - \hat{y}_{ensemble})^2
\end{equation}
, where $\hat{y}_{ensemble}$ is the output of the neural-network-aggregator. 


When a neural-network-aggregator is used, the loss function of the neural-network-aggregator is updated using the formula below: 
\begin{equation}
    \label{equation:recursive_nn_loss_function_aggregator}
    \mathcal{L}_{ensemble} = (\hat{y}_{ensemble} - y)^2
\end{equation}
, where $\hat{y}_{ensemble}$ is the output of the neural-network-aggregator. 
\\
\\
Backpropagation on the neural-network-aggregator should not be propagated to base learners to avoid double-backpropagation. This can be achieved by adding \texttt{.detach()} to the output tensors (which will be passed to the neural-network-aggregator) returned by the base learners. 


\subsubsection{Model construction}
Each ensemble model, such as \texttt{TraditionalEnsembleNet} and \texttt{StaticNCLEnsembleNet}, used a variable called \texttt{\_\_voter} to define how the base learners' results were combined. While \texttt{ArithmeticMeanVoter} was used in previous tests to simply average the results, two new options had been introduced to explore different aggregation methods. 

\texttt{MedianVoter} acted as a median-aggregator by using \texttt{torch.median(y\_predictions, dim=0).values} to find the middle value of all base learner predictions. This approach could be more robust than a simple average because it was less affected by extreme outliers that produced very high or very low errors. 

\texttt{NNVoter} used a simple neural network to learn the best way to combine predictions. It consisted of $m$ input features (one for each base learner) and a single output for the final ensemble prediction. Since it had no hidden layers and no dropout, it essentially learnt a specific weight for each base learner based on their individual performance, allowing the ensemble to trust its most accurate members more than the others. 


\subsubsection{Experiment}
There were 3 different types of aggregators, i.e. \texttt{ArithmeticMeanVoter}, \texttt{MedianVoter} and \texttt{NNVoter}. There were 3 different types of ensemble model, i.e. a traditional ensemble model in Section~\ref{sec:stage_2}, a NCL ensemble model in Section~\ref{sec:stage_3} and a NEAT-NCL ensemble model in Section~\ref{sec:stage_4}. However, only the NEAT-NCL ensemble model was not given \texttt{NNVoter} as an aggregator. Unlike traditional ensemble model and NCL ensemble model, base learners in a NEAT-NCL ensemble model change their architectures dynamically over generations and their optimisations are not gradient-based. It does not make sense to apply backpropagation to update the neural-network-aggregator every generation because every generation gives a new population of base learners which the updated aggregator no longer applies. Thus, the NEAT-NCL ensemble model was only given \texttt{ArithmeticMeanVoter} and \texttt{MedianVoter} as its aggregator. 

\paragraph{Varied hyperparameters}
Table~\ref{table:experiment_5_varied_hyperparameters} contains all varied hyperparameters in this experiment: 

\begin{table}[H]
    \centering
    \begin{tabular}{|c|c|p{3cm}|p{5cm}|}
        \hline
        Hyperparameter & Type & Choices / Ranges & Description \\
        \hline
        \small\texttt{correlation\_penalty\_coefficients} & Float & $[0, 1.5]\text{, step}=0.1$ & $\lambda$ (for NCL ensemble model only) \\
        \hline
        \small\texttt{ensemble\_sizes} & Integer & \{$[2, 16]\text{, step}=1$, $[50, 250]\text{, step}=50$\} & $m$ ($[50, 250]$ for NEAT-NCL ensemble model and $[2, 16]$ for the others) \\
        \hline
        \small\texttt{hidden\_sizes} & Python list & \{\texttt{[2]}, \texttt{[2, 2, 2]}, \texttt{[2, 2, 2, 2, 2, 2]}, \texttt{[6]}, \texttt{[12]}\} & Model architecture \\
        \hline
        \small\texttt{voter} & Categorical & \{\texttt{arithmetic\_mean}, \texttt{median}, \texttt{nn}\} & Aggregator type (\texttt{nn} is not applicable for NEAT-NCL ensemble model) \\
        \hline
    \end{tabular}
    \caption{Varied hyperparameters and their configurations}
    \label{table:experiment_5_varied_hyperparameters}
\end{table}

\paragraph{Fixed hyperparameters}
Table~\ref{table:experiment_5_fixed_hyperparameters} contains all fixed hyperparameters in this experiment: 

\begin{table}[H]
    \centering
    \begin{tabular}{|c|c|c|p{5cm}|}
        \hline
        Hyperparameter & Type & Choices / Ranges & Description \\
        \hline
        \small\texttt{input\_size} & Integer & $13$ & Number of input features of the MLP \\
        \hline
        \small\texttt{output\_size} & Integer & $1$ & Number of labels of the MLP \\
        \hline
        \small\texttt{optimizer} & Categorical & \texttt{torch.optim.Adam} & Optimisation algorithm used for training \\
        \hline
        \small\texttt{learning\_rate} & Float & $0.001$ & Learning rate of the optimisation algorithm \\
        \hline
        \small\texttt{dropout\_rate} & Float & $0.5$ & Dropout rate applied to the data tensor after the input layer \\
        \hline
        \small\texttt{activation} & Categorical & ReLU & Activation function used in each dense layer of the MLP \\
        \hline
        \small\texttt{epoch\_nums} & Integer & $400$ & Number of training epochs (for arithmetic-mean-aggregator and NCL ensemble model only) \\
        \hline
        \small\texttt{evolution\_epoch} & Integer & $400$ & Number of evolutionary generations during training (for NEAT-NCL ensemble model only) \\
        \hline
        \small\texttt{batch\_size} & Integer & \{$32$, $506$\} & Number of training samples passed to the model in one go ($506$ for NEAT-NCL ensemble model and $32$ for the others) \\
        \hline
    \end{tabular}
    \caption{Fixed hyperparameters and their configurations}
    \label{table:experiment_5_fixed_hyperparameters}
\end{table}

\paragraph{Metric}
Table~\ref{table:experiment_5_metric} contains all evaluation metrics collected in this experiment: 

\begin{table}[H]
    \centering
    \begin{tabular}{|c|c|c|}
        \hline
        Metric & Type & Description \\
        \hline
        \small\texttt{loss\_test} & Float & MSE of ensemble model during testing \\
        \hline
    \end{tabular}
    \caption{Evaluation metrics}
    \label{table:experiment_5_metric}
\end{table}

\paragraph{Further notes}
All output metrics collected are stored in \path{experiment_runner_traditional_arithmetic_mean.csv}, \path{experiment_runner_traditional_median.csv}, \path{experiment_runner_traditional_nn.csv}, \path{experiment_runner_static_ncl_arithmetic_mean.csv}, \path{experiment_runner_static_ncl_median.csv}, \path{experiment_runner_static_ncl_nn.csv}, \path{experiment_runner_neat_ncl_arithmetic_mean.csv} and \path{experiment_runner_neat_ncl_median.csv}. 


\subsection{Expected Observations}


\subsubsection{Arithmetic-mean-aggregator and median-aggregator are more similar and better}
Among the three types of aggregators, the arithmetic-mean-aggregator and the median-aggregator are expected to produce more similar predictions than the neural-network-aggregator. This is because both the mean and median methods treat the contributions of base learners relatively equally. 

In contrast, the neural-network-aggregator assigns different weights to each base learner, causing the final prediction to rely more heavily on specific groups. Applying a neural-network-aggregator on a NCL ensemble model, where base learners are forced to be diverse to cover different aspects of the training data, should lead to worse ensemble performance because it can end up putting too much attention on inaccurate base learners, which prevents the individual mistakes from being averaged out effectively. This lead to a higher ensemble MSE compared to the simpler mean or median methods. 


\subsubsection{Small ensemble MSE despite aggregator type used}
Since both the training and testing data were normalised using \texttt{MinMaxScaler}, all input features and labels stayed within a very small, consistent range. Since the target values were scaled down, the errors $\left|\hat{y}_i - y\right|$ (and thus MSE) are expected to be small as well. Regardless of whether the ensemble uses an arithmetic-mean-, a median- or a neural-network-aggregator to combine predictions, the final MSE is expected to remain low because the mathematical scale of the data itself is restricted. Refer to Section~\ref{subsec:project_settings_datasets} for more details on these specific value ranges. 


\subsection{Evaluation and Empirical Comparisons}


\subsubsection{Similar ensemble performance between arithmetic-mean-aggregator and median-aggregator}
\paragraph{Experiments and Observations}
Regardless of the ensemble model used, the arithmetic-mean-aggregator and the median-aggregator showed very similar performance. In Figure~\ref{fig:voter_traditional_mse_against_ensemble_size_2_16_fixed_architecture_2_fixed_epoch_num_400} through Figure~\ref{fig:voter_traditional_mse_against_ensemble_size_2_16_fixed_architecture_2_2_2_2_2_2_fixed_epoch_num_400}, the blue and orange lines representing these two methods stayed extremely close to each other. This same pattern appeared in the NCL and NEAT-NCL experiments (Figure~\ref{fig:voter_static_ncl_mse_against_correlation_penalty_coefficient_0_1_fixed_architecture_2_fixed_epoch_num_400_fixed_ensemble_size_2} to Figure~\ref{fig:voter_static_ncl_mse_against_correlation_penalty_coefficient_0_1_fixed_architecture_2_2_2_2_2_2_fixed_epoch_num_400_fixed_ensemble_size_12} and Figure~\ref{fig:voter_neat_ncl_mse_against_generation_index_0_400_fixed_ensemble_size_50_unconnected} to Figure~\ref{fig:voter_neat_ncl_mse_against_generation_index_0_400_fixed_ensemble_size_250_full_direct} respectively), where both aggregators resulted in nearly identical MSE values. 

In certain cases, particularly in the NEAT-NCL tests shown in Figure~\ref{fig:voter_neat_ncl_mse_against_generation_index_0_400_fixed_ensemble_size_250_unconnected}, Figure~\ref{fig:voter_neat_ncl_mse_against_generation_index_0_400_fixed_ensemble_size_250_partial_direct} and Figure~\ref{fig:voter_neat_ncl_mse_against_generation_index_0_400_fixed_ensemble_size_250_full_direct}, the median-aggregator actually performed better than the arithmetic-mean-aggregator. This suggested that the median could be more effective at ignoring outliers in larger populations. However, since the two methods were so evenly matched across most other experiments, it was difficult to definitively say which one was superior overall. 

\paragraph{Conclusions from Experiment}
Overall, an arithmetic-mean-aggregator and a median-aggregator give similar ensemble performance. In a few experiments of the NEAT-NCL ensemble model, a median-aggregator performs better than an arithmetic-mean-aggregator consistently. 


\subsubsection{Poor and fluctuated ensemble performance for neural-network-aggregator}
\paragraph{Experiments and Observations}
The neural-network-aggregator consistently showed the worst performance of the three options in both traditional and NCL ensemble models. As seen in Figure~\ref{fig:voter_traditional_mse_against_ensemble_size_2_16_fixed_architecture_2_fixed_epoch_num_400} through Figure~\ref{fig:voter_static_ncl_mse_against_correlation_penalty_coefficient_0_1_fixed_architecture_2_2_2_2_2_2_fixed_epoch_num_400_fixed_ensemble_size_12}, the green line representing this method stayed much higher than the others, indicating a significantly higher MSE. 

In addition to having a higher error, the neural-network-aggregator was also much less stable. From Figure~\ref{fig:voter_traditional_mse_against_ensemble_size_2_16_fixed_architecture_2_fixed_epoch_num_400} to Figure~\ref{fig:voter_static_ncl_mse_against_correlation_penalty_coefficient_0_1_fixed_architecture_2_2_2_2_2_2_fixed_epoch_num_400_fixed_ensemble_size_12}, while the mean and median methods stayed relatively flat, the green line fluctuated heavily as hyperparameters like ensemble size $m$ or the penalty coefficient $\lambda$ changed. This high level of fluctuation suggested that using a neural network to combine predictions made the ensemble model unstable. 

\paragraph{Conclusions from Experiment}
A neural-network-aggregator gives the poorest and most unstable ensemble performance compared to other aggregators. 


\subsubsection{More stable NEAT-NCL ensemble performance for median-aggregator}
\paragraph{Experiments and Observations}
Using a median-aggregator within a NEAT-NCL ensemble model resulted in much greater stability compared to using an arithmetic-mean-aggregator. This effect became especially clear as the ensemble size $m$ increased. As shown in Figure~\ref{fig:voter_neat_ncl_mse_against_generation_index_0_400_fixed_ensemble_size_250_unconnected}, Figure~\ref{fig:voter_neat_ncl_mse_against_generation_index_0_400_fixed_ensemble_size_150_partial_direct}, Figure~\ref{fig:voter_neat_ncl_mse_against_generation_index_0_400_fixed_ensemble_size_250_partial_direct} and Figure~\ref{fig:voter_neat_ncl_mse_against_generation_index_0_400_fixed_ensemble_size_250_full_direct}, the orange line stayed much steadier than the blue line. The reason for this improved stability was how the median handles outliers. In a NEAT-NCL model, the high level of diversity could sometimes lead certain base learners to produce very inaccurate results. While an arithmetivc mean is easily pulled away by these extreme errors, the median ignores them and focuses on the middle value of the group. This allowed the ensemble to maintain a consistent performance even when some individual models failed, making it a more reliable choice for large, diverse ensembles. 

\paragraph{Conclusions from Experiment}
A median-aggregator gives a more stable NEAT-NCL ensemble model performance than an arithmetic-mean-aggregator. This phenomenon is more visible when $m$ is larger. 


% ================================ Images ================================


% 2 Image
\begin{figure}[H]
    \centering

    \begin{minipage}[t]{0.48\textwidth}
        \vspace{0pt}
        \centering
        \includegraphics[width=\linewidth]{img/voter/traditional/voter_traditional_mse_against_ensemble_size_2_16_fixed_architecture_2_fixed_epoch_num_400.png}
        \caption{Impact of ensemble size on MSE of a traditional ensemble model. Base learners had 1 layer of 2 hidden nodes and were trained with 400 epochs. }
        \label{fig:voter_traditional_mse_against_ensemble_size_2_16_fixed_architecture_2_fixed_epoch_num_400}
    \end{minipage}
    \hfill
    \begin{minipage}[t]{0.48\textwidth}
        \vspace{0pt}
        \centering
        \includegraphics[width=\linewidth]{img/voter/traditional/voter_traditional_mse_against_ensemble_size_2_16_fixed_architecture_6_fixed_epoch_num_400.png}
        \caption{Impact of ensemble size on MSE of a traditional ensemble model. Base learners had 1 layer of 6 hidden nodes and were trained with 400 epochs. }
        \label{fig:voter_traditional_mse_against_ensemble_size_2_16_fixed_architecture_6_fixed_epoch_num_400}
    \end{minipage}
\end{figure}


% 2 Image
\begin{figure}[H]
    \centering

    \begin{minipage}[t]{0.48\textwidth}
        \vspace{0pt}
        \centering
        \includegraphics[width=\linewidth]{img/voter/traditional/voter_traditional_mse_against_ensemble_size_2_16_fixed_architecture_12_fixed_epoch_num_400.png}
        \caption{Impact of ensemble size on MSE of a traditional ensemble model. Base learners had 1 layer of 12 hidden nodes and were trained with 400 epochs. }
        \label{fig:voter_traditional_mse_against_ensemble_size_2_16_fixed_architecture_12_fixed_epoch_num_400}
    \end{minipage}
    \hfill
    \begin{minipage}[t]{0.48\textwidth}
        \vspace{0pt}
        \centering
        \includegraphics[width=\linewidth]{img/voter/traditional/voter_traditional_mse_against_ensemble_size_2_16_fixed_architecture_2_2_2_fixed_epoch_num_400.png}
        \caption{Impact of ensemble size on MSE of a traditional ensemble model. Base learners had 3 layers of 2 hidden nodes and were trained with 400 epochs. }
        \label{fig:voter_traditional_mse_against_ensemble_size_2_16_fixed_architecture_2_2_2_fixed_epoch_num_400}
    \end{minipage}
\end{figure}


% 2 Image
\begin{figure}[H]
    \centering

    \begin{minipage}[t]{0.48\textwidth}
        \vspace{0pt}
        \centering
        \includegraphics[width=\linewidth]{img/voter/traditional/voter_traditional_mse_against_ensemble_size_2_16_fixed_architecture_2_2_2_2_2_2_fixed_epoch_num_400.png}
        \caption{Impact of ensemble size on MSE of a traditional ensemble model. Base learners had 6 layers of 2 hidden nodes and were trained with 400 epochs. }
        \label{fig:voter_traditional_mse_against_ensemble_size_2_16_fixed_architecture_2_2_2_2_2_2_fixed_epoch_num_400}
    \end{minipage}
    \hfill
    \begin{minipage}[t]{0.48\textwidth}
        \vspace{0pt}
        \centering
        \includegraphics[width=\linewidth]{img/voter/static-ncl/voter_static_ncl_mse_against_correlation_penalty_coefficient_0_1_fixed_architecture_2_fixed_epoch_num_400_fixed_ensemble_size_2.png}
        \caption{Impact of correlation penalty coefficient (in the range $[0, 1]$) on MSE of a NCL ensemble model. Base learners had 1 layer of 2 hidden nodes and were trained with 400 epochs. Ensemble size was 2. }
        \label{fig:voter_static_ncl_mse_against_correlation_penalty_coefficient_0_1_fixed_architecture_2_fixed_epoch_num_400_fixed_ensemble_size_2}
    \end{minipage}
\end{figure}


% 2 Image
\begin{figure}[H]
    \centering

    \begin{minipage}[t]{0.48\textwidth}
        \vspace{0pt}
        \centering
        \includegraphics[width=\linewidth]{img/voter/static-ncl/voter_static_ncl_mse_against_correlation_penalty_coefficient_0_1_fixed_architecture_6_fixed_epoch_num_400_fixed_ensemble_size_2.png}
        \caption{Impact of correlation penalty coefficient (in the range $[0, 1]$) on MSE of a NCL ensemble model. Base learners had 1 layer of 6 hidden nodes and were trained with 400 epochs. Ensemble size was 2. }
        \label{fig:voter_static_ncl_mse_against_correlation_penalty_coefficient_0_1_fixed_architecture_6_fixed_epoch_num_400_fixed_ensemble_size_2}
    \end{minipage}
    \hfill
    \begin{minipage}[t]{0.48\textwidth}
        \vspace{0pt}
        \centering
        \includegraphics[width=\linewidth]{img/voter/static-ncl/voter_static_ncl_mse_against_correlation_penalty_coefficient_0_1_fixed_architecture_12_fixed_epoch_num_400_fixed_ensemble_size_2.png}
        \caption{Impact of correlation penalty coefficient (in the range $[0, 1]$) on MSE of a NCL ensemble model. Base learners had 1 layer of 12 hidden nodes and were trained with 400 epochs. Ensemble size was 2. }
        \label{fig:voter_static_ncl_mse_against_correlation_penalty_coefficient_0_1_fixed_architecture_12_fixed_epoch_num_400_fixed_ensemble_size_2}
    \end{minipage}
\end{figure}


% 2 Image
\begin{figure}[H]
    \centering

    \begin{minipage}[t]{0.48\textwidth}
        \vspace{0pt}
        \centering
        \includegraphics[width=\linewidth]{img/voter/static-ncl/voter_static_ncl_mse_against_correlation_penalty_coefficient_0_1_fixed_architecture_2_2_2_fixed_epoch_num_400_fixed_ensemble_size_2.png}
        \caption{Impact of correlation penalty coefficient (in the range $[0, 1]$) on MSE of a NCL ensemble model. Base learners had 3 layers of 2 hidden nodes and were trained with 400 epochs. Ensemble size was 2. }
        \label{fig:voter_static_ncl_mse_against_correlation_penalty_coefficient_0_1_fixed_architecture_2_2_2_fixed_epoch_num_400_fixed_ensemble_size_2}
    \end{minipage}
    \hfill
    \begin{minipage}[t]{0.48\textwidth}
        \vspace{0pt}
        \centering
        \includegraphics[width=\linewidth]{img/voter/static-ncl/voter_static_ncl_mse_against_correlation_penalty_coefficient_0_1_fixed_architecture_2_2_2_2_2_2_fixed_epoch_num_400_fixed_ensemble_size_2.png}
        \caption{Impact of correlation penalty coefficient (in the range $[0, 1]$) on MSE of a NCL ensemble model. Base learners had 6 layers of 2 hidden nodes and were trained with 400 epochs. Ensemble size was 2. }
        \label{fig:voter_static_ncl_mse_against_correlation_penalty_coefficient_0_1_fixed_architecture_2_2_2_2_2_2_fixed_epoch_num_400_fixed_ensemble_size_2}
    \end{minipage}
\end{figure}


% 2 Image
\begin{figure}[H]
    \centering

    \begin{minipage}[t]{0.48\textwidth}
        \vspace{0pt}
        \centering
        \includegraphics[width=\linewidth]{img/voter/static-ncl/voter_static_ncl_mse_against_correlation_penalty_coefficient_0_1_fixed_architecture_2_fixed_epoch_num_400_fixed_ensemble_size_6.png}
        \caption{Impact of correlation penalty coefficient (in the range $[0, 1]$) on MSE of a NCL ensemble model. Base learners had 1 layer of 2 hidden nodes and were trained with 400 epochs. Ensemble size was 6. }
        \label{fig:voter_static_ncl_mse_against_correlation_penalty_coefficient_0_1_fixed_architecture_2_fixed_epoch_num_400_fixed_ensemble_size_6}
    \end{minipage}
    \hfill
    \begin{minipage}[t]{0.48\textwidth}
        \vspace{0pt}
        \centering
        \includegraphics[width=\linewidth]{img/voter/static-ncl/voter_static_ncl_mse_against_correlation_penalty_coefficient_0_1_fixed_architecture_6_fixed_epoch_num_400_fixed_ensemble_size_6.png}
        \caption{Impact of correlation penalty coefficient (in the range $[0, 1]$) on MSE of a NCL ensemble model. Base learners had 1 layers of 6 hidden nodes and were trained with 400 epochs. Ensemble size was 6. }
        \label{fig:voter_static_ncl_mse_against_correlation_penalty_coefficient_0_1_fixed_architecture_6_fixed_epoch_num_400_fixed_ensemble_size_6}
    \end{minipage}
\end{figure}


% 2 Image
\begin{figure}[H]
    \centering

    \begin{minipage}[t]{0.48\textwidth}
        \vspace{0pt}
        \centering
        \includegraphics[width=\linewidth]{img/voter/static-ncl/voter_static_ncl_mse_against_correlation_penalty_coefficient_0_1_fixed_architecture_12_fixed_epoch_num_400_fixed_ensemble_size_6.png}
        \caption{Impact of correlation penalty coefficient (in the range $[0, 1]$) on MSE of a NCL ensemble model. Base learners had 1 layer of 12 hidden nodes and were trained with 400 epochs. Ensemble size was 6. }
        \label{fig:voter_static_ncl_mse_against_correlation_penalty_coefficient_0_1_fixed_architecture_12_fixed_epoch_num_400_fixed_ensemble_size_6}
    \end{minipage}
    \hfill
    \begin{minipage}[t]{0.48\textwidth}
        \vspace{0pt}
        \centering
        \includegraphics[width=\linewidth]{img/voter/static-ncl/voter_static_ncl_mse_against_correlation_penalty_coefficient_0_1_fixed_architecture_2_2_2_fixed_epoch_num_400_fixed_ensemble_size_6.png}
        \caption{Impact of correlation penalty coefficient (in the range $[0, 1]$) on MSE of a NCL ensemble model. Base learners had 3 layers of 2 hidden nodes and were trained with 400 epochs. Ensemble size was 6. }
        \label{fig:voter_static_ncl_mse_against_correlation_penalty_coefficient_0_1_fixed_architecture_2_2_2_fixed_epoch_num_400_fixed_ensemble_size_6}
    \end{minipage}
\end{figure}


% 2 Image
\begin{figure}[H]
    \centering

    \begin{minipage}[t]{0.48\textwidth}
        \vspace{0pt}
        \centering
        \includegraphics[width=\linewidth]{img/voter/static-ncl/voter_static_ncl_mse_against_correlation_penalty_coefficient_0_1_fixed_architecture_2_2_2_2_2_2_fixed_epoch_num_400_fixed_ensemble_size_6.png}
        \caption{Impact of correlation penalty coefficient (in the range $[0, 1]$) on MSE of a NCL ensemble model. Base learners had 6 layer of 2 hidden nodes and were trained with 400 epochs. Ensemble size was 6. }
        \label{fig:voter_static_ncl_mse_against_correlation_penalty_coefficient_0_1_fixed_architecture_2_2_2_2_2_2_fixed_epoch_num_400_fixed_ensemble_size_6}
    \end{minipage}
    \hfill
    \begin{minipage}[t]{0.48\textwidth}
        \vspace{0pt}
        \centering
        \includegraphics[width=\linewidth]{img/voter/static-ncl/voter_static_ncl_mse_against_correlation_penalty_coefficient_0_1_fixed_architecture_2_fixed_epoch_num_400_fixed_ensemble_size_12.png}
        \caption{Impact of correlation penalty coefficient (in the range $[0, 1]$) on MSE of a NCL ensemble model. Base learners had 1 layer of 2 hidden nodes and were trained with 400 epochs. Ensemble size was 12. }
        \label{fig:voter_static_ncl_mse_against_correlation_penalty_coefficient_0_1_fixed_architecture_2_fixed_epoch_num_400_fixed_ensemble_size_12}
    \end{minipage}
\end{figure}


% 2 Image
\begin{figure}[H]
    \centering

    \begin{minipage}[t]{0.48\textwidth}
        \vspace{0pt}
        \centering
        \includegraphics[width=\linewidth]{img/voter/static-ncl/voter_static_ncl_mse_against_correlation_penalty_coefficient_0_1_fixed_architecture_6_fixed_epoch_num_400_fixed_ensemble_size_12.png}
        \caption{Impact of correlation penalty coefficient (in the range $[0, 1]$) on MSE of a NCL ensemble model. Base learners had 1 layer of 6 hidden nodes and were trained with 400 epochs. Ensemble size was 12. }
        \label{fig:voter_static_ncl_mse_against_correlation_penalty_coefficient_0_1_fixed_architecture_6_fixed_epoch_num_400_fixed_ensemble_size_12}
    \end{minipage}
    \hfill
    \begin{minipage}[t]{0.48\textwidth}
        \vspace{0pt}
        \centering
        \includegraphics[width=\linewidth]{img/voter/static-ncl/voter_static_ncl_mse_against_correlation_penalty_coefficient_0_1_fixed_architecture_12_fixed_epoch_num_400_fixed_ensemble_size_12.png}
        \caption{Impact of correlation penalty coefficient (in the range $[0, 1]$) on MSE of a NCL ensemble model. Base learners had 1 layer of 12 hidden nodes and were trained with 400 epochs. Ensemble size was 12. }
        \label{fig:voter_static_ncl_mse_against_correlation_penalty_coefficient_0_1_fixed_architecture_12_fixed_epoch_num_400_fixed_ensemble_size_12}
    \end{minipage}
\end{figure}


% 2 Image
\begin{figure}[H]
    \centering

    \begin{minipage}[t]{0.48\textwidth}
        \vspace{0pt}
        \centering
        \includegraphics[width=\linewidth]{img/voter/static-ncl/voter_static_ncl_mse_against_correlation_penalty_coefficient_0_1_fixed_architecture_2_2_2_fixed_epoch_num_400_fixed_ensemble_size_12.png}
        \caption{Impact of correlation penalty coefficient (in the range $[0, 1]$) on MSE of a NCL ensemble model. Base learners had 3 layers of 2 hidden nodes and were trained with 400 epochs. Ensemble size was 12. }
        \label{fig:voter_static_ncl_mse_against_correlation_penalty_coefficient_0_1_fixed_architecture_2_2_2_fixed_epoch_num_400_fixed_ensemble_size_12}
    \end{minipage}
    \hfill
    \begin{minipage}[t]{0.48\textwidth}
        \vspace{0pt}
        \centering
        \includegraphics[width=\linewidth]{img/voter/static-ncl/voter_static_ncl_mse_against_correlation_penalty_coefficient_0_1_fixed_architecture_2_2_2_2_2_2_fixed_epoch_num_400_fixed_ensemble_size_12.png}
        \caption{Impact of correlation penalty coefficient (in the range $[0, 1]$) on MSE of a NCL ensemble model. Base learners had 6 layers of 2 hidden nodes and were trained with 400 epochs. Ensemble size was 12. }
        \label{fig:voter_static_ncl_mse_against_correlation_penalty_coefficient_0_1_fixed_architecture_2_2_2_2_2_2_fixed_epoch_num_400_fixed_ensemble_size_12}
    \end{minipage}
\end{figure}


% 2 Image
\begin{figure}[H]
    \centering

    \begin{minipage}[t]{0.48\textwidth}
        \vspace{0pt}
        \centering
        \includegraphics[width=\linewidth]{img/voter/neat-ncl/unconnected/voter_neat_ncl_mse_against_generation_index_0_400_fixed_ensemble_size_50_unconnected.png}
        \caption{Ensemble MSE of a NEAT-NCL ensemble model throughout generations of evolution when the ensemble size is 50 (\lstinline[language=ini]|initial_connection = unconnected|)}
        \label{fig:voter_neat_ncl_mse_against_generation_index_0_400_fixed_ensemble_size_50_unconnected}
    \end{minipage}
    \hfill
    \begin{minipage}[t]{0.48\textwidth}
        \vspace{0pt}
        \centering
        \includegraphics[width=\linewidth]{img/voter/neat-ncl/unconnected/voter_neat_ncl_mse_against_generation_index_0_400_fixed_ensemble_size_100_unconnected.png}
        \caption{Ensemble MSE of a NEAT-NCL ensemble model throughout generations of evolution when the ensemble size is 100 (\lstinline[language=ini]|initial_connection = unconnected|)}
        \label{fig:voter_neat_ncl_mse_against_generation_index_0_400_fixed_ensemble_size_100_unconnected}
    \end{minipage}
\end{figure}


% 2 Image
\begin{figure}[H]
    \centering

    \begin{minipage}[t]{0.48\textwidth}
        \vspace{0pt}
        \centering
        \includegraphics[width=\linewidth]{img/voter/neat-ncl/unconnected/voter_neat_ncl_mse_against_generation_index_0_400_fixed_ensemble_size_150_unconnected.png}
        \caption{Ensemble MSE of a NEAT-NCL ensemble model throughout generations of evolution when the ensemble size is 150 (\lstinline[language=ini]|initial_connection = unconnected|)}
        \label{fig:voter_neat_ncl_mse_against_generation_index_0_400_fixed_ensemble_size_150_unconnected}
    \end{minipage}
    \hfill
    \begin{minipage}[t]{0.48\textwidth}
        \vspace{0pt}
        \centering
        \includegraphics[width=\linewidth]{img/voter/neat-ncl/unconnected/voter_neat_ncl_mse_against_generation_index_0_400_fixed_ensemble_size_200_unconnected.png}
        \caption{Ensemble MSE of a NEAT-NCL ensemble model throughout generations of evolution when the ensemble size is 200 (\lstinline[language=ini]|initial_connection = unconnected|)}
        \label{fig:voter_neat_ncl_mse_against_generation_index_0_400_fixed_ensemble_size_200_unconnected}
    \end{minipage}
\end{figure}


% 2 Image
\begin{figure}[H]
    \centering

    \begin{minipage}[t]{0.48\textwidth}
        \vspace{0pt}
        \centering
        \includegraphics[width=\linewidth]{img/voter/neat-ncl/unconnected/voter_neat_ncl_mse_against_generation_index_0_400_fixed_ensemble_size_250_unconnected.png}
        \caption{Ensemble MSE of a NEAT-NCL ensemble model throughout generations of evolution when the ensemble size is 250 (\lstinline[language=ini]|initial_connection = unconnected|)}
        \label{fig:voter_neat_ncl_mse_against_generation_index_0_400_fixed_ensemble_size_250_unconnected}
    \end{minipage}
    \hfill
    \begin{minipage}[t]{0.48\textwidth}
        \vspace{0pt}
        \centering
        \includegraphics[width=\linewidth]{img/voter/neat-ncl/partial-direct/voter_neat_ncl_mse_against_generation_index_0_400_fixed_ensemble_size_50_partial_direct.png}
        \caption{Ensemble MSE of a NEAT-NCL ensemble model throughout generations of evolution when the ensemble size is 50 (\lstinline[language=ini]|initial_connection = partial_direct 0.5|)}
        \label{fig:voter_neat_ncl_mse_against_generation_index_0_400_fixed_ensemble_size_50_partial_direct}
    \end{minipage}
\end{figure}


% 2 Image
\begin{figure}[H]
    \centering

    \begin{minipage}[t]{0.48\textwidth}
        \vspace{0pt}
        \centering
        \includegraphics[width=\linewidth]{img/voter/neat-ncl/partial-direct/voter_neat_ncl_mse_against_generation_index_0_400_fixed_ensemble_size_100_partial_direct.png}
        \caption{Ensemble MSE of a NEAT-NCL ensemble model throughout generations of evolution when the ensemble size is 100 (\lstinline[language=ini]|initial_connection = partial 0.5|)}
        \label{fig:voter_neat_ncl_mse_against_generation_index_0_400_fixed_ensemble_size_100_partial_direct}
    \end{minipage}
    \hfill
    \begin{minipage}[t]{0.48\textwidth}
        \vspace{0pt}
        \centering
        \includegraphics[width=\linewidth]{img/voter/neat-ncl/partial-direct/voter_neat_ncl_mse_against_generation_index_0_400_fixed_ensemble_size_150_partial_direct.png}
        \caption{Ensemble MSE of a NEAT-NCL ensemble model throughout generations of evolution when the ensemble size is 150 (\lstinline[language=ini]|initial_connection = partial_direct 0.5|)}
        \label{fig:voter_neat_ncl_mse_against_generation_index_0_400_fixed_ensemble_size_150_partial_direct}
    \end{minipage}
\end{figure}


% 2 Image
\begin{figure}[H]
    \centering

    \begin{minipage}[t]{0.48\textwidth}
        \vspace{0pt}
        \centering
        \includegraphics[width=\linewidth]{img/voter/neat-ncl/partial-direct/voter_neat_ncl_mse_against_generation_index_0_400_fixed_ensemble_size_200_partial_direct.png}
        \caption{Ensemble MSE of a NEAT-NCL ensemble model throughout generations of evolution when the ensemble size is 200 (\lstinline[language=ini]|initial_connection = partial 0.5|)}
        \label{fig:voter_neat_ncl_mse_against_generation_index_0_400_fixed_ensemble_size_200_partial_direct}
    \end{minipage}
    \hfill
    \begin{minipage}[t]{0.48\textwidth}
        \vspace{0pt}
        \centering
        \includegraphics[width=\linewidth]{img/voter/neat-ncl/partial-direct/voter_neat_ncl_mse_against_generation_index_0_400_fixed_ensemble_size_250_partial_direct.png}
        \caption{Ensemble MSE of a NEAT-NCL ensemble model throughout generations of evolution when the ensemble size is 250 (\lstinline[language=ini]|initial_connection = partial_direct 0.5|)}
        \label{fig:voter_neat_ncl_mse_against_generation_index_0_400_fixed_ensemble_size_250_partial_direct}
    \end{minipage}
\end{figure}


% 2 Image
\begin{figure}[H]
    \centering

    \begin{minipage}[t]{0.48\textwidth}
        \vspace{0pt}
        \centering
        \includegraphics[width=\linewidth]{img/voter/neat-ncl/full-direct/voter_neat_ncl_mse_against_generation_index_0_400_fixed_ensemble_size_50_full_direct.png}
        \caption{Ensemble MSE of a NEAT-NCL ensemble model throughout generations of evolution when the ensemble size is 50 (\lstinline[language=ini]|initial_connection = full_direct|)}
        \label{fig:voter_neat_ncl_mse_against_generation_index_0_400_fixed_ensemble_size_50_full_direct}
    \end{minipage}
    \hfill
    \begin{minipage}[t]{0.48\textwidth}
        \vspace{0pt}
        \centering
        \includegraphics[width=\linewidth]{img/voter/neat-ncl/full-direct/voter_neat_ncl_mse_against_generation_index_0_400_fixed_ensemble_size_100_full_direct.png}
        \caption{Ensemble MSE of a NEAT-NCL ensemble model throughout generations of evolution when the ensemble size is 100 (\lstinline[language=ini]|initial_connection = full_direct|)}
        \label{fig:voter_neat_ncl_mse_against_generation_index_0_400_fixed_ensemble_size_100_full_direct}
    \end{minipage}
\end{figure}


% 2 Image
\begin{figure}[H]
    \centering

    \begin{minipage}[t]{0.48\textwidth}
        \vspace{0pt}
        \centering
        \includegraphics[width=\linewidth]{img/voter/neat-ncl/full-direct/voter_neat_ncl_mse_against_generation_index_0_400_fixed_ensemble_size_150_full_direct.png}
        \caption{Ensemble MSE of a NEAT-NCL ensemble model throughout generations of evolution when the ensemble size is 150 (\lstinline[language=ini]|initial_connection = full_direct|)}
        \label{fig:voter_neat_ncl_mse_against_generation_index_0_400_fixed_ensemble_size_150_full_direct}
    \end{minipage}
    \hfill
    \begin{minipage}[t]{0.48\textwidth}
        \vspace{0pt}
        \centering
        \includegraphics[width=\linewidth]{img/voter/neat-ncl/full-direct/voter_neat_ncl_mse_against_generation_index_0_400_fixed_ensemble_size_200_full_direct.png}
        \caption{Ensemble MSE of a NEAT-NCL ensemble model throughout generations of evolution when the ensemble size is 200 (\lstinline[language=ini]|initial_connection = full_direct|)}
        \label{fig:voter_neat_ncl_mse_against_generation_index_0_400_fixed_ensemble_size_200_full_direct}
    \end{minipage}
\end{figure}


% 1 Image
\begin{figure}[H]
    \centering
    \includegraphics[width=0.48\textwidth]{img/voter/neat-ncl/full-direct/voter_neat_ncl_mse_against_generation_index_0_400_fixed_ensemble_size_250_full_direct.png}
    \caption{Ensemble MSE of a NEAT-NCL ensemble model throughout generations of evolution when the ensemble size is 250 (\lstinline[language=ini]|initial_connection = full_direct|)}
    \label{fig:voter_neat_ncl_mse_against_generation_index_0_400_fixed_ensemble_size_250_full_direct}
\end{figure}


\section{Stage 6: Recursive architecture}
\label{sec:stage_6}


\subsection{Motivations}
In Section~\ref{sec:stage_5}, different aggregators had been explored, which gave different ensemble performance. However, each ensemble model always had 1 depth, i.e. all base learners' predictions contribute to the final ensemble predictions immediately. In previous sections, each ensemble model was treated as 1 model from an external perspective. However, it was actually composed of $m$ base learners from an internal perspective. What if this idea is expaned that each $i$th base learner actually consists of $m_{i}$ smaller base learners and so on etc. until the base case (1 MLP) is reached? This idea is shown in Figure~\ref{fig:recursive_idea}. From an internal view, each model \texttt{Model i} can either be a MLP or an ensemble model (of any types) itself. This gives an infinite number of ensemble architectures, i.e. ways of organising base learners. 2 particular ensemble architectures were explored in this section. 

% 1 Image
\begin{figure}[H]
    \centering
    \includegraphics[width=0.64\textwidth]{img/recursive/recursive_idea.png}
    \caption{The idea of applying recusion to an ensemble model. }
    \label{fig:recursive_idea}
\end{figure}


\subsection{Design and Implementation}


\subsubsection{Representing recursively using a tree}
Each recursive architecture can be represented in the form of a tree, where each node represents an ensemble model and each edge represents a connection from this ensemble model to 1 of its base learners. Thus, there is an aggregator (can be of different types) in each node. Figure~\ref{fig:recursive_tree_example} shows an example recursive architecture. The height of the tree is $2$. It is composed of 4 ensemble models, represented by Node~2, Node~3, Node~5 and Node~6. Node~1 and Node~4 are simply particular aggregators that combine predictions from their successors. If they are neural-network-aggregators, they have different number of input features because they have different number of successors ($3$ for Node~1 and $2$ for Node~4). Node~2, Node~3, Node~5 and Node~6 consist of $m_2$, $m_3$, $m_5$ and $m_6$ base learners respectively. Each of them aggregates the predictions with their own aggregator (which can be of different types). In this experiment, $m_1 = m_2 = m_3 = m_4$ and have range $[2, 4]$. Furthermore, all nodes in the tree use the same type of aggregator but this is not a must. 

% 1 Image
\begin{figure}[H]
    \centering
    \includegraphics[width=0.64\textwidth]{img/recursive/recursive_tree_example.png}
    \caption{The tree representing an example recursive architecture of a recursive ensemble model. Each node is an instance of \texttt{RecursiveEnsembleNet}. The grey area represents an aggregator. The blue area represents an ensemble instance (e.g. \texttt{RecursiveStaticNCLEnsembleNet} instance) acting as a base learner. The blue arrows represent the data flows of predictions between \texttt{RecursiveEnsembleNet} instances. }
    \label{fig:recursive_tree_example}
\end{figure}


\subsubsection{Model construction}
New versions of the ensemble models, \texttt{RecursiveTraditionalEnsembleNet} and \texttt{RecursiveStaticNCLEnsembleNet}, had been created. They functioned exactly like the original versions (\texttt{TraditionalEnsembleNet} and \texttt{StaticNCLEnsembleNet}) from Section~\ref{sec:stage_2} and Section~\ref{sec:stage_3} but they did not use \textit{DistributedDataParallel} or \textit{Multiprocessing}. Instead of splitting the work across two GPUs, the entire model now ran on a single GPU. 

The overall structure was managed by \texttt{RecursiveEnsembleNet}, which acted as a coordinator. It used two main components, \texttt{\_\_base\_learners} and \texttt{\_\_voter}. The \texttt{\_\_base\_learners} attribute was a list that can hold other recursive networks (non-leaf nodes) or the specific ensemble models (leaf nodes) like \texttt{RecursiveStaticNCLEnsembleNet} or \texttt{NEATNCLEnsembleNet}. 

In this setup, each \texttt{RecursiveEnsembleNet} at a non-leaf level used its \texttt{\_\_voter} to combine the predictions coming from the layer right below it. Because of this layered design, the total number of aggregators in the entire model was equal to the height of the tree plus the number of leaf nodes, i.e. $n_{aggregator} = height_{tree} + n_{leaf\_node}$. This allowed for a nested, hierarchical way to combine multiple different ensemble types into one final prediction. 


\subsubsection{Experiment}
There were 2 different types of ensemble architectures used in the experiment, a \textit{perfect binary tree} and a \textit{flat tree}. Figure~\ref{fig:recursive_perfect_binary_tree} and Figure~\ref{fig:recursive_flat_tree} show their corresponding ensemble tree structures. Their ensemble performance was compared against the traditional ensemble model in Section~\ref{sec:stage_2} and the NCL ensemble model in Section~\ref{sec:stage_3}. 
Note that $m$ of all ensemble models in this experiment were multiples of $4$. The perfect binary tree picked was composed of 4 ensemble models. For fair comparison with a flat tree, $m$ was set to be the same, i.e. multiples of $4$. 

% 2 Image
\begin{figure}[H]
    \centering

    \begin{minipage}[t]{0.48\textwidth}
        \vspace{0pt}
        \centering
        \includegraphics[width=\linewidth]{img/recursive/recursive_perfect_binary_tree.png}
        \caption{The tree representation of the recursive ensemble model whose architecture is a perfect binary tree. Each node is an instance of \texttt{RecursiveEnsembleNet}. The grey area represents an aggregator. The blue area represents an ensemble instance (e.g. \texttt{RecursiveStaticNCLEnsembleNet} instance) acting as a base learner. The blue arrows represent the data flows of predictions between \texttt{RecursiveEnsembleNet} instances. }
        \label{fig:recursive_perfect_binary_tree}
    \end{minipage}
    \hfill
    \begin{minipage}[t]{0.48\textwidth}
        \vspace{0pt}
        \centering
        \includegraphics[width=\linewidth]{img/recursive/recursive_flat_tree.png}
        \caption{The tree representation of the recursive ensemble model whose architecture is a flat tree. Each node is an instance of \texttt{RecursiveEnsembleNet}. The grey area represents an aggregator. The blue area represents an ensemble instance (e.g. \texttt{RecursiveStaticNCLEnsembleNet} instance) acting as a base learner. The blue arrows represent the data flows of predictions between \texttt{RecursiveEnsembleNet} instances. }
        \label{fig:recursive_flat_tree}
    \end{minipage}
\end{figure}

\paragraph{Varied hyperparameters}
Table~\ref{table:experiment_6_varied_hyperparameters} contains all varied hyperparameters in this experiment: 

\begin{table}[H]
    \centering
    \begin{tabular}{|c|c|p{3cm}|p{5cm}|}
        \hline
        Hyperparameter & Type & Choices / Ranges & Description \\
        \hline
        \small\texttt{architectures} & Categorical & \{perfect binary tree, flat tree\} & Ensemble model architecture \\
        \hline
        \small\texttt{correlation\_penalty\_coefficients} & Float & $[0.4, 0.9]\text{, step}=0.1$ & $\lambda$ (for NCL ensemble model only) \\
        \hline
        \small\texttt{ensemble\_sizes} & Integer & \{$[2, 16]\text{, step}=1$, $[50, 250]\text{, step}=50$\} & $m$ ($[50, 250]$ for recursively NEAT-NCL ensemble model and $[2, 16]$ for the others) \\
        \hline
        \small\texttt{voter} & Categorical & \{\texttt{arithmetic\_mean}, \texttt{median}, \texttt{nn}\} & Aggregator type (\texttt{nn} is not applicable for recursively NEAT-NCL ensemble model) \\
        \hline
    \end{tabular}
    \caption{Varied hyperparameters and their configurations}
    \label{table:experiment_6_varied_hyperparameters}
\end{table}

\paragraph{Fixed hyperparameters}
Table~\ref{table:experiment_6_fixed_hyperparameters} contains all fixed hyperparameters in this experiment: 

\begin{table}[H]
    \centering
    \begin{tabular}{|c|c|c|p{7cm}|}
        \hline
        Hyperparameter & Type & Choices / Ranges & Description \\
        \hline
        \small\texttt{input\_size} & Integer & $13$ & Number of input features of the MLP \\
        \hline
        \small\texttt{output\_size} & Integer & $1$ & Number of labels of the MLP \\
        \hline
        \small\texttt{optimizer} & Categorical & \texttt{torch.optim.Adam} & Optimisation algorithm used for training \\
        \hline
        \small\texttt{hidden\_sizes} & Python list & \texttt{[6]} & Model architecture \\
        \hline
        \small\texttt{learning\_rate} & Float & $0.001$ & Learning rate of the optimisation algorithm \\
        \hline
        \small\texttt{dropout\_rate} & Float & $0.5$ & Dropout rate applied to the data tensor after the input layer \\
        \hline
        \small\texttt{activation} & Categorical & ReLU & Activation function used in each dense layer of the MLP \\
        \hline
        \small\texttt{epoch\_nums} & Integer & $400$ & Number of training epochs (for arithmetic-mean-aggregator and NCL ensemble model only) \\
        \hline
        \small\texttt{evolution\_epoch} & Integer & $400$ & Number of evolutionary generations during training (for recursively NEAT-NCL ensemble model only) \\
        \hline
        \small\texttt{ensemble\_sizes} & Integer & $200$ & $m$ (for recursively NEAT-NCL ensemble model only) \\
        \hline
        \small\texttt{batch\_size} & Integer & \{$32$, $506$\} & Number of training samples passed to the model in one go ($506$ for recursively NEAT-NCL ensemble model and $32$ for the others) \\
        \hline
    \end{tabular}
    \caption{Fixed hyperparameters and their configurations}
    \label{table:experiment_6_fixed_hyperparameters}
\end{table}

\paragraph{Metric}
Table~\ref{table:experiment_6_metric} contains all evaluation metrics collected in this experiment: 

\begin{table}[H]
    \centering
    \begin{tabular}{|c|c|c|}
        \hline
        Metric & Type & Description \\
        \hline
        \small\texttt{loss\_test} & Float & MSE of ensemble model during testing \\
        \hline
    \end{tabular}
    \caption{Evaluation metrics}
    \label{table:experiment_6_metric}
\end{table}

\paragraph{Potential improvements}
Recursive ensemble model composed of different types of base learners could be explored with additional time. For example, 1 base learner in the recursive model could  be a traditional ensemble model while the 2 base learners could be a NCL ensemble model and a recursively NEAT-NCL ensemble model respectively. 
Furthermore, more variations in ensemble architecture could be included in the experiment. In this experiment, all base learners either formed a tree-shape or were either at the same depth as shown in Figure~\ref{fig:recursive_perfect_binary_tree} and Figure~\ref{fig:recursive_flat_tree}. Architectures shown in Figure~\ref{fig:recursive_tree_potential_1} and Figure~\ref{fig:recursive_tree_potential_2} could be also attempted to investigate the best ensemble architecture for this problem. 
Lastly, each MLP was fixed to have 1 hidden layer of 6 hidden nodes. Other MLP architectures, e.g. 1 hidden layer of 2 or 12 hidden nodes, could be attempted. 

% 2 Image
\begin{figure}[H]
    \centering

    \begin{minipage}[t]{0.48\textwidth}
        \vspace{0pt}
        \centering
        \includegraphics[width=\linewidth]{img/recursive/recursive_tree_potential_1.png}
        \caption{The tree representing an example recursive architecture of a potential recursive ensemble model that can be experimented with additional time. Each node is an instance of \texttt{RecursiveEnsembleNet}. The grey area represents an aggregator. The blue area represents an ensemble instance (e.g. \texttt{RecursiveStaticNCLEnsembleNet} instance) acting as a base learner. The blue arrows represent the data flows of predictions between \texttt{RecursiveEnsembleNet} instances. }
        \label{fig:recursive_tree_potential_1}
    \end{minipage}
    \hfill
    \begin{minipage}[t]{0.48\textwidth}
        \vspace{0pt}
        \centering
        \includegraphics[width=\linewidth]{img/recursive/recursive_tree_potential_2.png}
        \caption{The tree representing an example recursive architecture of a potential recursive ensemble model that can be experimented with additional time. Each node is an instance of \texttt{RecursiveEnsembleNet}. The grey area represents an aggregator. The blue area represents an ensemble instance (e.g. \texttt{RecursiveStaticNCLEnsembleNet} instance) acting as a base learner. The blue arrows represent the data flows of predictions between \texttt{RecursiveEnsembleNet} instances. }
        \label{fig:recursive_tree_potential_2}
    \end{minipage}
\end{figure}

\paragraph{Further notes}
All output metrics collected were stored in \path{experiment_runner_recursive_traditional.csv}, \path{experiment_runner_recursive_static_ncl.csv}, \path{experiment_runner_recursive_neat_ncl_unconnected.csv}, \path{experiment_runner_recursive_neat_ncl_partial_direct.csv}, \path{experiment_runner_recursive_neat_ncl_full_direct.csv}. 


% TODO: Is it ok to have no expected observations? 
% \subsection{Expected Observations}


% \subsubsection{...}
% ...


\subsection{Evaluation and Empirical Comparisons}


\subsubsection{Similar recursively traditional ensemble performance for different recursive architectures}
\paragraph{Experiments and Observations}
In a recursive version of a traditional ensemble model (from Section~\ref{sec:stage_2}), the specific structure of the recursive tree (whether it is a perfect binary tree or a flat tree) did not significantly change the final ensemble performance. According to Figure~\ref{fig:recursive_traditional_mse_against_ensemble_size_8_16_fixed_architecture_6_fixed_epoch_num_400_fixed_aggregator_arithmetic_mean}, Figure~\ref{fig:recursive_traditional_mse_against_ensemble_size_8_16_fixed_architecture_6_fixed_epoch_num_400_fixed_aggregator_median} and Figure~\ref{fig:recursive_traditional_mse_against_ensemble_size_8_16_fixed_architecture_6_fixed_epoch_num_400_fixed_aggregator_nn}, the blue and orange lines representing these two shapes were almost identical. 

This suggested that for recursively traditional ensembles, the way of layering the groups of models was less important than the quality of the base learners themselves. Even when using different aggregators like the arithmetic-mean-aggregator, median-aggregator or a neural-network-aggregator, switching from a simple flat list of models to a multi-layered perfect binary tree did not provide a noticeable improvement in accuracy or a lower ensemble MSE. 

\paragraph{Conclusions from Experiment}
Using different available recursive architectures does not give performance improvement for a recursively traditional ensemble model. 


\subsubsection{Recursive architecture combined with aggregator type affect recursively NCL ensemble performance}
\paragraph{Experiments and Observations}
For a recursively NCL ensemble model, the effectiveness of recursive architectures depended heavily on the type of aggregator used. 

When using an arithmetic-mean-aggregator, there was no significant difference between a flat tree and a perfect binary tree. As shown in Figure~\ref{fig:recursive_static_ncl_mse_against_correlation_penalty_coefficient_04_09_fixed_architecture_6_fixed_ensemble_size_8_fixed_aggregator_arithmetic_mean} and Figure~\ref{fig:recursive_static_ncl_mse_against_correlation_penalty_coefficient_04_09_fixed_architecture_6_fixed_ensemble_size_16_fixed_aggregator_arithmetic_mean}, the blue and orange lines overlapped frequently. While increasing the correlation penalty coefficient $\lambda$ helped lower the MSE for both lines, the standard NCL model from Section~\ref{sec:stage_3} (the green line) performed much better. This suggested that recursive structures actually hindered performance when using simple averaging. 

With a median-aggregator, a flat tree consistently outperformed a perfect binary tree. Figure~\ref{fig:recursive_static_ncl_mse_against_correlation_penalty_coefficient_04_09_fixed_architecture_6_fixed_ensemble_size_8_fixed_aggregator_median} and Figure~\ref{fig:recursive_static_ncl_mse_against_correlation_penalty_coefficient_04_09_fixed_architecture_6_fixed_ensemble_size_16_fixed_aggregator_median} showed the orange line staying below the blue line. Interestingly, for a perfect binary tree, increasing $\lambda$ actually made the MSE worse while it had little effect on the flat tree. Again, the standard non-recursively NCL ensemble model from Section~\ref{sec:stage_3} remained superior, indicating that recursive architectures was also not a good fit for median-based voting. 

In contrast, when a neural-network-aggregator was used, recursive architectures became beneficial. A flat tree generally performed better than a perfect binary tree, as seen in Figure~\ref{fig:recursive_static_ncl_mse_against_correlation_penalty_coefficient_04_09_fixed_architecture_6_fixed_ensemble_size_8_fixed_aggregator_nn} and Figure~\ref{fig:recursive_static_ncl_mse_against_correlation_penalty_coefficient_04_09_fixed_architecture_6_fixed_ensemble_size_16_fixed_aggregator_nn}. Crucially, both recursive versions (blue and orange lines) often outperformed the control NCL model from Section~\ref{sec:stage_3}. This suggested that the hierarchical structure helped the neural network learn more effective ways to combine diverse predictions, despite some increased fluctuations in performance (comparing to lines in Figure~\ref{fig:recursive_traditional_mse_against_ensemble_size_8_16_fixed_architecture_6_fixed_epoch_num_400_fixed_aggregator_arithmetic_mean}, Figure~\ref{fig:recursive_static_ncl_mse_against_correlation_penalty_coefficient_04_09_fixed_architecture_6_fixed_ensemble_size_8_fixed_aggregator_arithmetic_mean}, Figure~\ref{fig:recursive_static_ncl_mse_against_correlation_penalty_coefficient_04_09_fixed_architecture_6_fixed_ensemble_size_8_fixed_aggregator_median} and Figure~\ref{fig:recursive_static_ncl_mse_against_correlation_penalty_coefficient_04_09_fixed_architecture_6_fixed_ensemble_size_16_fixed_aggregator_median}). 

\paragraph{Conclusions from Experiment}
Both recursive architecture and aggregator type have an impact on recursively NCL ensemble performance. When a neural-network-aggregator is used, a recursive architecture gives ensemble performance improvement but not for other aggregators. Among the two recursive architecture, a flat tree seems to be a more favourable candidate than a perfect binary tree for any given aggregator type. 


\subsubsection{Perfect binary tree is better when an arithmetic-mean-aggregator is used in a recursively NEAT-NCL ensemble model}
\paragraph{Experiments and Observations}
When an arithmetic-mean-aggregator was used in a recursively NEAT-NCL ensemble, a perfect binary tree structure performed better than a flat tree. According to the results in Table~\ref{table:experiment_6_neat_ncl_unconnected}, Table~\ref{table:experiment_6_neat_ncl_partial_direct} and Table~\ref{table:experiment_6_neat_ncl_full_direct}, the test loss (\texttt{loss\_test}) was consistently lower for the binary tree setup (first row) compared to the flat tree setup (second row). 

This suggested that for the highly diverse base learners produced by NEAT-NCL, a hierarchical approach to averaging was more effective. By grouping and averaging predictions in stages rather than all at once, the perfect binary tree likely did a better job of managing the structural and functional variety of the base learners, leading to a more accurate final prediction. 

\paragraph{Conclusions from Experiment}
When an arithmetic-mean-aggregator was in a recursively NEAT-NCL ensemble model, arranging base learners in a perfect binary tree gives better performance than a flat tree. 


\subsubsection{Recursive architectures give similar performance when a median-aggregator is used in a recursively NEAT-NCL ensemble model}
\paragraph{Experiments and Observations}
When a median-aggregator was used in a recursively NEAT-NCL ensemble, the specific tree architecture (whether a perfect binary tree or a flat tree) did not significantly impact ensemble performance. 

Data from Table~\ref{table:experiment_6_neat_ncl_unconnected}, Table~\ref{table:experiment_6_neat_ncl_partial_direct} and Table~\ref{table:experiment_6_neat_ncl_full_direct} showed that the test loss (\texttt{loss\_test}) in the third row was approximately the same as the value in the fourth row. This suggested that the median is a naturally robust aggregator for NEAT-NCL models and its ability to filter out extreme predictions worked equally well regardless of whether the models were organised in a hierarchy or a simple flat list. 

\paragraph{Conclusions from Experiment}
When a median-aggregator was in a recursively NEAT-NCL ensemble model, arranging base learners in a perfect binary tree gives similar ensemble performance as arranging base learners in a flat tree. 


\subsubsection{A median-aggregator is better than an arithmetic-mean-aggregator for a recursively NEAT-NCL ensemble model}
\paragraph{Experiments and Observations}
The median-aggregator consistently outperformed the arithmetic-mean-aggregator in NEAT-NCL experiments. Data from Table~\ref{table:experiment_6_neat_ncl_unconnected}, Table~\ref{table:experiment_6_neat_ncl_partial_direct} and Table~\ref{table:experiment_6_neat_ncl_full_direct} showed that when the arithmetic-mean-aggregator was used, the test loss (\texttt{loss\_test}) stayed within $[0.331. 0.417]$. In contrast, the median-aggregator achieved a significantly lower test loss of approximately $0.286$. This improvement was likely due to the high diversity in NEAT-NCL models. Since these models were forced to be very different, some individual learners might produce extreme errors or ``outlier'' predictions. An arithmetic-mean-aggregator is sensitive to these extremes and gets pulled away from the correct value. The median-agggregator, however, effectively ignores these outliers by selecting the middle value, leading to a more accurate and stable ensemble prediction. 

\paragraph{Conclusions from Experiment}
For a recursively NEAT-NCL ensemble model, a median-aggregator gives better ensemble performance than an arithmetic-mean-aggregator and should be used. 


\subsubsection{\texttt{loss\_test} when a median-aggregator and a flat tree are used are the same in all recursively NEAT-NCL ensemble models}
\paragraph{Experiments and Observations}
The test loss (\texttt{loss\_test}) remained exactly $0.28602060675621$ across all experiments where a median-aggregator and a flat tree were used, regardless of the initial connection settings. This consistency, shown in the last rows of Table~\ref{table:experiment_6_neat_ncl_unconnected}, Table~\ref{table:experiment_6_neat_ncl_partial_direct} and Table~\ref{table:experiment_6_neat_ncl_full_direct}, was a direct result of the evolutionary convergence of the models. 

As discussed in Section~\ref{sec:stage_4}, the most effective architecture for this problem was a base learner with no hidden nodes. While the models (phenomes) started with different configurations (e.g. different number of connections), they all evolved over 400 generations toward this same simple structure. Since base learners of different recursively NEAT-NCL ensemble models converged to nearly identical structures, they produced similar distributions of predictions. 

The median-aggregator on a flat tree played a key role in this result. By selecting the middle value, it effectively ignored the tiny variations in the distributions of the base learners' predictions in different recursively NEAT-NCL ensemble models. It consistently picked the prediction from a specific, dominant genome at each depth. Since this process happened the same way across all depths of all recursively NEAT-NCL models, the final ensemble prediction (and therefore the final test loss) ended up being exactly the same every time. 

\paragraph{Conclusions from Experiment}
When a median-aggregator and a flat tree is used, the performance of the recursively NEAT-NCL ensemble model remains unchanged despite the value of \lstinline[language=ini]|initial_connection|. 


% ================================ Images ================================


% 2 Image
\begin{figure}[H]
    \centering

    \begin{minipage}[t]{0.48\textwidth}
        \vspace{0pt}
        \centering
        \includegraphics[width=\linewidth]{img/recursive/recursive_traditional_mse_against_ensemble_size_8_16_fixed_architecture_6_fixed_epoch_num_400_fixed_aggregator_arithmetic_mean.png}
        \caption{Impact of ensemble size on MSE of a recursively traditional ensemble model. Base learners had 1 layer of 6 hidden nodes and were trained with 400 epochs. An arithmetic-mean-aggregator was used. }
        \label{fig:recursive_traditional_mse_against_ensemble_size_8_16_fixed_architecture_6_fixed_epoch_num_400_fixed_aggregator_arithmetic_mean}
    \end{minipage}
    \hfill
    \begin{minipage}[t]{0.48\textwidth}
        \vspace{0pt}
        \centering
        \includegraphics[width=\linewidth]{img/recursive/recursive_traditional_mse_against_ensemble_size_8_16_fixed_architecture_6_fixed_epoch_num_400_fixed_aggregator_median.png}
        \caption{Impact of ensemble size on MSE of a recursively traditional ensemble model. Base learners had 1 layer of 6 hidden nodes and were trained with 400 epochs. A median-aggregator was used. }
        \label{fig:recursive_traditional_mse_against_ensemble_size_8_16_fixed_architecture_6_fixed_epoch_num_400_fixed_aggregator_median}
    \end{minipage}
\end{figure}


% 2 Image
\begin{figure}[H]
    \centering

    \begin{minipage}[t]{0.48\textwidth}
        \vspace{0pt}
        \centering
        \includegraphics[width=\linewidth]{img/recursive/recursive_traditional_mse_against_ensemble_size_8_16_fixed_architecture_6_fixed_epoch_num_400_fixed_aggregator_nn.png}
        \caption{Impact of ensemble size on MSE of a recursively traditional ensemble model. Base learners had 1 layer of 6 hidden nodes and were trained with 400 epochs. A neural-network-aggregator was used. }
        \label{fig:recursive_traditional_mse_against_ensemble_size_8_16_fixed_architecture_6_fixed_epoch_num_400_fixed_aggregator_nn}
    \end{minipage}
    \hfill
    \begin{minipage}[t]{0.48\textwidth}
        \vspace{0pt}
        \centering
        \includegraphics[width=\linewidth]{img/recursive/recursive_static_ncl_mse_against_correlation_penalty_coefficient_04_09_fixed_architecture_6_fixed_ensemble_size_8_fixed_aggregator_arithmetic_mean.png}
        \caption{Impact of correlation penalty coefficient (in range $[0.4, 0.9]$) on MSE of a recursively NCL ensemble model. Base learners had 1 layer of 6 hidden nodes and were trained with 400 epochs. Ensemble size was 8. An arithmetic-mean-aggregator was used. }
        \label{fig:recursive_static_ncl_mse_against_correlation_penalty_coefficient_04_09_fixed_architecture_6_fixed_ensemble_size_8_fixed_aggregator_arithmetic_mean}
    \end{minipage}
\end{figure}


% 2 Image
\begin{figure}[H]
    \centering

    \begin{minipage}[t]{0.48\textwidth}
        \vspace{0pt}
        \centering
        \includegraphics[width=\linewidth]{img/recursive/recursive_static_ncl_mse_against_correlation_penalty_coefficient_04_09_fixed_architecture_6_fixed_ensemble_size_8_fixed_aggregator_median.png}
        \caption{Impact of correlation penalty coefficient (in range $[0.4, 0.9]$) on MSE of a recursively NCL ensemble model. Base learners had 1 layer of 6 hidden nodes and were trained with 400 epochs. Ensemble size was 8. A median-aggregator was used. }
        \label{fig:recursive_static_ncl_mse_against_correlation_penalty_coefficient_04_09_fixed_architecture_6_fixed_ensemble_size_8_fixed_aggregator_median}
    \end{minipage}
    \hfill
    \begin{minipage}[t]{0.48\textwidth}
        \vspace{0pt}
        \centering
        \includegraphics[width=\linewidth]{img/recursive/recursive_static_ncl_mse_against_correlation_penalty_coefficient_04_09_fixed_architecture_6_fixed_ensemble_size_8_fixed_aggregator_nn.png}
        \caption{Impact of correlation penalty coefficient (in range $[0.4, 0.9]$) on MSE of a recursively NCL ensemble model. Base learners had 1 layer of 6 hidden nodes and were trained with 400 epochs. Ensemble size was 8. A neural-network-aggregator was used. }
        \label{fig:recursive_static_ncl_mse_against_correlation_penalty_coefficient_04_09_fixed_architecture_6_fixed_ensemble_size_8_fixed_aggregator_nn}
    \end{minipage}
\end{figure}


% 2 Image
\begin{figure}[H]
    \centering

    \begin{minipage}[t]{0.48\textwidth}
        \vspace{0pt}
        \centering
        \includegraphics[width=\linewidth]{img/recursive/recursive_static_ncl_mse_against_correlation_penalty_coefficient_04_09_fixed_architecture_6_fixed_ensemble_size_16_fixed_aggregator_arithmetic_mean.png}
        \caption{Impact of correlation penalty coefficient (in range $[0.4, 0.9]$) on MSE of a recursively NCL ensemble model. Base learners had 1 layer of 6 hidden nodes and were trained with 400 epochs. Ensemble size was 16. An arithmetic-mean-aggregator was used. }
        \label{fig:recursive_static_ncl_mse_against_correlation_penalty_coefficient_04_09_fixed_architecture_6_fixed_ensemble_size_16_fixed_aggregator_arithmetic_mean}
    \end{minipage}
    \hfill
    \begin{minipage}[t]{0.48\textwidth}
        \vspace{0pt}
        \centering
        \includegraphics[width=\linewidth]{img/recursive/recursive_static_ncl_mse_against_correlation_penalty_coefficient_04_09_fixed_architecture_6_fixed_ensemble_size_16_fixed_aggregator_median.png}
        \caption{Impact of correlation penalty coefficient (in range $[0.4, 0.9]$) on MSE of a recursively NCL ensemble model. Base learners had 1 layer of 6 hidden nodes and were trained with 400 epochs. Ensemble size was 16. A median-aggregator was used. }
        \label{fig:recursive_static_ncl_mse_against_correlation_penalty_coefficient_04_09_fixed_architecture_6_fixed_ensemble_size_16_fixed_aggregator_median}
    \end{minipage}
\end{figure}


% 1 Image
\begin{figure}[H]
    \centering
    \includegraphics[width=0.48\textwidth]{img/recursive/recursive_static_ncl_mse_against_correlation_penalty_coefficient_04_09_fixed_architecture_6_fixed_ensemble_size_16_fixed_aggregator_nn.png}
    \caption{Impact of correlation penalty coefficient (in range $[0.4, 0.9]$) on MSE of a recursively NCL ensemble model. Base learners had 1 layer of 6 hidden nodes and were trained with 400 epochs. Ensemble size was 16. A neural-network-aggregator was used. }
    \label{fig:recursive_static_ncl_mse_against_correlation_penalty_coefficient_04_09_fixed_architecture_6_fixed_ensemble_size_16_fixed_aggregator_nn}
\end{figure}


\begin{table}[H]
    \centering
    \begin{tabular}{|c|c|p{4cm}|p{5cm}|c|}
        \hline
        Metric & Type & Description & Configurations & Value \\
        \hline
        \small\texttt{loss\_test} & Float & MSE of a recursively NEAT-NCL ensemble model during testing & \{\texttt{ensemble\_architecture}: perfect binary tree, \texttt{aggregator\_type}: arithmetic mean\} & $0.36160245537757874$ \\
        \hline
        \small\texttt{loss\_test} & Float & MSE of a recursively NEAT-NCL ensemble model during testing & \{\texttt{ensemble\_architecture}: flat tree, \texttt{aggregator\_type}: arithmetic mean\} & $0.32550227642059326$ \\
        \hline
        \small\texttt{loss\_test} & Float & MSE of a recursively NEAT-NCL ensemble model during testing & \{\texttt{ensemble\_architecture}: perfect binary tree, \texttt{aggregator\_type}: median\} & $0.2860206067562103$ \\
        \hline
        \small\texttt{loss\_test} & Float & MSE of a recursively NEAT-NCL ensemble model during testing & \{\texttt{ensemble\_architecture}: flat tree, \texttt{aggregator\_type}: median\} & $0.2860206067562103$ \\
        \hline
    \end{tabular}
    \caption{Testing loss of a recursively NEAT-NCL ensemble model after 400 generations (\lstinline[language=ini]|initial_connection = unconnected|). }
    \label{table:experiment_6_neat_ncl_unconnected}
\end{table}


\begin{table}[H]
    \centering
    \begin{tabular}{|c|c|p{4cm}|p{5cm}|c|}
        \hline
        Metric & Type & Description & Configurations & Value \\
        \hline
        \small\texttt{loss\_test} & Float & MSE of a recursively NEAT-NCL ensemble model during testing & \{\texttt{ensemble\_architecture}: perfect binary tree, \texttt{aggregator\_type}: arithmetic mean\} & $0.3384256660938263$ \\
        \hline
        \small\texttt{loss\_test} & Float & MSE of a recursively NEAT-NCL ensemble model during testing & \{\texttt{ensemble\_architecture}: flat tree, \texttt{aggregator\_type}: arithmetic mean\} & $0.32768353819847107$ \\
        \hline
        \small\texttt{loss\_test} & Float & MSE of a recursively NEAT-NCL ensemble model during testing & \{\texttt{ensemble\_architecture}: perfect binary tree, \texttt{aggregator\_type}: median\} & $0.2860206067562103$ \\
        \hline
        \small\texttt{loss\_test} & Float & MSE of a recursively NEAT-NCL ensemble model during testing & \{\texttt{ensemble\_architecture}: flat tree, \texttt{aggregator\_type}: median\} & $0.2860206067562103$ \\
        \hline
    \end{tabular}
    \caption{Testing loss of a recursively NEAT-NCL ensemble model after 400 generations (\lstinline[language=ini]|initial_connection = partial_direct 0.5|). }
    \label{table:experiment_6_neat_ncl_partial_direct}
\end{table}


\begin{table}[H]
    \centering
    \begin{tabular}{|c|c|p{4cm}|p{5cm}|c|}
        \hline
        Metric & Type & Description & Configurations & Value \\
        \hline
        \small\texttt{loss\_test} & Float & MSE of a recursively NEAT-NCL ensemble model during testing & \{\texttt{ensemble\_architecture}: perfect binary tree, \texttt{aggregator\_type}: arithmetic mean\} & $0.34439876675605774$ \\
        \hline
        \small\texttt{loss\_test} & Float & MSE of a recursively NEAT-NCL ensemble model during testing & \{\texttt{ensemble\_architecture}: flat tree, \texttt{aggregator\_type}: arithmetic mean\} & $0.38286858797073364$ \\
        \hline
        \small\texttt{loss\_test} & Float & MSE of a recursively NEAT-NCL ensemble model during testing & \{\texttt{ensemble\_architecture}: perfect binary tree, \texttt{aggregator\_type}: median\} & $0.2860206067562103$ \\
        \hline
        \small\texttt{loss\_test} & Float & MSE of a recursively NEAT-NCL ensemble model during testing & \{\texttt{ensemble\_architecture}: flat tree, \texttt{aggregator\_type}: median\} & $0.2860206067562103$ \\
        \hline
    \end{tabular}
    \caption{Testing loss of a recursively NEAT-NCL ensemble model after 400 generations (\lstinline[language=ini]|initial_connection = full_direct|). }
    \label{table:experiment_6_neat_ncl_full_direct}
\end{table}


\appendix


\section{Loss Function of a NCL Base Learner}
\label{appendix:loss_function_ncl}
\[
\mathcal{L}_i = (\hat{y}_i - y)^2 - \lambda (\hat{y}_i - \bar{y})^2
\]
\begin{center}
    \begin{minipage}{0.8\linewidth}
        \small
        where $\hat{y}_i$ is the prediction made by the $i$th base learner,  
        $\bar{y}$ is the mean prediction of all base learners and  
        $\lambda$ is the correlation penalty coefficient (typically in $[0,1]$ when ensemble size is large). 
    \end{minipage}
\end{center}


\section{Diversity Coefficient}
\label{appendix:diversity_coefficient}
\[
\rho = \frac{1}{m} \sum_{i=1}^{m} (\hat{y}_i - \bar{y})^{2}
\]
\begin{center}
    \begin{minipage}{0.8\linewidth}
        \small
        where $\hat{y}_i$ is the prediction made by the $i$th base learner and  
        $\bar{y}$ is the mean prediction of all base learners. 
    \end{minipage}
\end{center}


\section{Mean Squared Error}
\label{appendix:mse}
\[
MSE = \frac{1}{m} \sum_{i=1}^{m} (\hat{y}_i - y)^{2}
\]
\begin{center}
    \begin{minipage}{0.8\linewidth}
        \small
        where $\hat{y}_i$ is the prediction made by the $i$th base learner and  
        $y$ is the true value. 
    \end{minipage}
\end{center}


\section{Root Mean Squared Error}
\label{appendix:rmse}
\[
RMSE = \sqrt{\frac{1}{m} \sum_{i=1}^{m} (\hat{y}_i - y)^{2}}
\]
\begin{center}
    \begin{minipage}{0.8\linewidth}
        \small
        where $\hat{y}_i$ is the prediction made by the $i$th base learner and  
        $y$ is the true value. 
    \end{minipage}
\end{center}


\section{Correlation Penalty Function Conversion}
\label{appendix:correlation_penalty_function_conversion}
\begin{align*}
    &p_i = (\hat{y}_i - \bar{y}) \sum_{j \neq i} (\hat{y}_j - \bar{y})
    \\
    &= (\hat{y}_i - \bar{y}) \left[\sum_{j \neq i} (\hat{y}_j - \bar{y}) + \hat{y}_i - \bar{y} - \hat{y}_i + \bar{y}\right]
    \\
    &= (\hat{y}_i - \bar{y}) \left[0 - \hat{y}_i + \bar{y}\right] \quad \text{(}\sum_{j \neq i} (\hat{y}_j - \bar{y}) + \hat{y}_i - \bar{y} = 0\text{)}
    \\
    &= (\hat{y}_i - \bar{y})^2
\end{align*}
\begin{center}
    \begin{minipage}{0.8\linewidth}
        \small
        where $\hat{y}_i$ is the prediction made by the $i$th base learner and  
        $\bar{y}$ the mean prediction of all base learners. 
    \end{minipage}
\end{center}


\section{$\lambda_{upper\_bound}$}
\label{appendix:lambda_upper_bound}
\[
\lambda_{upper\_bound} = \frac{m^2}{2(m - 1)^2}
\]
\begin{center}
    \begin{minipage}{0.8\linewidth}
        \small
        where $m$ is the ensemble size. This formula is copied from Brown et al.\ \cite{brown2005managing}. 
    \end{minipage}
\end{center}


\section{Proof of Theorem~\ref{theorem:bias_variance_decomposition}}
\label{appendix:bias_variance_decomposition}
\[
\underbrace{\mathbb{E}_{\mathcal{S}}\left[\mathbb{E}_{\mathbf{x}y}\left[(y - \hat{y})^2\right]\right]}_{\text{expected risk}} = \underbrace{\mathbb{E}_{\mathbf{x}y}\left[(y - \hat{y^*})^2\right]}_{\text{noise}} + \underbrace{\mathbb{E}_{\mathbf{x}y}\left[\left(\hat{y^*} - \mathbb{E}_{\mathcal{S}}[\hat{y}]\right)^2\right]}_{\text{bias}} + \underbrace{\mathbb{E}_{\mathbf{x}y}\left[\mathbb{E}_{\mathcal{S}}\left[(\hat{y} - \mathbb{E}_{\mathcal{S}}[\hat{y}])^2\right]\right]}_{\text{variance}}
\]
\begin{center}
    \begin{minipage}{0.8\linewidth}
        \small
        where $\hat{y}$ is the prediction made by a model and $\hat{y^*}$ is the prediction of the Bayes Model. 
    \end{minipage}
    \begin{proof}
        \begin{align*}
            &\mathbb{E}_{\mathcal{S}}\left[\mathbb{E}_{\mathbf{x}y}\left[(y - \hat{y})^2\right]\right] - \mathbb{E}_{\mathbf{x}y}\left[(y - \hat{y^*})^2\right]
            \\
            &= \underbrace{\mathbb{E}_{\mathcal{S}}\left[\mathbb{E}_{\mathbf{x}y}\left[(y - \hat{y})^2\right]\right] - R(\mathbb{E}_{\mathcal{S}}\left[\hat{y}\right])}_{\text{variance}} + \underbrace{R(\mathbb{E}_{\mathcal{S}}\left[\hat{y}\right]) - \mathbb{E}_{\mathbf{x}y}\left[(y - \hat{y^*})^2\right]}_{\text{bias}}
            \\
            % ========================================
            &\mathbb{E}_{\mathcal{S}}\left[\mathbb{E}_{\mathbf{x}y}\left[(y - \hat{y})^2\right]\right] - R(\mathbb{E}_{\mathcal{S}}\left[\hat{y}\right])
            \\
            &= \mathbb{E}_{\mathcal{S}}\left[\mathbb{E}_{\mathbf{x}y}\left[(y - \hat{y})^2\right]\right] - \mathbb{E}_{\mathbf{x}y}\left[(y - \mathbb{E}_{\mathcal{S}}[\hat{y}])^2\right]
            \\
            &= \mathbb{E}_{\mathcal{S}}\left[\mathbb{E}_{\mathbf{x}y}\left[(y - \hat{y})^2 - (y - \mathbb{E}_{\mathcal{S}}[\hat{y}])^2\right]\right]
            \\
            &= \mathbb{E}_{\mathcal{S}}\left[\mathbb{E}_{\mathbf{x}y}\left[y^2 + \hat{y}^2 - 2y\hat{y} - y^2 - \mathbb{E}_{\mathcal{S}}[\hat{y}]^2 + 2y\mathbb{E}_{\mathcal{S}}[\hat{y}]\right]\right]
            \\
            &= \mathbb{E}_{\mathcal{S}}\left[\mathbb{E}_{\mathbf{x}y}\left[\hat{y}^2 - 2y\hat{y} - \mathbb{E}_{\mathcal{S}}[\hat{y}]^2 + 2y\mathbb{E}_{\mathcal{S}}[\hat{y}]\right]\right]
            \\
            &= \mathbb{E}_{\mathbf{x}y}\left[\mathbb{E}_{\mathcal{S}}[\hat{y}^2] - 2y\mathbb{E}_{\mathcal{S}}[\hat{y}] - \mathbb{E}_{\mathcal{S}}[\hat{y}]^2 + 2y\mathbb{E}_{\mathcal{S}}[\hat{y}]\right]
            \\
            &= \mathbb{E}_{\mathbf{x}y}\left[\mathbb{E}_{\mathcal{S}}[\hat{y}^2] - \mathbb{E}_{\mathcal{S}}[\hat{y}]^2\right]
            \\
            &= \underbrace{\mathbb{E}_{\mathbf{x}y}\left[\mathbb{E}_{\mathcal{S}}\left[(\hat{y} - \mathbb{E}_{\mathcal{S}}[\hat{y}])^2\right]\right]}_{\text{variance}} \quad \text{(}\mathbb{E}\left[Z^2\right] - (\mathbb{E}\left[Z^2\right])^2 = \mathbb{E}\left[(Z - \mathbb{E}\left[Z\right])^2\right]\text{)}
            \\
            % ========================================
            &R(\mathbb{E}_{\mathcal{S}}\left[\hat{y}\right]) - \mathbb{E}_{\mathbf{x}y}\left[(y - \hat{y^*})^2\right]
            \\
            &=\mathbb{E}_{\mathbf{x}y}\left[(y - \mathbb{E}_{\mathcal{S}}[\hat{y}])^2\right] - \mathbb{E}_{\mathbf{x}y}\left[(y - \hat{y^*})^2\right]
            \\
            &= \mathbb{E}_{\mathbf{x}y}\left[(y - \mathbb{E}_{\mathcal{S}}[\hat{y}])^2 - (y - \hat{y^*})^2\right]
            \\
            &= \mathbb{E}_{\mathbf{x}y}\left[y^2 + \mathbb{E}_{\mathcal{S}}[\hat{y}]^2 - 2y\mathbb{E}_{\mathcal{S}}[\hat{y}] - y^2 - (\hat{y^*})^2 + 2y\hat{y^*}\right]
            \\
            &= \mathbb{E}_{\mathbf{x}}\left[\mathbb{E}_{y \mid \mathbf{x}}\left[\mathbb{E}_{\mathcal{S}}[\hat{y}]^2 - 2y\mathbb{E}_{\mathcal{S}}[\hat{y}] - (\hat{y^*})^2 + 2y\hat{y^*}\right]\right] \quad \text{(}\mathbb{E}_{\mathbf{x}y}\left[...\right] = \mathbb{E}_{\mathbf{x}}\left[\mathbb{E}_{y \mid \mathbf{x}}\left[...\right]\right]\text{)}
            \\
            &= \mathbb{E}_{\mathbf{x}}\left[\mathbb{E}_{\mathcal{S}}[\hat{y}]^2 - 2\hat{y^*}\mathbb{E}_{\mathcal{S}}[\hat{y}] - (\hat{y^*})^2 + 2(\hat{y^*})^2\right] \quad \text{(}\hat{y^*} = \mathbb{E}_{y \mid \mathbf{x}}\left[y\right]\text{)}
            \\
            &= \mathbb{E}_{\mathbf{x}}\left[\mathbb{E}_{\mathcal{S}}[\hat{y}]^2 - 2\hat{y^*}\mathbb{E}_{\mathcal{S}}[\hat{y}] + (\hat{y^*})^2\right]
            \\
            &= \mathbb{E}_{\mathbf{x}}\left[\left(\hat{y^*} - \mathbb{E}_{\mathcal{S}}[\hat{y}]\right)^2\right]
            \\
            &= \underbrace{\mathbb{E}_{\mathbf{x}y}\left[\left(\hat{y^*} - \mathbb{E}_{\mathcal{S}}[\hat{y}]\right)^2\right]}_{\text{bias}}
            \\
            % ========================================
            &\therefore \mathbb{E}_{\mathcal{S}}\left[\mathbb{E}_{\mathbf{x}y}\left[(y - \hat{y})^2\right]\right] - \mathbb{E}_{\mathbf{x}y}\left[(y - \hat{y^*})^2\right] = \underbrace{\mathbb{E}_{\mathbf{x}y}\left[\mathbb{E}_{\mathcal{S}}\left[(\hat{y} - \mathbb{E}_{\mathcal{S}}[\hat{y}])^2\right]\right]}_{\text{variance}} + \underbrace{\mathbb{E}_{\mathbf{x}y}\left[\left(\hat{y^*} - \mathbb{E}_{\mathcal{S}}[\hat{y}]\right)^2\right]}_{\text{bias}}
            \\
            % ========================================
            &\therefore \underbrace{\mathbb{E}_{\mathcal{S}}\left[\mathbb{E}_{\mathbf{x}y}\left[(y - \hat{y})^2\right]\right]}_{\text{expected risk}} = \underbrace{\mathbb{E}_{\mathbf{x}y}\left[(y - \hat{y^*})^2\right]}_{\text{noise}} + \underbrace{\mathbb{E}_{\mathbf{x}y}\left[\left(\hat{y^*} - \mathbb{E}_{\mathcal{S}}[\hat{y}]\right)^2\right]}_{\text{bias}} + \underbrace{\mathbb{E}_{\mathbf{x}y}\left[\mathbb{E}_{\mathcal{S}}\left[(\hat{y} - \mathbb{E}_{\mathcal{S}}[\hat{y}])^2\right]\right]}_{\text{variance}}
        \end{align*}
    \end{proof}
\end{center}


\section{Proof of Theorem~\ref{theorem:ncl_differentiation_incorrect}}
\label{appendix:ncl_differentiation_incorrect}
\[
\frac{\partial{p_i}}{\partial{\hat{y}_i}} = -(\hat{y}_i - \bar{y})
\]
\begin{center}
    \begin{minipage}{0.8\linewidth}
        \small
        where it is assumed that $\bar{y}$ has a constant value with respect to $\hat{y}_i$. The formula is proposed by Liu et al.\ \cite{liu1999negative}. 
    \end{minipage}
    \begin{proof}
        \begin{align*}
            &\frac{\partial{p_i}}{\partial{\hat{y}_i}}
            \\
            &= \frac{\partial{(\hat{y}_i - \bar{y})\sum_{j \neq i} (\hat{y}_j - \bar{y})}}{\partial{\hat{y}_i}}
            \\
            &= \sum_{j \neq i} (\hat{y}_j - \bar{y}) \frac{\partial{(\hat{y}_i - \bar{y})}}{\partial{\hat{y}_i}} \quad \text{(}\bar{y}\text{ is assumed to have a constant value with respect to }\hat{y}_i\text{)}
            \\
            &= \sum_{j \neq i} (\hat{y}_j - \bar{y}) (1 - 0)
            \\
            &= \sum_{j \neq i} (\hat{y}_j - \bar{y})
            \\
            &= \sum_{j \neq i} (\hat{y}_j - \bar{y}) + \hat{y}_i - \bar{y} - \hat{y}_i + \bar{y}
            \\
            &= 0 - \hat{y}_i + \bar{y}
            \\
            &= -(\hat{y}_i - \bar{y})
        \end{align*}
    \end{proof}
\end{center}


\section{Proof of Theorem~\ref{theorem:ncl_differentiation_correct}}
\label{appendix:ncl_differentiation_correct}
\[
\frac{\partial{p_i}}{\partial{\hat{y}_i}} = \frac{-2 (m - 1)}{m} (\hat{y}_i - \bar{y})
\]
\begin{center}
    \begin{minipage}{0.8\linewidth}
        \small
        where $\bar{y}$ is not assumed to have a constant value with respect to $\hat{y}_i$. It is a modified version of the formula proposed by Liu et al.\ \cite{liu1999negative}. 
    \end{minipage}
    \begin{proof}
        \begin{align*}
            &\frac{\partial{p_i}}{\partial{\hat{y}_i}}
            \\
            &= \frac{\partial{(\hat{y}_i - \bar{y})^2}}{\partial{\hat{y}_i}}
            \\
            &= -2 (\hat{y}_i - \bar{y}) \frac{\partial{(\hat{y}_i - \bar{y})}}{\partial{\hat{y}_i}}
            \\
            &= -2 (\hat{y}_i - \bar{y}) \frac{\partial{(\hat{y}_i - \frac{1}{m} \hat{y}_1 - \frac{1}{m} \hat{y}_2 - ... - \frac{1}{m} \hat{y}_i - ... - \frac{1}{m} \hat{y}_m)}}{\partial{\hat{y}_i}} \quad \text{(}\bar{y} = \frac{1}{m} \sum_{i = 1}^m \hat{y}_i\text{)}
            \\
            &= -2 (\hat{y}_i - \bar{y}) (1 - 0 - 0 ... - \frac{1}{m} - ... - 0)
            \\
            &= \frac{-2 (m - 1)}{m} (\hat{y}_i - \bar{y})
        \end{align*}
    \end{proof}
\end{center}


\section{Proof of Theorem~\ref{theorem:ensemble_population_risk_le_average_population_risk}}
\label{appendix:ensemble_population_risk_le_average_population_risk}
\[
R(\bar{y}) \le \frac{1}{m} \sum_{i=1}^{m} R(\hat{y}_i)
\]
\begin{center}
    \begin{minipage}{0.8\linewidth}
        \small
        where $R(\bar{y})$ is ensemble population risk and $R(\hat{y}_i)$ is the population risk of the $i$th base learner. 
    \end{minipage}
    \begin{proof}
        \begin{align*}
            &R(\bar{y})
            \\
            &= \mathbb{E}_{\bar{y}y}[(\bar{y} - y)^2] \quad \text{(MSE is used)}
            \\
            &= \mathbb{E}_{\bar{y}y}[\frac{1}{m} \sum_{i = 1}^{m} (\bar{y}_i - y)^2 - \frac{1}{m} \sum_{i = 1}^{m} (\bar{y}_i - \hat{y})^2] \quad \text{(According to Theorem~\ref{theorem:ambiguity_decomposition})}
            \\
            &\leq \mathbb{E}_{\bar{y}y}[\frac{1}{m} \sum_{i = 1}^{m} (\bar{y}_i - y)^2]
        \end{align*}
        \begin{align*}
            &\mathbb{E}_{\bar{y}y}[\frac{1}{m} \sum_{i = 1}^{m} (\bar{y}_i - y)^2]
            \\
            &= \frac{1}{m} \sum_{i = 1}^{m} \mathbb{E}_{\bar{y}y}[(\bar{y}_i - y)^2]
            \\
            &= \frac{1}{m} \sum_{i = 1}^{m} R(\hat{y}_i)
            \\
            &\therefore R(\bar{y}) \leq \frac{1}{m} \sum_{i = 1}^{m} R(\hat{y}_i)
        \end{align*}
    \end{proof}
\end{center}


\section{Proof of Theorem~\ref{theorem:ambiguity_decomposition}}
\label{appendix:ambiguity_decomposition}
\[
\underbrace{(\bar{y} - y)^2}_{\text{ensemble loss}} = \underbrace{\frac{1}{m} \sum_{i=1}^{m} (\hat{y}_i - y)^2}_{\text{average loss}} - \underbrace{\frac{1}{m} \sum_{i=1}^{m} (\hat{y}_i - \bar{y})^2}_{ambiguity}
\]
\begin{center}
    \begin{minipage}{0.8\linewidth}
        \small
        where $\hat{y}_i$ is the prediction made by the $i$th base learner, $y$ is the true value and $\bar{y}$ is the mean prediction of all base learners. 
    \end{minipage}
    \begin{proof}
        \begin{align*}
            &(\bar{y} - y)^2
            \\
            &= (\bar{y} - y)(\bar{y} - y)
            \\
            &= (\bar{y} - y)\left[\frac{1}{m} \sum_{i=1}^{m} \hat{y}_i + \bar{y} - y - \frac{1}{m} \sum_{i=1}^{m} \hat{y}_i\right]
            \\
            &= (\bar{y} - y)\left[\frac{1}{m} \sum_{i=1}^{m} \hat{y}_i + \frac{1}{m} \sum_{i=1}^{m} \hat{y}_i - \frac{1}{m} \sum_{i=1}^{m} y - \frac{1}{m} \sum_{i=1}^{m} \bar{y}\right] \quad \text{(By the definition of arithmetic mean combiner }\bar{y}\text{)}
            \\
            &= (\bar{y} - y) \left[\frac{1}{m} \sum_{i=1}^{m} (\hat{y}_i + \hat{y}_i - y - \bar{y})\right]
            \\
            &= \frac{1}{m} \sum_{i=1}^{m} \left[(\bar{y} - y) (\hat{y}_i + \hat{y}_i - y - \bar{y})\right]
            \\
            &= \frac{1}{m} \sum_{i=1}^{m} \left[((\hat{y}_i - y) - (\hat{y}_i - \bar{y})) ((\hat{y}_i - y) + (\hat{y}_i - \bar{y}))\right]
            \\
            &= \frac{1}{m} \sum_{i=1}^{m} \left[(\hat{y}_i - y)^2 - (\hat{y}_i - \bar{y})^2\right]
            \\
            &= \underbrace{\frac{1}{m} \sum_{i=1}^{m} (\hat{y}_i - y)^2}_{\text{average loss}} - \underbrace{\frac{1}{m} \sum_{i=1}^{m} (\hat{y}_i - \bar{y})^2}_{ambiguity}
        \end{align*}
    \end{proof}
\end{center}


\section{Proof of Theorem~\ref{theorem:bias_variance_diversity_decomposition}}
\label{appendix:bias_variance_diversity_decomposition}
\begin{align*}
    \underbrace{\mathbb{E}_{\mathcal{S}}\left[\mathbb{E}_{\mathbf{x}y}\left[(y - \bar{y})^2\right]\right]}_{\text{expected ensemble risk}} = 
    \\
    &\phantom{=} \underbrace{\mathbb{E}_{\mathbf{x}y}\left[(y - \hat{y^*})^2\right]}_{\text{noise}} + \underbrace{\frac{1}{m} \sum_{i=1}^{m} \mathbb{E}_{\mathbf{x}y}\left[\left(\hat{y^*} - \mathbb{E}_{\mathcal{S}}[\hat{y}_i]\right)^2\right]}_{\text{average bias}} + \underbrace{\frac{1}{m} \sum_{i=1}^{m} \mathbb{E}_{\mathbf{x}y}\left[\mathbb{E}_{\mathcal{S}}\left[(\hat{y}_i - \mathbb{E}_{\mathcal{S}}[\hat{y}_i])^2\right]\right]}_{\text{average variance}}
    \\
    &\phantom{=}  - \underbrace{\mathbb{E}_{\mathbf{x}y}\left[\mathbb{E}_{\mathcal{S}}\left[\frac{1}{m} \sum_{i=1}^{m} (\hat{y}_i - \bar{y})^2\right]\right]}_{\text{diversity}}
\end{align*}
\begin{center}
    \begin{minipage}{0.8\linewidth}
        \small
        where $\hat{y}_i$ is the prediction made by the $i$th base learner, $\bar{y}$ is the mean prediction of all base learners and $\hat{y^*}$ is the prediction of the Bayes Model. 
    \end{minipage}
    \begin{proof}
        \begin{align*}
            &\mathbb{E}_{\mathcal{S}}\left[\mathbb{E}_{\mathbf{x}y}\left[(y - \bar{y})^2\right]\right]
            \\
            &= \mathbb{E}_{\mathcal{S}}\left[\mathbb{E}_{\mathbf{x}y}\left[(\bar{y} - y)^2\right]\right]
            \\
            &= \mathbb{E}_{\mathcal{S}}\left[\mathbb{E}_{\mathbf{x}y}\left[\frac{1}{m} \sum_{i=1}^{m} (\hat{y}_i - y)^2 - \frac{1}{m} \sum_{i=1}^{m} (\hat{y}_i - \bar{y})^2\right]\right]
            \\
            &= \frac{1}{m} \sum_{i=1}^{m} \mathbb{E}_{\mathcal{S}}\left[\mathbb{E}_{\mathbf{x}y}\left[(\hat{y}_i - y)^2\right]\right] - \mathbb{E}_{\mathcal{S}}\left[\mathbb{E}_{\mathbf{x}y}\left[\frac{1}{m} \sum_{i=1}^{m} (\hat{y}_i - \bar{y})^2\right]\right]
            \\
            &= \frac{1}{m} \sum_{i=1}^{m} \mathbb{E}_{\mathbf{x}y}\left[(y - \hat{y^*})^2\right] + \frac{1}{m} \sum_{i=1}^{m} \mathbb{E}_{\mathbf{x}y}\left[\left(\hat{y^*} - \mathbb{E}_{\mathcal{S}}[\hat{y}_i]\right)^2\right]
            \\
            &\phantom{=} \quad + \frac{1}{m} \sum_{i=1}^{m} \mathbb{E}_{\mathbf{x}y}\left[\mathbb{E}_{\mathcal{S}}\left[(\hat{y}_i - \mathbb{E}_{\mathcal{S}}[\hat{y}_i])^2\right]\right] - \mathbb{E}_{\mathcal{S}}\left[\mathbb{E}_{\mathbf{x}y}\left[\frac{1}{m} \sum_{i=1}^{m} (\hat{y}_i - \bar{y})^2\right]\right] \quad \text{(Using Theorem~\ref{theorem:bias_variance_decomposition})}
            \\
            &= \mathbb{E}_{\mathbf{x}y}\left[(y - \hat{y^*})^2\right] + \frac{1}{m} \sum_{i=1}^{m} \mathbb{E}_{\mathbf{x}y}\left[\left(\hat{y^*} - \mathbb{E}_{\mathcal{S}}[\hat{y}_i]\right)^2\right] + \frac{1}{m} \sum_{i=1}^{m} \mathbb{E}_{\mathbf{x}y}\left[\mathbb{E}_{\mathcal{S}}\left[(\hat{y}_i - \mathbb{E}_{\mathcal{S}}[\hat{y}_i])^2\right]\right] - \mathbb{E}_{\mathbf{x}y}\left[\mathbb{E}_{\mathcal{S}}\left[\frac{1}{m} \sum_{i=1}^{m} (\hat{y}_i - \bar{y})^2\right]\right]
            \\
            % ========================================
            &\therefore \underbrace{\mathbb{E}_{\mathcal{S}}\left[\mathbb{E}_{\mathbf{x}y}\left[(y - \bar{y})^2\right]\right]}_{\text{expected ensemble risk}} = 
            \\
            &\phantom{=} \underbrace{\mathbb{E}_{\mathbf{x}y}\left[(y - \hat{y^*})^2\right]}_{\text{noise}} + \underbrace{\frac{1}{m} \sum_{i=1}^{m} \mathbb{E}_{\mathbf{x}y}\left[\left(\hat{y^*} - \mathbb{E}_{\mathcal{S}}[\hat{y}_i]\right)^2\right]}_{\text{average bias}} + \underbrace{\frac{1}{m} \sum_{i=1}^{m} \mathbb{E}_{\mathbf{x}y}\left[\mathbb{E}_{\mathcal{S}}\left[(\hat{y}_i - \mathbb{E}_{\mathcal{S}}[\hat{y}_i])^2\right]\right]}_{\text{average variance}} - \underbrace{\mathbb{E}_{\mathbf{x}y}\left[\mathbb{E}_{\mathcal{S}}\left[\frac{1}{m} \sum_{i=1}^{m} (\hat{y}_i - \bar{y})^2\right]\right]}_{\text{diversity}}
        \end{align*}
    \end{proof}
\end{center}


\section{Proof of Theorem~\ref{theorem:loss_function_ncl_lambda_1}}
\label{appendix:loss_function_ncl_lambda_1}
\[
\underbrace{\mathcal{L}_i}_{\text{ith base learner loss}} = \frac{1}{n} \sum^{n}_{i = 1} \underbrace{(\bar{y} - y)^2}_{\text{squared loss of 1 data point}}
\]
\begin{center}
    \begin{minipage}{0.8\linewidth}
        \small
        for ith base learner, when $\lambda = 1$, MSE is used and the ensemble model is an arithmetic-mean-aggregator. $n$ is the number of data points in the training dataset. 
    \end{minipage}
    \begin{proof}
        \begin{align*}
            &\mathcal{L}_{i} = \frac{1}{n} \sum^{n}_{i = 1} \left[(\hat{y}_i - y)^2 - \lambda (\hat{y}_i - \bar{y})^2\right]
            \\
            &= \frac{1}{n} \sum^{n}_{i = 1} \left[(\hat{y}_i - y)^2 - (\hat{y}_i - \bar{y})^2\right] \quad \text{(}\lambda = 1\text{)}
            \\
            &= \frac{1}{n} \sum^{n}_{i = 1} \left[(\hat{y}_i - y + \hat{y}_i - \bar{y})(\hat{y}_i - y - \hat{y}_i + \bar{y})\right]
            \\
            &= \frac{1}{n} \sum^{n}_{i = 1} \left[(2\hat{y}_i - y - \bar{y})(\bar{y} - y)\right]
            \\
            &= (\bar{y} - y) \frac{1}{n} \sum^{n}_{i = 1} \left[(2\hat{y}_i - y - \bar{y})\right]
            \\
            &= (\bar{y} - y) \frac{1}{n} \sum^{n}_{i = 1} \left[(2\bar{y} - y - \bar{y})\right] \quad \text{(By the definition of arithmetic mean combiner }\bar{y}\text{)}
            \\
            &= (\bar{y} - y) \frac{1}{n} \sum^{n}_{i = 1} \left[(\bar{y} - y)\right]
            \\
            &= \frac{1}{n} \sum^{n}_{i = 1} \left[(\bar{y} - y)(\bar{y} - y)\right]
            \\
            &= \frac{1}{n} \sum^{n}_{i = 1} \underbrace{(\bar{y} - y)^2}_{\text{squared loss of 1 data point}}
        \end{align*}
    \end{proof}
\end{center}


\bibliographystyle{IEEEtran} % Or plain, alpha, apalike, etc.
\bibliography{references}


\end{document}
